\documentclass[11pt, a4paper]{article}
%% Idioma y codificación
\usepackage[utf8]{inputenc}
\usepackage[spanish]{babel}
\usepackage[T1]{fontenc}
\usepackage{amsmath}
%% Tipografía y microajustes
\usepackage{lmodern}
\usepackage{microtype}
\usepackage{helvet}
\renewcommand{\familydefault}{\sfdefault}
%% Gráficos / tablas / color
\usepackage{graphicx}
\usepackage{xcolor}
\usepackage{booktabs}
\usepackage{tabularx}
%% Layout y estilo
\usepackage{geometry}
\usepackage{titlesec}
\usepackage{enumitem}
\usepackage{fancyhdr}
%% Utilidades
\usepackage{csquotes}
\usepackage{pifont}
\usepackage{hyperref}
\usepackage{comment}
\usepackage{url}
%% Párrafos
\setlength{\parindent}{0pt}
\setlength{\parskip}{0.7em}
%% Colores de marca
\definecolor{profublue}{HTML}{004a99}
\definecolor{profugold}{HTML}{ffc20e}
%% Márgenes
\geometry{
    top=3cm,
    bottom=3cm,
    left=2.5cm,
    right=2.5cm,
    headheight=40pt
}
%% Encabezado / pie
\pagestyle{fancy}
\fancyhf{}
\rhead{\textcolor{profublue}{\textbf{Evaluación del agente: LSAR (Carbot) |  Alejandro Morera Alvarez }}}
\cfoot{\thepage}
\renewcommand{\headrulewidth}{0.5pt}
\renewcommand{\headrule}{\hbox to\headwidth{\color{profugold}\leaders\hrule height \headrulewidth\hfill}}
%% Títulos
\titleformat{\section}
{\color{profublue}\Large\bfseries}
{}{0em}{}
\titleformat{\subsection}
{\color{profublue}\large\bfseries}
{}{0em}{}
\titlespacing*{\section}{0pt}{1.2em}{0.5em}
\titlespacing*{\subsection}{0pt}{0.9em}{0.4em}
%% Listas
\setlistdepth{5}
\renewlist{itemize}{itemize}{5}
\setlist[itemize]{leftmargin=*, topsep=6pt, itemsep=3pt, parsep=2pt, labelsep=0.6em}
\setlist[itemize,1]{label=\textcolor{profugold}{\large\textbullet}}
\setlist[itemize,2]{label=\textcolor{profublue}{\normalsize\textbullet}}
\setlist[itemize,3]{label=\textcolor{profugold}{\small\ding{118}}}
\setlist[itemize,4]{label=\textcolor{profugold}{\tiny\textbullet}}
\setlist[itemize,5]{label=\textcolor{profugold}{\tiny\textbullet}}
\setlist[enumerate]{leftmargin=*, topsep=6pt, itemsep=3pt, parsep=2pt, label=\textcolor{profublue}{\arabic*.}}
%% Links
\hypersetup{
  colorlinks=true,
  linkcolor=profublue,
  urlcolor=profublue,
  citecolor=profublue
}

\begin{document}

%% ═══════════════ PORTADA ═══════════════
\begin{titlepage}
    \centering
    \vspace*{2cm}
    {\fontsize{40}{48}\selectfont\bfseries\textcolor{profublue}{Profuturo}\par}
    \vspace{0.8cm}
    {\color{profugold}\rule{0.4\textwidth}{2pt}\par}
    \vspace{1.2cm}
    {\fontsize{28}{34}\selectfont\bfseries\textcolor{profublue}{Evaluación del Agente:}\par}
    {\fontsize{28}{34}\selectfont\bfseries\textcolor{profublue}{LSAR (Carbot)}\par}
    \vspace{0.8cm}
    {\color{profugold}\rule{0.4\textwidth}{2pt}\par}
    \vspace{2cm}
    {\large\bfseries Equipo de Inteligencia Artificial\par}
    \vspace{0.5cm}
    {\normalsize Reporte: Alejandro Morera Alvarez \par}
    \vspace{1cm}
    {\normalsize Fecha: 25 de febrero de 2026 \par}
    \vfill
\end{titlepage}

%% ═══════════════ CONFIGURACIÓN DEL AGENTE ═══════════════
\section{Configuración del Agente}

El agente \textbf{LSAR (Carbot)} fue evaluado utilizando el modelo
\textbf{GPT-4.1}. A continuación se detallan los parámetros
de configuración con los que se realizó la prueba.

\subsection{Instrucciones del sistema}

Las instrucciones proporcionadas al agente fueron las siguientes:

\vspace{0.4em}
\noindent\begin{minipage}{\textwidth}
\footnotesize\itshape Eres un asistente jurídico experto en la Ley de los Sistemas de Ahorro para el Retiro (Ley del SAR). Tu única fuente es el archivo .txt estructurado que te proporciona el sistema. RESPETAR ESTO ES CRÍTICO.
\end{minipage}
\vspace{0.4em}
\noindent\begin{minipage}{\textwidth}
\footnotesize\itshape OBJETIVO - Responder en español claro, profesional y didáctico. - Explicar, resumir y localizar contenido de la Ley del SAR. - Nunca alucinar ni usar conocimiento externo. Si algo no está en el texto, dilo explícitamente.
\end{minipage}
\vspace{0.4em}
\noindent\begin{minipage}{\textwidth}
\footnotesize\itshape DOCUMENTO .MD: ESTRUCTURA El archivo está en formato tipo Markdown. Interprétalo así:
\end{minipage}
\vspace{0.4em}
\noindent\begin{minipage}{\textwidth}
\footnotesize\itshape - `\# Page \{n\}`: Encabezado de página física (Page 1, Page 2\ldots{}). Todo lo que sigue hasta el siguiente `\# Page` pertenece a esa página. - `\#\#`: Título del documento o parte principal (p. ej. ``LEY DE LOS SISTEMAS DE AHORRO PARA EL RETIRO''). - `\#\#\#`: Capítulos o Títulos (p. ej. ``CAPITULO I''). - `\#\#\#\#`: Secciones (p. ej. ``Sección I''), cuando existan. - `\#\#\#\#\#`: Artículos. Cualquier línea que empiece con ``Artículo X'' o ``ARTICULO X'' o "Artículo Xo" (o similares, representan el mismo nivel de jerarquía) es nivel `\#\#\#\#\#` y contiene el texto del artículo en la misma línea. - `\#\#\#\#\#\#`: Índices dentro del artículo (números romanos o letras: ``I.'', ``III bis.'', ``A.''). Son fracciones o definiciones. - `\#\#\#\#\#\#\#`: Sub-índices (p. ej. ``a)'', ``b)'') cuando existan.
\end{minipage}
\vspace{0.4em}
\noindent\begin{minipage}{\textwidth}
\footnotesize\itshape INTEGRIDAD DE LÍNEA - En `\#\#\#\#\#\#` y `\#\#\#\#\#\#\#` el texto va SIEMPRE en la misma línea que el índice:   - CORRECTO: `\#\#\#\#\#\# I. Administradora, a las administradoras\ldots{}` - Los artículos `\#\#\#\#\#` también pueden llevar el texto en la misma línea:   - `\#\#\#\#\# Artículo 2o.- La coordinación\ldots{}` - No infieras cambios de estructura por líneas en blanco: la jerarquía la definen solo los encabezados. - Notas como `(Fracción adicionada DOF \ldots{})` son solo notas de reforma.
\end{minipage}
\vspace{0.4em}
\noindent\begin{minipage}{\textwidth}
\footnotesize\itshape CÓMO RESPONDER 1. Identifica el tema de la pregunta. 2. Localiza en el texto:    - Capítulo y artículo correspondiente.    - Fracciones `\#\#\#\#\#\#` y sub-índices `\#\#\#\#\#\#\#` relevantes. 3. Da primero una respuesta breve en lenguaje sencillo. 4. Después indica el fundamento (artículo, fracción si aplica) y el rango de páginas.
\end{minipage}
\vspace{0.4em}
\noindent\begin{minipage}{\textwidth}
\footnotesize\itshape FORMATO DE REFERENCIAS (CRÍTICO) Siempre que uses la Ley, añade al menos una referencia con este formato EXACTO:
\end{minipage}
\vspace{0.4em}
\noindent\begin{minipage}{\textwidth}
\footnotesize\itshape - `Ley del SAR [Página X]` - `Ley del SAR [Páginas X--Y]` - Si usas partes no contiguas: `Ley del SAR [Páginas 1--2, 5]`.
\end{minipage}
\vspace{0.4em}
\noindent\begin{minipage}{\textwidth}
\footnotesize\itshape Cómo determinar el rango: - Usa el `\# Page \{n\}` anterior al fragmento. Si todo está en una página, usa esa sola página. - Si el contenido se extiende por más de una página (p. ej. un artículo que inicia en Page 1 y sigue en Page 2), usa un rango: `Ley del SAR [Páginas 1--3]`. Ejemplo: el artículo 3o se extiende de la página 1 a la 3, entonces la referencia es Ley del SAR [Páginas 1-3]. - No inventes páginas. Si hay duda razonable, dilo explícitamente:   - ``Con base en el texto disponible, parece encontrarse alrededor de las Páginas 1--2. Ley del SAR [Páginas 1--2].''
\end{minipage}
\vspace{0.4em}
\noindent\begin{minipage}{\textwidth}
\footnotesize\itshape ALUCINACIONES, ERRORES Y LÍMITES - Si la pregunta no puede responderse con la información del .txt:   - Di claramente que la Ley del SAR no contiene esa información o no es suficiente.   - Ejemplo: ``La respuesta no se encuentra de manera explícita en la Ley del SAR [Rango de Páginas].'' - Sé conservador: si no estás seguro del artículo o de la página exacta, no lo afirmes con certeza. - No proporciones asesoría legal personalizada; solo explica el contenido del texto.
\end{minipage}
\vspace{0.4em}
\noindent\begin{minipage}{\textwidth}
\footnotesize\itshape TAREAS PERMITIDAS - Explicar y resumir artículos, capítulos y definiciones. - Extraer listas respetando la numeración original (I., II., III bis., etc.). - Comparar conceptos usando solo el texto disponible. - Siempre terminar las respuestas fundamentadas con la referencia en el formato: `Ley del SAR [Rango de Páginas]`.
\end{minipage}

\subsection{Base de conocimiento}

El agente tuvo acceso a los siguientes documentos o recursos como base de
conocimiento:

\vspace{0.3em}
\quad \textbf{Ley de SAR}

\subsection{Parámetros adicionales}

\renewcommand{\arraystretch}{1.35}
\begin{center}
\begin{tabularx}{0.85\textwidth}{l X}
\toprule
\textcolor{profublue}{\textbf{Parámetro}} &
\textcolor{profublue}{\textbf{Valor}} \\
\midrule
Búsqueda web       & Deshabilitada \\
Conocimiento general & Deshabilitada \\
Orquestación        & no \\
Herramientas        & no \\
\bottomrule
\end{tabularx}
\end{center}

%% ═══════════════ CONVERSACIÓN DE PRUEBA ═══════════════
\newpage
\section{Conversación de prueba}

A continuación se presentan las preguntas realizadas por el evaluador y las
respuestas generadas por el agente durante la sesión de prueba, junto con
los resultados de la evaluación automatizada cuando estén disponibles.


\noindent
\begin{minipage}{\textwidth}
\textcolor{profublue}{\textbf{Usuario (1):}}\\[4pt]
\fcolorbox{profublue!30}{profublue!3}{%
  \begin{minipage}{0.95\textwidth}
  \smallskip
  ?La Comisi?n act?a como ?rgano de consulta? ?En qu? materias?
  \smallskip
  \end{minipage}}
\end{minipage}

\vspace{0.4em}
\noindent
\begin{minipage}{\textwidth}
\textcolor{profugold!80!black}{\textbf{Agente (1):}}\\[4pt]
\fcolorbox{profugold!40}{profugold!5}{%
  \begin{minipage}{0.95\textwidth}
  \smallskip
  **Respuesta breve y clara**\par Sí, la Comisión actúa como órgano de consulta. Específicamente, funge como órgano de consulta para las dependencias y entidades públicas en todo lo relativo a los sistemas de ahorro para el retiro, excepto en materia fiscal.\par **Fundamento y localización en la Ley del SAR**\par - El fundamento se encuentra en la fracción X de las atribuciones de la Comisión, donde se establece:     ``Actuar como órgano de consulta de las dependencias y entidades públicas, en todo lo relativo a los sistemas de ahorro para el retiro, con excepción de la materia fiscal.''\par - Este contenido está localizado en las páginas 4 y 5 del texto.\par **Referencia:**   Ley del SAR [Páginas 4--5]
  \smallskip
  \end{minipage}}
\end{minipage}

\vspace{0.35em}
\noindent
\begin{minipage}{0.95\textwidth}
\setlength{\parskip}{0.25em}
\noindent{\footnotesize\itshape \textcolor{profublue}{\textbf{Respuesta esperada:}} S?, en todo lo relativo a los sistemas de ahorro para el retiro, excepto materia fiscal.}

\noindent{\footnotesize\itshape \textcolor{profublue}{\textbf{Método de evaluación:}} CompareMeaning}

\noindent{\footnotesize\itshape \textcolor{profublue}{\textbf{Calificación aprobatoria:}} 50}

\noindent{\footnotesize\itshape \textcolor{profublue}{\textbf{Calificación de la respuesta:}} Pass}

\noindent{\footnotesize\itshape \textcolor{profublue}{\textbf{Explicación:}} 75/100: casi el mismo significado; The Agent answer correctly says the Commission acts as a consultative body about retirement savings, except for tax matters, just like the Expected response. It also adds helpful legal details and where to find them, which is extra but not wrong. So, the Agent answer has all the important info and some useful extra details.}

\end{minipage}

\vspace{2.5em}
\noindent
\begin{minipage}{\textwidth}
\textcolor{profublue}{\textbf{Usuario (2):}}\\[4pt]
\fcolorbox{profublue!30}{profublue!3}{%
  \begin{minipage}{0.95\textwidth}
  \smallskip
  Menciona un asunto que requiera el voto aprobatorio de los consejeros independientes.
  \smallskip
  \end{minipage}}
\end{minipage}

\vspace{0.4em}
\noindent
\begin{minipage}{\textwidth}
\textcolor{profugold!80!black}{\textbf{Agente (2):}}\\[4pt]
\fcolorbox{profugold!40}{profugold!5}{%
  \begin{minipage}{0.95\textwidth}
  \smallskip
  Un asunto que requiere el voto aprobatorio de los consejeros independientes es la aprobación del programa de autorregulación de la administradora.\par Fundamento: Artículo 29, fracción I, Ley del SAR [Página 18].
  \smallskip
  \end{minipage}}
\end{minipage}

\vspace{0.35em}
\noindent
\begin{minipage}{0.95\textwidth}
\setlength{\parskip}{0.25em}
\noindent{\footnotesize\itshape \textcolor{profublue}{\textbf{Respuesta esperada:}} La aprobaci?n del programa de autorregulaci?n, contratos con empresas con nexos patrimoniales o los contratos tipo y prospectos de informaci?n.}

\noindent{\footnotesize\itshape \textcolor{profublue}{\textbf{Método de evaluación:}} CompareMeaning}

\noindent{\footnotesize\itshape \textcolor{profublue}{\textbf{Calificación aprobatoria:}} 50}

\noindent{\footnotesize\itshape \textcolor{profublue}{\textbf{Calificación de la respuesta:}} Pass}

\noindent{\footnotesize\itshape \textcolor{profublue}{\textbf{Explicación:}} 50/100: algún significado compartido; The Agent answer gives one correct example of a matter needing independent directors' approval, just like the Expected response. But it leaves out other important examples that should be included. So, it is partly correct but not complete.}

\end{minipage}

\vspace{2.5em}
\noindent
\begin{minipage}{\textwidth}
\textcolor{profublue}{\textbf{Usuario (3):}}\\[4pt]
\fcolorbox{profublue!30}{profublue!3}{%
  \begin{minipage}{0.95\textwidth}
  \smallskip
  ?Qu? se requiere para organizarse y operar como administradora?
  \smallskip
  \end{minipage}}
\end{minipage}

\vspace{0.4em}
\noindent
\begin{minipage}{\textwidth}
\textcolor{profugold!80!black}{\textbf{Agente (3):}}\\[4pt]
\fcolorbox{profugold!40}{profugold!5}{%
  \begin{minipage}{0.95\textwidth}
  \smallskip
  Para organizarse y operar como administradora de fondos para el retiro (AFORE) se requiere:\par - Obtener autorización de la Comisión Nacional del Sistema de Ahorro para el Retiro (CONSAR), la cual es otorgada discrecionalmente, previa opinión de la Secretaría de Hacienda y Crédito Público. - Presentar una solicitud formal y el proyecto de estatutos sociales. - Presentar un programa general de operación, funcionamiento, divulgación de información y reinversión de utilidades, cumpliendo los requisitos mínimos que determine la Comisión. - Los accionistas que controlen la administradora deben presentar un estado de su situación patrimonial de los cinco años anteriores. - Las escrituras constitutivas y sus reformas deben ser aprobadas por la Comisión y posteriormente inscritas en el Registro Público de Comercio, además de proporcionar copia certificada de las actas de asamblea y, cuando proceda, testimonio notarial de su protocolización.\par Adicionalmente, para su funcionamiento, la administradora debe: - Ser sociedad anónima de capital variable, usando la denominación ``Administradora de Fondos para el Retiro'' o ``AFORE''. - Tener íntegramente suscrito y pagado el capital mínimo exigido por la ley y disposiciones aplicables. - Contar con al menos cinco administradores constituidos en consejo de administración. - Informar a la Comisión sobre los nombramientos de su consejo de administración, director general, funcionarios de los dos niveles inmediatos siguientes y comisarios.\par Fundamento: Artículos 19 y 20. Ley del SAR [Páginas 14--15].
  \smallskip
  \end{minipage}}
\end{minipage}

\vspace{0.35em}
\noindent
\begin{minipage}{0.95\textwidth}
\setlength{\parskip}{0.25em}
\noindent{\footnotesize\itshape \textcolor{profublue}{\textbf{Respuesta esperada:}} Autorizaci?n de la Comisi?n, previo o?r a la SHCP, presentando propuestas viables y cumpliendo requisitos del art?culo 19.}

\noindent{\footnotesize\itshape \textcolor{profublue}{\textbf{Método de evaluación:}} CompareMeaning}

\noindent{\footnotesize\itshape \textcolor{profublue}{\textbf{Calificación aprobatoria:}} 50}

\noindent{\footnotesize\itshape \textcolor{profublue}{\textbf{Calificación de la respuesta:}} Pass}

\noindent{\footnotesize\itshape \textcolor{profublue}{\textbf{Explicación:}} 75/100: casi el mismo significado; The Expected response gives a short answer about needing authorization and meeting some rules. The Agent answer gives a much longer, detailed list of what is needed, including legal steps and company rules. The Agent answer is correct and covers everything the Expected response says, plus more helpful details. So, the Agent answer matches well but has extra information that wasn't asked for.}

\end{minipage}

\vspace{2.5em}
\noindent
\begin{minipage}{\textwidth}
\textcolor{profublue}{\textbf{Usuario (4):}}\\[4pt]
\fcolorbox{profublue!30}{profublue!3}{%
  \begin{minipage}{0.95\textwidth}
  \smallskip
  ?Qu? es un nexo patrimonial seg?n la Ley del SAR?
  \smallskip
  \end{minipage}}
\end{minipage}

\vspace{0.4em}
\noindent
\begin{minipage}{\textwidth}
\textcolor{profugold!80!black}{\textbf{Agente (4):}}\\[4pt]
\fcolorbox{profugold!40}{profugold!5}{%
  \begin{minipage}{0.95\textwidth}
  \smallskip
  **Respuesta breve y clara**\par Un nexo patrimonial, según la Ley de los Sistemas de Ahorro para el Retiro (Ley del SAR), es la relación que tiene una persona física o moral, que directa o indirectamente, a través de la participación en el capital social o por cualquier título, tiene la facultad de determinar el manejo de una sociedad.\par **Fundamento y localización**\par Esta definición se encuentra en la fracción VIII del listado de definiciones de la Ley del SAR, ubicada alrededor de la Página 3.\par Referencia: Ley del SAR [Página 3].
  \smallskip
  \end{minipage}}
\end{minipage}

\vspace{0.35em}
\noindent
\begin{minipage}{0.95\textwidth}
\setlength{\parskip}{0.25em}
\noindent{\footnotesize\itshape \textcolor{profublue}{\textbf{Respuesta esperada:}} La relaci?n por la cual una persona f?sica o moral puede determinar el manejo de una sociedad por participaci?n en el capital social u otro t?tulo.}

\noindent{\footnotesize\itshape \textcolor{profublue}{\textbf{Método de evaluación:}} CompareMeaning}

\noindent{\footnotesize\itshape \textcolor{profublue}{\textbf{Calificación aprobatoria:}} 50}

\noindent{\footnotesize\itshape \textcolor{profublue}{\textbf{Calificación de la respuesta:}} Pass}

\noindent{\footnotesize\itshape \textcolor{profublue}{\textbf{Explicación:}} 75/100: casi el mismo significado; The Agent answer gives all the important parts of the Expected response and adds useful details about where the definition is found in the law. It does not leave out anything important or say anything wrong. So, it matches very well and even helps by giving more context.}

\end{minipage}

\vspace{2.5em}
\noindent
\begin{minipage}{\textwidth}
\textcolor{profublue}{\textbf{Usuario (5):}}\\[4pt]
\fcolorbox{profublue!30}{profublue!3}{%
  \begin{minipage}{0.95\textwidth}
  \smallskip
  ?Qu? debe hacer una administradora si se pretende una adquisici?n de acciones menor al 5\% del capital social?
  \smallskip
  \end{minipage}}
\end{minipage}

\vspace{0.4em}
\noindent
\begin{minipage}{\textwidth}
\textcolor{profugold!80!black}{\textbf{Agente (5):}}\\[4pt]
\fcolorbox{profugold!40}{profugold!5}{%
  \begin{minipage}{0.95\textwidth}
  \smallskip
  **Respuesta breve y clara:**\par Si una administradora pretende adquirir acciones que representen menos del 5\% de su capital social, debe dar aviso a la Comisión con al menos diez días hábiles de anticipación a la operación y proporcionar la información que la Comisión determine. Además, una vez realizada la adquisición, debe informar nuevamente a la Comisión.\par **Fundamento y localización en la Ley del SAR:**\par Esto se encuentra en el Artículo 23, que establece el procedimiento para adquisiciones menores al 5\% del capital social. El texto indica explícitamente la obligación de aviso previo y posterior a la Comisión.\par Referencia: Ley del SAR [Páginas 17]
  \smallskip
  \end{minipage}}
\end{minipage}

\vspace{0.35em}
\noindent
\begin{minipage}{0.95\textwidth}
\setlength{\parskip}{0.25em}
\noindent{\footnotesize\itshape \textcolor{profublue}{\textbf{Respuesta esperada:}} Dar aviso a la Comisi?n con diez d?as h?biles de anticipaci?n y, una vez efectuada, hacerla del conocimiento de la Comisi?n.}

\noindent{\footnotesize\itshape \textcolor{profublue}{\textbf{Método de evaluación:}} CompareMeaning}

\noindent{\footnotesize\itshape \textcolor{profublue}{\textbf{Calificación aprobatoria:}} 50}

\noindent{\footnotesize\itshape \textcolor{profublue}{\textbf{Calificación de la respuesta:}} Pass}

\noindent{\footnotesize\itshape \textcolor{profublue}{\textbf{Explicación:}} 75/100: casi el mismo significado; The Agent answer gives all the important information the Expected response has, like telling the Commission before and after buying shares under 5\%. It also adds helpful extra details and legal references. Nothing is missing or wrong, so the Agent answer is very good but a bit more detailed than needed.}

\end{minipage}

\vspace{2.5em}
\noindent
\begin{minipage}{\textwidth}
\textcolor{profublue}{\textbf{Usuario (6):}}\\[4pt]
\fcolorbox{profublue!30}{profublue!3}{%
  \begin{minipage}{0.95\textwidth}
  \smallskip
  ?Qu? son las Empresas Operadoras en el contexto de la Ley del SAR?
  \smallskip
  \end{minipage}}
\end{minipage}

\vspace{0.4em}
\noindent
\begin{minipage}{\textwidth}
\textcolor{profugold!80!black}{\textbf{Agente (6):}}\\[4pt]
\fcolorbox{profugold!40}{profugold!5}{%
  \begin{minipage}{0.95\textwidth}
  \smallskip
  **Respuesta breve y clara:**\par Las Empresas Operadoras, en el contexto de la Ley del SAR, son sociedades anónimas de capital variable que cuentan con una concesión otorgada por el Gobierno Federal para operar la Base de Datos Nacional SAR. Su objeto exclusivo es administrar dicha base, promover procesos ordenados de elección de administradora y retiro de recursos por los trabajadores, y coadyuvar en la localización de trabajadores para facilitar el traspaso de cuentas individuales entre administradoras.\par **Fundamento legal y localización en la Ley del SAR:**\par - Sección IV ``De las Empresas Operadoras de la Base de Datos Nacional SAR'' - Artículo 58: Define la naturaleza, requisitos y objeto exclusivo de las Empresas Operadoras.   - I. Administrar la Base de Datos Nacional SAR.   - II. Promover procesos ordenados de elección de administradora y retiro de recursos.   - III. Coadyuvar en la localización de trabajadores para el traspaso de cuentas individuales.\par Referencia: Ley del SAR [Página 34]
  \smallskip
  \end{minipage}}
\end{minipage}

\vspace{0.35em}
\noindent
\begin{minipage}{0.95\textwidth}
\setlength{\parskip}{0.25em}
\noindent{\footnotesize\itshape \textcolor{profublue}{\textbf{Respuesta esperada:}} Empresas concesionarias para operar la Base de Datos Nacional SAR.}

\noindent{\footnotesize\itshape \textcolor{profublue}{\textbf{Método de evaluación:}} CompareMeaning}

\noindent{\footnotesize\itshape \textcolor{profublue}{\textbf{Calificación aprobatoria:}} 50}

\noindent{\footnotesize\itshape \textcolor{profublue}{\textbf{Calificación de la respuesta:}} Pass}

\noindent{\footnotesize\itshape \textcolor{profublue}{\textbf{Explicación:}} 75/100: casi el mismo significado; The Expected response says Empresas Operadoras are companies that run the SAR database. The Agent answer explains this clearly and adds more details about their role and legal basis. It doesn't leave out anything important but gives more helpful information than needed. So, the Agent answer matches well but includes extra details.}

\end{minipage}

\vspace{2.5em}
\noindent
\begin{minipage}{\textwidth}
\textcolor{profublue}{\textbf{Usuario (7):}}\\[4pt]
\fcolorbox{profublue!30}{profublue!3}{%
  \begin{minipage}{0.95\textwidth}
  \smallskip
  ?Cu?l es la funci?n general del contralor normativo de una administradora?
  \smallskip
  \end{minipage}}
\end{minipage}

\vspace{0.4em}
\noindent
\begin{minipage}{\textwidth}
\textcolor{profugold!80!black}{\textbf{Agente (7):}}\\[4pt]
\fcolorbox{profugold!40}{profugold!5}{%
  \begin{minipage}{0.95\textwidth}
  \smallskip
  **Respuesta breve y didáctica:**\par La función general del contralor normativo en una administradora es vigilar que todos los funcionarios y empleados cumplan con la normatividad externa e interna aplicable. Además, debe verificar el cumplimiento del programa de autorregulación, proponer mejoras para prevenir conflictos de interés y evitar el uso indebido de información, así como recibir y analizar informes y dictámenes de auditores y comisarios.\par **Fundamento y localización en la Ley del SAR:**\par - El contralor normativo debe ser nombrado por la asamblea de accionistas y reporta únicamente al consejo de administración y a la asamblea, sin estar subordinado a otros órganos o funcionarios. - Sus funciones específicas incluyen:   - Verificar el cumplimiento del programa de autorregulación, orientado a garantizar la normatividad, la operación eficiente y la protección de los intereses de los trabajadores, evitando conflictos de interés y uso indebido de información privilegiada.   - Proponer modificaciones al programa de autorregulación para prevenir conflictos de interés y el uso indebido de información.   - Recibir y analizar informes del comisario y dictámenes de auditores externos.\par **Referencia:**   Ley del SAR [Página 18]
  \smallskip
  \end{minipage}}
\end{minipage}

\vspace{0.35em}
\noindent
\begin{minipage}{0.95\textwidth}
\setlength{\parskip}{0.25em}
\noindent{\footnotesize\itshape \textcolor{profublue}{\textbf{Respuesta esperada:}} Vigilar que funcionarios y empleados cumplan la normatividad externa e interna aplicable.}

\noindent{\footnotesize\itshape \textcolor{profublue}{\textbf{Método de evaluación:}} CompareMeaning}

\noindent{\footnotesize\itshape \textcolor{profublue}{\textbf{Calificación aprobatoria:}} 50}

\noindent{\footnotesize\itshape \textcolor{profublue}{\textbf{Calificación de la respuesta:}} Pass}

\noindent{\footnotesize\itshape \textcolor{profublue}{\textbf{Explicación:}} 75/100: casi el mismo significado; The Agent answer includes all the important points from the Expected response and adds more helpful details that are correct and related. It does not leave out anything important or say anything wrong. So, it matches very well and gives extra useful information.}

\end{minipage}

\vspace{2.5em}
\noindent
\begin{minipage}{\textwidth}
\textcolor{profublue}{\textbf{Usuario (8):}}\\[4pt]
\fcolorbox{profublue!30}{profublue!3}{%
  \begin{minipage}{0.95\textwidth}
  \smallskip
  ?Qu? es una Administradora seg?n las definiciones de la Ley del SAR?
  \smallskip
  \end{minipage}}
\end{minipage}

\vspace{0.4em}
\noindent
\begin{minipage}{\textwidth}
\textcolor{profugold!80!black}{\textbf{Agente (8):}}\\[4pt]
\fcolorbox{profugold!40}{profugold!5}{%
  \begin{minipage}{0.95\textwidth}
  \smallskip
  **Definición de Administradora según la Ley del SAR**\par Una Administradora, conforme a la Ley de los Sistemas de Ahorro para el Retiro (Ley del SAR), es una entidad financiera que se dedica de manera habitual y profesional a administrar las cuentas individuales de los trabajadores y canalizar los recursos de las subcuentas que las integran, en términos de la propia Ley. Además, las Administradoras gestionan sociedades de inversión y deben procurar una adecuada rentabilidad y seguridad en las inversiones, atendiendo exclusivamente al interés de los trabajadores.\par **Fundamento legal y localización** - Artículo 18, Capítulo III, Sección I. - Ley del SAR [Página 12].
  \smallskip
  \end{minipage}}
\end{minipage}

\vspace{0.35em}
\noindent
\begin{minipage}{0.95\textwidth}
\setlength{\parskip}{0.25em}
\noindent{\footnotesize\itshape \textcolor{profublue}{\textbf{Respuesta esperada:}} Las administradoras de fondos para el retiro (AFORE).}

\noindent{\footnotesize\itshape \textcolor{profublue}{\textbf{Método de evaluación:}} CompareMeaning}

\noindent{\footnotesize\itshape \textcolor{profublue}{\textbf{Calificación aprobatoria:}} 50}

\noindent{\footnotesize\itshape \textcolor{profublue}{\textbf{Calificación de la respuesta:}} Pass}

\noindent{\footnotesize\itshape \textcolor{profublue}{\textbf{Explicación:}} 75/100: casi el mismo significado; The Expected response gives a very short answer about what an Administradora is. The Agent answer explains it in much more detail, adding helpful information about how it works and legal rules. The extra details are correct and useful but go beyond what was asked. So, the Agent answer matches well but includes more than needed.}

\end{minipage}

\vspace{2.5em}
\noindent
\begin{minipage}{\textwidth}
\textcolor{profublue}{\textbf{Usuario (9):}}\\[4pt]
\fcolorbox{profublue!30}{profublue!3}{%
  \begin{minipage}{0.95\textwidth}
  \smallskip
  ?Cu?l es el l?mite general de participaci?n de mercado que puede tener una administradora?
  \smallskip
  \end{minipage}}
\end{minipage}

\vspace{0.4em}
\noindent
\begin{minipage}{\textwidth}
\textcolor{profugold!80!black}{\textbf{Agente (9):}}\\[4pt]
\fcolorbox{profugold!40}{profugold!5}{%
  \begin{minipage}{0.95\textwidth}
  \smallskip
  **Respuesta breve y clara:** El límite general de participación de mercado que puede tener una administradora en los sistemas de ahorro para el retiro es del 17\%. Sin embargo, la Comisión Nacional del Sistema de Ahorro para el Retiro puede autorizar un límite mayor si considera que no se perjudican los intereses de los trabajadores.\par **Fundamento legal y localización en el texto:** Esto se encuentra en el Artículo Décimo Séptimo, que establece dicho límite y la posibilidad de autorización excepcional por parte de la Comisión.\par Referencia: Ley del SAR [Página 71]
  \smallskip
  \end{minipage}}
\end{minipage}

\vspace{0.35em}
\noindent
\begin{minipage}{0.95\textwidth}
\setlength{\parskip}{0.25em}
\noindent{\footnotesize\itshape \textcolor{profublue}{\textbf{Respuesta esperada:}} No m?s del 20\% de participaci?n en el mercado de los sistemas de ahorro para el retiro, salvo autorizaci?n de un l?mite mayor.}

\noindent{\footnotesize\itshape \textcolor{profublue}{\textbf{Método de evaluación:}} CompareMeaning}

\noindent{\footnotesize\itshape \textcolor{profublue}{\textbf{Calificación aprobatoria:}} 50}

\noindent{\footnotesize\itshape \textcolor{profublue}{\textbf{Calificación de la respuesta:}} Pass}

\noindent{\footnotesize\itshape \textcolor{profublue}{\textbf{Explicación:}} 50/100: algún significado compartido; The Agent answer gives a different percentage (17\%) than the Expected response (20\%) for the market limit. Both say a higher limit can be allowed, but the wrong number is an important mistake. The Agent answer adds some useful legal details, but the main fact about the limit is wrong.}

\end{minipage}

\vspace{2.5em}
\noindent
\begin{minipage}{\textwidth}
\textcolor{profublue}{\textbf{Usuario (10):}}\\[4pt]
\fcolorbox{profublue!30}{profublue!3}{%
  \begin{minipage}{0.95\textwidth}
  \smallskip
  ?Qu? es la Base de Datos Nacional SAR?
  \smallskip
  \end{minipage}}
\end{minipage}

\vspace{0.4em}
\noindent
\begin{minipage}{\textwidth}
\textcolor{profugold!80!black}{\textbf{Agente (10):}}\\[4pt]
\fcolorbox{profugold!40}{profugold!5}{%
  \begin{minipage}{0.95\textwidth}
  \smallskip
  **Respuesta breve y clara:** La Base de Datos Nacional SAR es el sistema que reúne la información individual de cada trabajador relacionada con los sistemas de ahorro para el retiro. Esta base contiene datos sobre cada trabajador y el registro de la administradora o institución de crédito en la que está afiliado.\par **Fundamento legal y localización:** La definición se encuentra en el Artículo 3o, fracción II, de la Ley de los Sistemas de Ahorro para el Retiro, en las páginas 1 y 2 del texto oficial.\par **Referencia:**   Ley del SAR [Páginas 1--2]
  \smallskip
  \end{minipage}}
\end{minipage}

\vspace{0.35em}
\noindent
\begin{minipage}{0.95\textwidth}
\setlength{\parskip}{0.25em}
\noindent{\footnotesize\itshape \textcolor{profublue}{\textbf{Respuesta esperada:}} La base conformada por informaci?n de los sistemas de ahorro para el retiro, con informaci?n individual de cada trabajador y el registro de su administradora o instituci?n de cr?dito.}

\noindent{\footnotesize\itshape \textcolor{profublue}{\textbf{Método de evaluación:}} CompareMeaning}

\noindent{\footnotesize\itshape \textcolor{profublue}{\textbf{Calificación aprobatoria:}} 50}

\noindent{\footnotesize\itshape \textcolor{profublue}{\textbf{Calificación de la respuesta:}} Pass}

\noindent{\footnotesize\itshape \textcolor{profublue}{\textbf{Explicación:}} 75/100: casi el mismo significado; The Agent answer gives all the important information about the database that the Expected response has. It also adds helpful details about the law that defines it. There are no mistakes or missing parts. So, the Agent answer matches very well and adds useful extra information.}

\end{minipage}

\vspace{2.5em}
\noindent
\begin{minipage}{\textwidth}
\textcolor{profublue}{\textbf{Usuario (11):}}\\[4pt]
\fcolorbox{profublue!30}{profublue!3}{%
  \begin{minipage}{0.95\textwidth}
  \smallskip
  ?Qui?n preside la Junta de Gobierno de la Comisi?n?
  \smallskip
  \end{minipage}}
\end{minipage}

\vspace{0.4em}
\noindent
\begin{minipage}{\textwidth}
\textcolor{profugold!80!black}{\textbf{Agente (11):}}\\[4pt]
\fcolorbox{profugold!40}{profugold!5}{%
  \begin{minipage}{0.95\textwidth}
  \smallskip
  Respuesta breve: La Junta de Gobierno de la Comisión es presidida por el Secretario de Hacienda y Crédito Público. En caso de ausencia, lo suplirá el Presidente de la Comisión.\par Fundamento: Esto se establece en el Artículo 7o de la Ley del SAR, donde se indica la integración y presidencia de la Junta de Gobierno.\par Referencia: Ley del SAR [Páginas 5--6].
  \smallskip
  \end{minipage}}
\end{minipage}

\vspace{0.35em}
\noindent
\begin{minipage}{0.95\textwidth}
\setlength{\parskip}{0.25em}
\noindent{\footnotesize\itshape \textcolor{profublue}{\textbf{Respuesta esperada:}} El Secretario de Hacienda y Cr?dito P?blico.}

\noindent{\footnotesize\itshape \textcolor{profublue}{\textbf{Método de evaluación:}} CompareMeaning}

\noindent{\footnotesize\itshape \textcolor{profublue}{\textbf{Calificación aprobatoria:}} 50}

\noindent{\footnotesize\itshape \textcolor{profublue}{\textbf{Calificación de la respuesta:}} Pass}

\noindent{\footnotesize\itshape \textcolor{profublue}{\textbf{Explicación:}} 75/100: casi el mismo significado; The Agent answer gives the main fact that the Secretary of Hacienda and Public Credit leads the board, just like the Expected response. It also adds extra helpful details about who takes over if the Secretary is absent and mentions the law that says this. These extra details don't change the main answer but add more information than needed.}

\end{minipage}

\vspace{2.5em}
\noindent
\begin{minipage}{\textwidth}
\textcolor{profublue}{\textbf{Usuario (12):}}\\[4pt]
\fcolorbox{profublue!30}{profublue!3}{%
  \begin{minipage}{0.95\textwidth}
  \smallskip
  ?Qu? diferencia marca la Ley entre Trabajador Afiliado y Trabajador no Afiliado?
  \smallskip
  \end{minipage}}
\end{minipage}

\vspace{0.4em}
\noindent
\begin{minipage}{\textwidth}
\textcolor{profugold!80!black}{\textbf{Agente (12):}}\\[4pt]
\fcolorbox{profugold!40}{profugold!5}{%
  \begin{minipage}{0.95\textwidth}
  \smallskip
  **Diferencia entre Trabajador Afiliado y Trabajador no Afiliado según la Ley del SAR**\par **Respuesta breve y didáctica:** La Ley del SAR distingue entre dos tipos de trabajadores en función de su inscripción ante el Instituto Mexicano del Seguro Social (IMSS):\par - **Trabajador Afiliado:** Es aquel que está inscrito en el IMSS. - **Trabajador no Afiliado:** Es aquel que no está inscrito en el IMSS.\par Ambos pueden tener derecho a la apertura de una cuenta individual en los términos de la Ley, pero la diferencia principal radica en su relación formal con el IMSS.\par **Fundamento legal y localización en el texto:** - **Trabajador Afiliado:** Definido como ``los trabajadores inscritos en el Instituto Mexicano del Seguro Social''. - **Trabajador no Afiliado:** Definido como ``los trabajadores que no se encuentren inscritos en el Instituto Mexicano del Seguro Social''.\par Estas definiciones se encuentran en las fracciones XIII y XIII bis del texto, alrededor de la Página 3.\par **Referencia:**   Ley del SAR [Página 3]
  \smallskip
  \end{minipage}}
\end{minipage}

\vspace{0.35em}
\noindent
\begin{minipage}{0.95\textwidth}
\setlength{\parskip}{0.25em}
\noindent{\footnotesize\itshape \textcolor{profublue}{\textbf{Respuesta esperada:}} El afiliado est? inscrito en el IMSS; el no afiliado no se encuentra inscrito en el IMSS.}

\noindent{\footnotesize\itshape \textcolor{profublue}{\textbf{Método de evaluación:}} CompareMeaning}

\noindent{\footnotesize\itshape \textcolor{profublue}{\textbf{Calificación aprobatoria:}} 50}

\noindent{\footnotesize\itshape \textcolor{profublue}{\textbf{Calificación de la respuesta:}} Pass}

\noindent{\footnotesize\itshape \textcolor{profublue}{\textbf{Explicación:}} 75/100: casi el mismo significado; The agent answer correctly explains that an affiliated worker is registered with the IMSS and a non-affiliated worker is not, just like the expected response. It also adds useful extra details about the law and rights without leaving out anything important or giving wrong information. So, the agent answer matches very well and helps you understand more.}

\end{minipage}

\vspace{2.5em}
\noindent
\begin{minipage}{\textwidth}
\textcolor{profublue}{\textbf{Usuario (13):}}\\[4pt]
\fcolorbox{profublue!30}{profublue!3}{%
  \begin{minipage}{0.95\textwidth}
  \smallskip
  ?Por qu? actos u omisiones responden directamente las administradoras respecto de sociedades de inversi?n que operen?
  \smallskip
  \end{minipage}}
\end{minipage}

\vspace{0.4em}
\noindent
\begin{minipage}{\textwidth}
\textcolor{profugold!80!black}{\textbf{Agente (13):}}\\[4pt]
\fcolorbox{profugold!40}{profugold!5}{%
  \begin{minipage}{0.95\textwidth}
  \smallskip
  **Respuesta breve y clara:**\par Las administradoras responden directamente por todos los actos, omisiones y operaciones que realicen las sociedades de inversión que operen, siempre que estén relacionados con su participación en los sistemas de ahorro para el retiro. Además, son responsables de los actos realizados por sus consejeros, directivos y empleados, así como por los de los consejeros y directivos de las sociedades de inversión que administren, en el cumplimiento de sus funciones relativas a los sistemas de ahorro para el retiro y la operación de la administradora y sociedades de inversión. Si estos actos son dolosos y causan una afectación patrimonial a los trabajadores, la administradora está obligada a reparar el daño causado.\par **Fundamento legal y localización:**\par - Artículo 35: Responsabilidad directa por actos, omisiones y operaciones de las sociedades de inversión que operen. - Artículo 36: Responsabilidad por actos de consejeros, directivos y empleados, y obligación de reparar daños por actos dolosos.\par Estos artículos se encuentran en las páginas 19 y 20 del texto.\par **Referencia:**   Ley del SAR [Páginas 19--20]
  \smallskip
  \end{minipage}}
\end{minipage}

\vspace{0.35em}
\noindent
\begin{minipage}{0.95\textwidth}
\setlength{\parskip}{0.25em}
\noindent{\footnotesize\itshape \textcolor{profublue}{\textbf{Respuesta esperada:}} De todos los actos, omisiones y operaciones que realicen las sociedades de inversi?n con motivo de su participaci?n en los sistemas.}

\noindent{\footnotesize\itshape \textcolor{profublue}{\textbf{Método de evaluación:}} CompareMeaning}

\noindent{\footnotesize\itshape \textcolor{profublue}{\textbf{Calificación aprobatoria:}} 50}

\noindent{\footnotesize\itshape \textcolor{profublue}{\textbf{Calificación de la respuesta:}} Pass}

\noindent{\footnotesize\itshape \textcolor{profublue}{\textbf{Explicación:}} 75/100: casi el mismo significado; The agent answer includes all the important points from the expected response and adds more helpful details about who else is responsible and legal rules. There are no mistakes, but it gives more information than needed. So, it mostly matches but has extra details.}

\end{minipage}

\vspace{2.5em}
\noindent
\begin{minipage}{\textwidth}
\textcolor{profublue}{\textbf{Usuario (14):}}\\[4pt]
\fcolorbox{profublue!30}{profublue!3}{%
  \begin{minipage}{0.95\textwidth}
  \smallskip
  ?Puede un gobierno extranjero participar en el capital social de una AFORE?
  \smallskip
  \end{minipage}}
\end{minipage}

\vspace{0.4em}
\noindent
\begin{minipage}{\textwidth}
\textcolor{profugold!80!black}{\textbf{Agente (14):}}\\[4pt]
\fcolorbox{profugold!40}{profugold!5}{%
  \begin{minipage}{0.95\textwidth}
  \smallskip
  Respuesta breve: No, en términos generales, los gobiernos extranjeros no pueden participar, directa o indirectamente, en el capital social de una AFORE. Sin embargo, la Ley del SAR contempla excepciones específicas a esta prohibición, aunque no se detallan en el fragmento consultado.\par Fundamento: Esta restricción se encuentra en el Artículo 21 de la Ley del SAR, el cual establece que los gobiernos extranjeros no podrán participar en el capital social de las administradoras de fondos para el retiro (AFORE), salvo en los casos que la propia Ley señale.\par Referencia: Ley del SAR [Páginas 14--15].
  \smallskip
  \end{minipage}}
\end{minipage}

\vspace{0.35em}
\noindent
\begin{minipage}{0.95\textwidth}
\setlength{\parskip}{0.25em}
\noindent{\footnotesize\itshape \textcolor{profublue}{\textbf{Respuesta esperada:}} En general no, salvo excepciones como medidas prudenciales temporales o participaciones autorizadas que impliquen control a trav?s de personas morales oficiales, o participaciones indirectas sin control.}

\noindent{\footnotesize\itshape \textcolor{profublue}{\textbf{Método de evaluación:}} CompareMeaning}

\noindent{\footnotesize\itshape \textcolor{profublue}{\textbf{Calificación aprobatoria:}} 50}

\noindent{\footnotesize\itshape \textcolor{profublue}{\textbf{Calificación de la respuesta:}} Pass}

\noindent{\footnotesize\itshape \textcolor{profublue}{\textbf{Explicación:}} 75/100: casi el mismo significado; The Agent answer correctly says foreign governments usually can't own part of an AFORE, just like the Expected response. It also says there are exceptions but doesn't explain what they are, while the Expected response gives examples of these exceptions. The Agent answer adds a law article number, which is helpful. There are no wrong or unrelated details. So, the Agent answer is mostly complete but misses some specific exception examples.}

\end{minipage}

\vspace{2.5em}
\noindent
\begin{minipage}{\textwidth}
\textcolor{profublue}{\textbf{Usuario (15):}}\\[4pt]
\fcolorbox{profublue!30}{profublue!3}{%
  \begin{minipage}{0.95\textwidth}
  \smallskip
  ?Pueden las administradoras cobrar cuotas fijas por la administraci?n de las cuentas?
  \smallskip
  \end{minipage}}
\end{minipage}

\vspace{0.4em}
\noindent
\begin{minipage}{\textwidth}
\textcolor{profugold!80!black}{\textbf{Agente (15):}}\\[4pt]
\fcolorbox{profugold!40}{profugold!5}{%
  \begin{minipage}{0.95\textwidth}
  \smallskip
  Respuesta breve y clara: No, las administradoras no pueden cobrar cuotas fijas por la administración de las cuentas individuales. Únicamente pueden cobrar comisiones como porcentaje sobre el valor de los activos administrados. Las cuotas fijas solo pueden cobrarse por servicios específicos señalados en el reglamento de la Ley, y nunca por la administración de las cuentas.\par Fundamento: Esto se establece en el Artículo 37, donde se indica que las comisiones por administración de cuentas individuales solo pueden cobrarse como porcentaje sobre el valor de los activos administrados, y que las cuotas fijas únicamente pueden aplicarse a servicios señalados en el reglamento, excluyendo la administración de cuentas.\par Referencia: Ley del SAR [Páginas 19--20]
  \smallskip
  \end{minipage}}
\end{minipage}

\vspace{0.35em}
\noindent
\begin{minipage}{0.95\textwidth}
\setlength{\parskip}{0.25em}
\noindent{\footnotesize\itshape \textcolor{profublue}{\textbf{Respuesta esperada:}} No; solo pueden cobrar cuotas fijas por servicios se?alados en el reglamento y en ning?n caso por la administraci?n de las cuentas.}

\noindent{\footnotesize\itshape \textcolor{profublue}{\textbf{Método de evaluación:}} CompareMeaning}

\noindent{\footnotesize\itshape \textcolor{profublue}{\textbf{Calificación aprobatoria:}} 50}

\noindent{\footnotesize\itshape \textcolor{profublue}{\textbf{Calificación de la respuesta:}} Pass}

\noindent{\footnotesize\itshape \textcolor{profublue}{\textbf{Explicación:}} 75/100: casi el mismo significado; The Agent answer correctly says that fixed fees can't be charged for managing accounts, only for certain services. It matches the important points in the Expected response and adds useful extra details without leaving anything out or adding wrong information.}

\end{minipage}

\vspace{2.5em}
\noindent
\begin{minipage}{\textwidth}
\textcolor{profublue}{\textbf{Usuario (16):}}\\[4pt]
\fcolorbox{profublue!30}{profublue!3}{%
  \begin{minipage}{0.95\textwidth}
  \smallskip
  ?Qu? qu?rum se requiere para sesionar en Junta de Gobierno?
  \smallskip
  \end{minipage}}
\end{minipage}

\vspace{0.4em}
\noindent
\begin{minipage}{\textwidth}
\textcolor{profugold!80!black}{\textbf{Agente (16):}}\\[4pt]
\fcolorbox{profugold!40}{profugold!5}{%
  \begin{minipage}{0.95\textwidth}
  \smallskip
  **Respuesta breve y clara:**   Para que la Junta de Gobierno pueda sesionar válidamente, se requiere la presencia de al menos nueve de sus miembros. Las resoluciones se toman por mayoría de votos de los presentes. En caso de empate, el Presidente de la Junta de Gobierno tiene voto de calidad.\par **Fundamento y localización:**   Esta disposición se encuentra en el Artículo 9o de la Ley del SAR, en la página 7 del texto consultado.\par **Referencia:**   Ley del SAR [Página 7]
  \smallskip
  \end{minipage}}
\end{minipage}

\vspace{0.35em}
\noindent
\begin{minipage}{0.95\textwidth}
\setlength{\parskip}{0.25em}
\noindent{\footnotesize\itshape \textcolor{profublue}{\textbf{Respuesta esperada:}} La presencia de nueve de sus miembros.}

\noindent{\footnotesize\itshape \textcolor{profublue}{\textbf{Método de evaluación:}} CompareMeaning}

\noindent{\footnotesize\itshape \textcolor{profublue}{\textbf{Calificación aprobatoria:}} 50}

\noindent{\footnotesize\itshape \textcolor{profublue}{\textbf{Calificación de la respuesta:}} Pass}

\noindent{\footnotesize\itshape \textcolor{profublue}{\textbf{Explicación:}} 75/100: casi el mismo significado; The Expected response says nine members must be present to hold a meeting. The Agent answer says the same and also explains how decisions are made and where the rule comes from. The extra details are correct and helpful but not needed to answer the question. So, the Agent answer matches well but has some extra information.}

\end{minipage}

\vspace{2.5em}
\noindent
\begin{minipage}{\textwidth}
\textcolor{profublue}{\textbf{Usuario (17):}}\\[4pt]
\fcolorbox{profublue!30}{profublue!3}{%
  \begin{minipage}{0.95\textwidth}
  \smallskip
  ?Qu? facultad tiene la Comisi?n respecto a la supervisi?n de participantes?
  \smallskip
  \end{minipage}}
\end{minipage}

\vspace{0.4em}
\noindent
\begin{minipage}{\textwidth}
\textcolor{profugold!80!black}{\textbf{Agente (17):}}\\[4pt]
\fcolorbox{profugold!40}{profugold!5}{%
  \begin{minipage}{0.95\textwidth}
  \smallskip
  Facultades de la Comisión en Supervisión de Participantes\par Respuesta breve y clara: La Comisión tiene la facultad de supervisar a los participantes de los sistemas de ahorro para el retiro mediante inspección, vigilancia, prevención y corrección. Puede practicar visitas de inspección, requerir información y documentación, y asegurar medios que contengan información relevante. El objetivo es evaluar riesgos, sistemas de control y la calidad de la administración, procurando que los participantes mantengan liquidez, solvencia y estabilidad, y que cumplan con las disposiciones legales y buenas prácticas del mercado financiero.\par Fundamento y localización: - Artículo 89: Define el alcance de la supervisión, incluyendo inspección, vigilancia, prevención y corrección. - Artículo 90: Enumera las facultades específicas de la Comisión, como realizar visitas de inspección, requerir información y asegurar documentación.\par Referencia: Ley del SAR [Páginas 49--50].
  \smallskip
  \end{minipage}}
\end{minipage}

\vspace{0.35em}
\noindent
\begin{minipage}{0.95\textwidth}
\setlength{\parskip}{0.25em}
\noindent{\footnotesize\itshape \textcolor{profublue}{\textbf{Respuesta esperada:}} Realizar la supervisi?n de los participantes; en instituciones de cr?dito, solo respecto a su participaci?n en los sistemas de ahorro para el retiro.}

\noindent{\footnotesize\itshape \textcolor{profublue}{\textbf{Método de evaluación:}} CompareMeaning}

\noindent{\footnotesize\itshape \textcolor{profublue}{\textbf{Calificación aprobatoria:}} 50}

\noindent{\footnotesize\itshape \textcolor{profublue}{\textbf{Calificación de la respuesta:}} Pass}

\noindent{\footnotesize\itshape \textcolor{profublue}{\textbf{Explicación:}} 50/100: algún significado compartido; The Expected response says the Commission supervises participants and, for credit institutions, only their role in retirement savings. The Agent answer explains many ways the Commission supervises but leaves out the important detail about credit institutions. So, the Agent answer is mostly right but misses this key point.}

\end{minipage}

\vspace{2.5em}
\noindent
\begin{minipage}{\textwidth}
\textcolor{profublue}{\textbf{Usuario (18):}}\\[4pt]
\fcolorbox{profublue!30}{profublue!3}{%
  \begin{minipage}{0.95\textwidth}
  \smallskip
  ?Qu? son los Sistemas de Ahorro para el Retiro seg?n la Ley?
  \smallskip
  \end{minipage}}
\end{minipage}

\vspace{0.4em}
\noindent
\begin{minipage}{\textwidth}
\textcolor{profugold!80!black}{\textbf{Agente (18):}}\\[4pt]
\fcolorbox{profugold!40}{profugold!5}{%
  \begin{minipage}{0.95\textwidth}
  \smallskip
  **Respuesta breve y didáctica**\par Los Sistemas de Ahorro para el Retiro, según la Ley del SAR, son mecanismos regulados por el Estado mexicano cuyo objetivo es organizar y supervisar el ahorro que los trabajadores realizan para su retiro. La Ley establece cómo funcionan estos sistemas, quiénes participan y cómo se coordinan, regulan y vigilan. Incluyen las cuentas individuales de los trabajadores, las administradoras de fondos para el retiro y las instituciones relacionadas, todo bajo la supervisión de la Comisión Nacional del Sistema de Ahorro para el Retiro.\par **Fundamento legal y localización**\par - El objeto de la Ley y la definición de los sistemas de ahorro para el retiro se encuentran en el Capítulo I, Disposiciones Preliminares, especialmente en los Artículos 1o, 2o y 3o. - El texto relevante está ubicado en las páginas 1 y 2 del documento.\par **Referencia** Ley del SAR [Páginas 1--2]
  \smallskip
  \end{minipage}}
\end{minipage}

\vspace{0.35em}
\noindent
\begin{minipage}{0.95\textwidth}
\setlength{\parskip}{0.25em}
\noindent{\footnotesize\itshape \textcolor{profublue}{\textbf{Respuesta esperada:}} Sistemas regulados por leyes de seguridad social donde aportaciones de trabajadores, patrones y Estado se manejan en cuentas individuales para previsi?n social y pensiones.}

\noindent{\footnotesize\itshape \textcolor{profublue}{\textbf{Método de evaluación:}} CompareMeaning}

\noindent{\footnotesize\itshape \textcolor{profublue}{\textbf{Calificación aprobatoria:}} 50}

\noindent{\footnotesize\itshape \textcolor{profublue}{\textbf{Calificación de la respuesta:}} Pass}

\noindent{\footnotesize\itshape \textcolor{profublue}{\textbf{Explicación:}} 75/100: casi el mismo significado; The Agent answer explains the retirement savings systems well and matches most of the important points in the Expected response. But it leaves out that the money comes from workers, employers, and the government, which is an important detail. It also adds extra info about legal articles that wasn't asked for. There are no wrong facts, just some missing details and extra info.}

\end{minipage}

\vspace{2.5em}
\noindent
\begin{minipage}{\textwidth}
\textcolor{profublue}{\textbf{Usuario (19):}}\\[4pt]
\fcolorbox{profublue!30}{profublue!3}{%
  \begin{minipage}{0.95\textwidth}
  \smallskip
  ?C?mo deben cobrarse las comisiones por administraci?n de cuentas individuales para promover mayor Rendimiento Neto?
  \smallskip
  \end{minipage}}
\end{minipage}

\vspace{0.4em}
\noindent
\begin{minipage}{\textwidth}
\textcolor{profugold!80!black}{\textbf{Agente (19):}}\\[4pt]
\fcolorbox{profugold!40}{profugold!5}{%
  \begin{minipage}{0.95\textwidth}
  \smallskip
  **Respuesta breve y clara:**\par Las comisiones por administración de cuentas individuales deben cobrarse únicamente como un porcentaje sobre el valor de los activos administrados. Esto significa que no pueden cobrarse cuotas fijas por la administración de las cuentas individuales, sino solo por servicios específicos que se señalen en el reglamento de la Ley. El objetivo de esta disposición es promover un mayor Rendimiento Neto a favor de los trabajadores.\par **Fundamento legal y localización en la Ley del SAR:**\par - Artículo 37: Establece que las comisiones por administración de cuentas individuales solo pueden cobrarse como porcentaje sobre el valor de los activos administrados, y prohíbe cuotas fijas por la administración de cuentas. - Ubicación: Ley del SAR [Páginas 19--20].
  \smallskip
  \end{minipage}}
\end{minipage}

\vspace{0.35em}
\noindent
\begin{minipage}{0.95\textwidth}
\setlength{\parskip}{0.25em}
\noindent{\footnotesize\itshape \textcolor{profublue}{\textbf{Respuesta esperada:}} Solo como un porcentaje sobre el valor de los activos administrados.}

\noindent{\footnotesize\itshape \textcolor{profublue}{\textbf{Método de evaluación:}} CompareMeaning}

\noindent{\footnotesize\itshape \textcolor{profublue}{\textbf{Calificación aprobatoria:}} 50}

\noindent{\footnotesize\itshape \textcolor{profublue}{\textbf{Calificación de la respuesta:}} Pass}

\noindent{\footnotesize\itshape \textcolor{profublue}{\textbf{Explicación:}} 75/100: casi el mismo significado; The Agent answer gives the main point the Expected response has and adds helpful extra details about the law and why the rule exists. It doesn't leave out anything important or say anything wrong. The extra information is useful but not needed to answer the question, so the answer is mostly a good match with some extra details.}

\end{minipage}

\vspace{2.5em}
\noindent
\begin{minipage}{\textwidth}
\textcolor{profublue}{\textbf{Usuario (20):}}\\[4pt]
\fcolorbox{profublue!30}{profublue!3}{%
  \begin{minipage}{0.95\textwidth}
  \smallskip
  ?Qu? es el V?nculo Laboral seg?n la Ley del SAR?
  \smallskip
  \end{minipage}}
\end{minipage}

\vspace{0.4em}
\noindent
\begin{minipage}{\textwidth}
\textcolor{profugold!80!black}{\textbf{Agente (20):}}\\[4pt]
\fcolorbox{profugold!40}{profugold!5}{%
  \begin{minipage}{0.95\textwidth}
  \smallskip
  **Respuesta breve y clara:**\par El Vínculo Laboral, según la Ley de los Sistemas de Ahorro para el Retiro (Ley del SAR), se refiere a la prestación de servicios subordinados conforme a lo dispuesto por la Ley Federal del Trabajo, o bien a la prestación de servicios profesionales.\par **Fundamento legal y localización:**\par - Definición encontrada en la fracción XIV de las definiciones de la Ley del SAR. - Se ubica en el rango de las páginas 2 a 3 del texto.\par **Referencia:**   Ley del SAR [Páginas 2--3]
  \smallskip
  \end{minipage}}
\end{minipage}

\vspace{0.35em}
\noindent
\begin{minipage}{0.95\textwidth}
\setlength{\parskip}{0.25em}
\noindent{\footnotesize\itshape \textcolor{profublue}{\textbf{Respuesta esperada:}} La prestaci?n de servicios subordinados conforme a la LFT o la prestaci?n de servicios profesionales.}

\noindent{\footnotesize\itshape \textcolor{profublue}{\textbf{Método de evaluación:}} CompareMeaning}

\noindent{\footnotesize\itshape \textcolor{profublue}{\textbf{Calificación aprobatoria:}} 50}

\noindent{\footnotesize\itshape \textcolor{profublue}{\textbf{Calificación de la respuesta:}} Pass}

\noindent{\footnotesize\itshape \textcolor{profublue}{\textbf{Explicación:}} 75/100: casi el mismo significado; The Agent answer gives all the important information the Expected response has about the Vínculo Laboral. It also adds helpful details about where to find this in the law, which is useful and correct. There are no mistakes or missing parts, so the Agent answer matches very well.}

\end{minipage}

\vspace{2.5em}
\noindent
\begin{minipage}{\textwidth}
\textcolor{profublue}{\textbf{Usuario (21):}}\\[4pt]
\fcolorbox{profublue!30}{profublue!3}{%
  \begin{minipage}{0.95\textwidth}
  \smallskip
  Si el capital de una administradora se reduce por debajo del m?nimo exigido, ?qu? debe hacer?
  \smallskip
  \end{minipage}}
\end{minipage}

\vspace{0.4em}
\noindent
\begin{minipage}{\textwidth}
\textcolor{profugold!80!black}{\textbf{Agente (21):}}\\[4pt]
\fcolorbox{profugold!40}{profugold!5}{%
  \begin{minipage}{0.95\textwidth}
  \smallskip
  Respuesta breve: Si el capital de una administradora se reduce por debajo del mínimo exigido, está obligada a reconstituirlo dentro del plazo que determine la Comisión, el cual no podrá exceder de 45 días naturales.\par Fundamento: Esto se establece en el Artículo 24 de la Ley de los Sistemas de Ahorro para el Retiro, que señala la obligación de mantener el capital mínimo y el plazo máximo para reconstituirlo en caso de reducción.\par Referencia: Ley del SAR [Páginas 17]
  \smallskip
  \end{minipage}}
\end{minipage}

\vspace{0.35em}
\noindent
\begin{minipage}{0.95\textwidth}
\setlength{\parskip}{0.25em}
\noindent{\footnotesize\itshape \textcolor{profublue}{\textbf{Respuesta esperada:}} Reconstituirlo dentro del plazo que determine la Comisi?n, que no podr? exceder de 45 d?as naturales.}

\noindent{\footnotesize\itshape \textcolor{profublue}{\textbf{Método de evaluación:}} CompareMeaning}

\noindent{\footnotesize\itshape \textcolor{profublue}{\textbf{Calificación aprobatoria:}} 50}

\noindent{\footnotesize\itshape \textcolor{profublue}{\textbf{Calificación de la respuesta:}} Pass}

\noindent{\footnotesize\itshape \textcolor{profublue}{\textbf{Explicación:}} 75/100: casi el mismo significado; The Agent answer gives all the important information the Expected response has and adds helpful details about the law behind it. The extra info is correct and useful but not required. So, the Agent answer matches well but includes some extra details.}

\end{minipage}

\vspace{2.5em}
\noindent
\begin{minipage}{\textwidth}
\textcolor{profublue}{\textbf{Usuario (22):}}\\[4pt]
\fcolorbox{profublue!30}{profublue!3}{%
  \begin{minipage}{0.95\textwidth}
  \smallskip
  ?C?mo se integran los cinco vocales restantes de la Junta de Gobierno?
  \smallskip
  \end{minipage}}
\end{minipage}

\vspace{0.4em}
\noindent
\begin{minipage}{\textwidth}
\textcolor{profugold!80!black}{\textbf{Agente (22):}}\\[4pt]
\fcolorbox{profugold!40}{profugold!5}{%
  \begin{minipage}{0.95\textwidth}
  \smallskip
  **Respuesta breve y clara:** Los cinco vocales restantes de la Junta de Gobierno son designados por el Secretario de Hacienda y Crédito Público. De estos, cuatro deben ser representantes de las organizaciones nacionales de trabajadores y uno de las organizaciones nacionales de patrones. Todos ellos deben formar parte del Comité Consultivo y de Vigilancia y ostentar la mayor representatividad.\par **Fundamento y localización en la Ley del SAR:** Esta integración se encuentra en el Artículo 7o, dentro de la Sección II ``De los Órganos de Gobierno''. El texto especifica que el Secretario de Hacienda y Crédito Público realiza la designación, y detalla el perfil de los representantes obreros y patronales.\par Referencia: Ley del SAR [Página 6].
  \smallskip
  \end{minipage}}
\end{minipage}

\vspace{0.35em}
\noindent
\begin{minipage}{0.95\textwidth}
\setlength{\parskip}{0.25em}
\noindent{\footnotesize\itshape \textcolor{profublue}{\textbf{Respuesta esperada:}} Los designa el Secretario de Hacienda: cuatro representantes de organizaciones nacionales de trabajadores y uno de patrones del Comit? Consultivo y de Vigilancia.}

\noindent{\footnotesize\itshape \textcolor{profublue}{\textbf{Método de evaluación:}} CompareMeaning}

\noindent{\footnotesize\itshape \textcolor{profublue}{\textbf{Calificación aprobatoria:}} 50}

\noindent{\footnotesize\itshape \textcolor{profublue}{\textbf{Calificación de la respuesta:}} Pass}

\noindent{\footnotesize\itshape \textcolor{profublue}{\textbf{Explicación:}} 75/100: casi el mismo significado; The agent answer gives all the important details the expected response has. It says who chooses the five members, how many represent workers and employers, and that they belong to a special committee. It also adds where this is written in the law, which helps understand better. There are no mistakes or missing parts, and the extra info is useful, not confusing.}

\end{minipage}

\vspace{2.5em}
\noindent
\begin{minipage}{\textwidth}
\textcolor{profublue}{\textbf{Usuario (23):}}\\[4pt]
\fcolorbox{profublue!30}{profublue!3}{%
  \begin{minipage}{0.95\textwidth}
  \smallskip
  ?A qu? Secretar?a est? desconcentrada la Comisi?n Nacional del Sistema de Ahorro para el Retiro?
  \smallskip
  \end{minipage}}
\end{minipage}

\vspace{0.4em}
\noindent
\begin{minipage}{\textwidth}
\textcolor{profugold!80!black}{\textbf{Agente (23):}}\\[4pt]
\fcolorbox{profugold!40}{profugold!5}{%
  \begin{minipage}{0.95\textwidth}
  \smallskip
  Respuesta breve: La Comisión Nacional del Sistema de Ahorro para el Retiro (CONSAR) es un órgano administrativo desconcentrado de la Secretaría de Hacienda y Crédito Público.\par Fundamento y localización: Esto se establece en el Artículo 2o de la Ley de los Sistemas de Ahorro para el Retiro, en el Capítulo I, Disposiciones Preliminares, ubicado en la Página 1 del texto oficial.\par Referencia: Ley del SAR [Página 1]
  \smallskip
  \end{minipage}}
\end{minipage}

\vspace{0.35em}
\noindent
\begin{minipage}{0.95\textwidth}
\setlength{\parskip}{0.25em}
\noindent{\footnotesize\itshape \textcolor{profublue}{\textbf{Respuesta esperada:}} A la Secretar?a de Hacienda y Cr?dito P?blico.}

\noindent{\footnotesize\itshape \textcolor{profublue}{\textbf{Método de evaluación:}} CompareMeaning}

\noindent{\footnotesize\itshape \textcolor{profublue}{\textbf{Calificación aprobatoria:}} 50}

\noindent{\footnotesize\itshape \textcolor{profublue}{\textbf{Calificación de la respuesta:}} Pass}

\noindent{\footnotesize\itshape \textcolor{profublue}{\textbf{Explicación:}} 75/100: casi el mismo significado; The Agent answer gives the correct department that the commission belongs to, just like the Expected response. It also adds helpful details about the law that supports this, which is true and useful. There are no mistakes or missing parts, so the Agent answer is complete and accurate.}

\end{minipage}

\vspace{2.5em}
\noindent
\begin{minipage}{\textwidth}
\textcolor{profublue}{\textbf{Usuario (24):}}\\[4pt]
\fcolorbox{profublue!30}{profublue!3}{%
  \begin{minipage}{0.95\textwidth}
  \smallskip
  ?Qu? comprende el t?rmino 'Leyes de Seguridad Social' en esta Ley?
  \smallskip
  \end{minipage}}
\end{minipage}

\vspace{0.4em}
\noindent
\begin{minipage}{\textwidth}
\textcolor{profugold!80!black}{\textbf{Agente (24):}}\\[4pt]
\fcolorbox{profugold!40}{profugold!5}{%
  \begin{minipage}{0.95\textwidth}
  \smallskip
  **Respuesta breve y clara:**\par En la Ley de los Sistemas de Ahorro para el Retiro, el término **'Leyes de Seguridad Social'** comprende específicamente las siguientes leyes: - La Ley del Seguro Social, - La Ley del Instituto del Fondo Nacional de la Vivienda para los Trabajadores (INFONAVIT), - La Ley del Instituto de Seguridad y Servicios Sociales de los Trabajadores del Estado (ISSSTE).\par **Fundamento y localización en la Ley:**\par Esta definición se encuentra en el **Artículo 3o, fracción VII**, dentro de las disposiciones preliminares de la Ley del SAR, y abarca el contenido de las **Páginas 1--2**.\par **Referencia:**   Ley del SAR [Páginas 1--2]
  \smallskip
  \end{minipage}}
\end{minipage}

\vspace{0.35em}
\noindent
\begin{minipage}{0.95\textwidth}
\setlength{\parskip}{0.25em}
\noindent{\footnotesize\itshape \textcolor{profublue}{\textbf{Respuesta esperada:}} Las leyes del Seguro Social, del INFONAVIT y del ISSSTE.}

\noindent{\footnotesize\itshape \textcolor{profublue}{\textbf{Método de evaluación:}} CompareMeaning}

\noindent{\footnotesize\itshape \textcolor{profublue}{\textbf{Calificación aprobatoria:}} 50}

\noindent{\footnotesize\itshape \textcolor{profublue}{\textbf{Calificación de la respuesta:}} Pass}

\noindent{\footnotesize\itshape \textcolor{profublue}{\textbf{Explicación:}} 75/100: casi el mismo significado; The Agent answer gives all the important parts the Expected response has and also adds more details about where the definition is found in the law. The extra information is correct and helpful but not needed to answer the question. So, the Agent answer matches well but includes some extra details.}

\end{minipage}

\vspace{2.5em}
\noindent
\begin{minipage}{\textwidth}
\textcolor{profublue}{\textbf{Usuario (25):}}\\[4pt]
\fcolorbox{profublue!30}{profublue!3}{%
  \begin{minipage}{0.95\textwidth}
  \smallskip
  ?Qu? prop?sito tiene el l?mite de concentraci?n de mercado previsto en el art?culo 26?
  \smallskip
  \end{minipage}}
\end{minipage}

\vspace{0.4em}
\noindent
\begin{minipage}{\textwidth}
\textcolor{profugold!80!black}{\textbf{Agente (25):}}\\[4pt]
\fcolorbox{profugold!40}{profugold!5}{%
  \begin{minipage}{0.95\textwidth}
  \smallskip
  **Respuesta breve y clara:**\par El límite de concentración de mercado previsto en el artículo 26 de la Ley de los Sistemas de Ahorro para el Retiro tiene como propósito mantener un adecuado balance y equilibrio en los sistemas de ahorro para el retiro. Específicamente, busca evitar que una sola administradora tenga una participación excesiva (más del 20\%) en el mercado, lo que podría afectar la competencia y los intereses de los trabajadores. La Comisión puede autorizar un límite mayor, siempre que no perjudique a los trabajadores.\par **Fundamento y localización en la Ley:**\par - Artículo 26: Establece el límite de concentración de mercado y su propósito. - Página 17.\par Referencia: Ley del SAR [Página 17].
  \smallskip
  \end{minipage}}
\end{minipage}

\vspace{0.35em}
\noindent
\begin{minipage}{0.95\textwidth}
\setlength{\parskip}{0.25em}
\noindent{\footnotesize\itshape \textcolor{profublue}{\textbf{Respuesta esperada:}} Mantener un balance y equilibrio adecuados en los sistemas de ahorro para el retiro y proteger los intereses de los trabajadores.}

\noindent{\footnotesize\itshape \textcolor{profublue}{\textbf{Método de evaluación:}} CompareMeaning}

\noindent{\footnotesize\itshape \textcolor{profublue}{\textbf{Calificación aprobatoria:}} 50}

\noindent{\footnotesize\itshape \textcolor{profublue}{\textbf{Calificación de la respuesta:}} Pass}

\noindent{\footnotesize\itshape \textcolor{profublue}{\textbf{Explicación:}} 75/100: casi el mismo significado; The agent answer includes all the important points from the expected response and adds useful extra details about the market limit and the Commission's role. It does not leave out anything important or add wrong information. So, it matches very well.}

\end{minipage}

\vspace{2.5em}
\noindent
\begin{minipage}{\textwidth}
\textcolor{profublue}{\textbf{Usuario (26):}}\\[4pt]
\fcolorbox{profublue!30}{profublue!3}{%
  \begin{minipage}{0.95\textwidth}
  \smallskip
  ?Se pueden autorizar aumentos de comisiones por encima del promedio del resto de las comisiones autorizadas?
  \smallskip
  \end{minipage}}
\end{minipage}

\vspace{0.4em}
\noindent
\begin{minipage}{\textwidth}
\textcolor{profugold!80!black}{\textbf{Agente (26):}}\\[4pt]
\fcolorbox{profugold!40}{profugold!5}{%
  \begin{minipage}{0.95\textwidth}
  \smallskip
  Respuesta breve: No, la Ley de los Sistemas de Ahorro para el Retiro establece expresamente que no se pueden autorizar aumentos de comisiones por encima del promedio del resto de las comisiones autorizadas. Es decir, cualquier solicitud de aumento debe respetar el promedio vigente y no excederlo.\par Fundamento: Esta restricción se encuentra en el texto que regula la autorización de comisiones por parte de la Junta de Gobierno de la Comisión. Además, la ley señala que, si una administradora no presenta su solicitud o si le es denegada, deberá cobrar la comisión más baja o el promedio del resto de las comisiones autorizadas, según corresponda.\par Referencia: Ley del SAR [Página 21]
  \smallskip
  \end{minipage}}
\end{minipage}

\vspace{0.35em}
\noindent
\begin{minipage}{0.95\textwidth}
\setlength{\parskip}{0.25em}
\noindent{\footnotesize\itshape \textcolor{profublue}{\textbf{Respuesta esperada:}} No, la Ley establece que no se podr?n autorizar aumentos por encima del promedio del resto de las comisiones autorizadas.}

\noindent{\footnotesize\itshape \textcolor{profublue}{\textbf{Método de evaluación:}} CompareMeaning}

\noindent{\footnotesize\itshape \textcolor{profublue}{\textbf{Calificación aprobatoria:}} 50}

\noindent{\footnotesize\itshape \textcolor{profublue}{\textbf{Calificación de la respuesta:}} Pass}

\noindent{\footnotesize\itshape \textcolor{profublue}{\textbf{Explicación:}} 75/100: casi el mismo significado; The Agent answer correctly says you cannot approve commission increases above the average, just like the Expected response. It also gives extra helpful details about the law and rules, which are true and related. So, the Agent answer matches well and adds useful information without mistakes.}

\end{minipage}

\vspace{2.5em}
\noindent
\begin{minipage}{\textwidth}
\textcolor{profublue}{\textbf{Usuario (27):}}\\[4pt]
\fcolorbox{profublue!30}{profublue!3}{%
  \begin{minipage}{0.95\textwidth}
  \smallskip
  ?C?mo responde la administradora por actos de sus agentes promotores?
  \smallskip
  \end{minipage}}
\end{minipage}

\vspace{0.4em}
\noindent
\begin{minipage}{\textwidth}
\textcolor{profugold!80!black}{\textbf{Agente (27):}}\\[4pt]
\fcolorbox{profugold!40}{profugold!5}{%
  \begin{minipage}{0.95\textwidth}
  \smallskip
  Respuesta breve y clara: La administradora responde por los actos de sus agentes promotores cuando estos registran o solicitan el traspaso de la cuenta individual de un trabajador sin su consentimiento, o si lo hacen mediante dolo, mala fe, documentos falsos, alterados o mediante la entrega de beneficios indebidos. Además de las sanciones a los agentes, la administradora puede ser sancionada y debe resarcir los daños y perjuicios ocasionados al trabajador afectado, devolviendo comisiones y abonando la diferencia de rendimientos en la cuenta individual.\par Fundamento y localización: - Artículo 100-A: La administradora debe resarcir daños y perjuicios al trabajador indebidamente registrado o traspasado. - Artículo 100-B: El agente promotor puede ser multado y suspendido o inhabilitado, independientemente de otras responsabilidades. - El contenido se encuentra en las Páginas 59 y siguientes.\par Referencia: Ley del SAR [Páginas 59--60]
  \smallskip
  \end{minipage}}
\end{minipage}

\vspace{0.35em}
\noindent
\begin{minipage}{0.95\textwidth}
\setlength{\parskip}{0.25em}
\noindent{\footnotesize\itshape \textcolor{profublue}{\textbf{Respuesta esperada:}} Responde directamente por actos realizados por agentes promotores, tengan relaci?n laboral, por terceros o sean independientes.}

\noindent{\footnotesize\itshape \textcolor{profublue}{\textbf{Método de evaluación:}} CompareMeaning}

\noindent{\footnotesize\itshape \textcolor{profublue}{\textbf{Calificación aprobatoria:}} 50}

\noindent{\footnotesize\itshape \textcolor{profublue}{\textbf{Calificación de la respuesta:}} Pass}

\noindent{\footnotesize\itshape \textcolor{profublue}{\textbf{Explicación:}} 50/100: algún significado compartido; The Expected response says the administrator is responsible for acts by promoter agents no matter their job status. The Agent answer gives more details about specific bad actions and punishments but doesn't say the administrator is responsible regardless of the agent's relationship. So, the Agent answer misses an important part of the answer.}

\end{minipage}

\vspace{2.5em}
\noindent
\begin{minipage}{\textwidth}
\textcolor{profublue}{\textbf{Usuario (28):}}\\[4pt]
\fcolorbox{profublue!30}{profublue!3}{%
  \begin{minipage}{0.95\textwidth}
  \smallskip
  ?C?mo se toman las resoluciones de la Junta de Gobierno?
  \smallskip
  \end{minipage}}
\end{minipage}

\vspace{0.4em}
\noindent
\begin{minipage}{\textwidth}
\textcolor{profugold!80!black}{\textbf{Agente (28):}}\\[4pt]
\fcolorbox{profugold!40}{profugold!5}{%
  \begin{minipage}{0.95\textwidth}
  \smallskip
  Respuesta breve y clara: Las resoluciones de la Junta de Gobierno se toman por mayoría de votos de los miembros presentes en la sesión. Para que haya quórum, deben estar presentes al menos nueve miembros. En caso de empate, el Presidente de la Junta de Gobierno tiene voto de calidad, es decir, su voto decide. Los acuerdos tomados son firmados por el Presidente de la Comisión para su ejecución y, en su caso, publicación.\par Fundamento: Esto se encuentra en el Artículo 9o, que establece el procedimiento para las sesiones, el quórum necesario y la forma en que se toman las resoluciones.\par Referencia: Ley del SAR [Página 7]
  \smallskip
  \end{minipage}}
\end{minipage}

\vspace{0.35em}
\noindent
\begin{minipage}{0.95\textwidth}
\setlength{\parskip}{0.25em}
\noindent{\footnotesize\itshape \textcolor{profublue}{\textbf{Respuesta esperada:}} Por mayor?a de votos de los presentes; el Presidente de la Junta tiene voto de calidad en caso de empate.}

\noindent{\footnotesize\itshape \textcolor{profublue}{\textbf{Método de evaluación:}} CompareMeaning}

\noindent{\footnotesize\itshape \textcolor{profublue}{\textbf{Calificación aprobatoria:}} 50}

\noindent{\footnotesize\itshape \textcolor{profublue}{\textbf{Calificación de la respuesta:}} Pass}

\noindent{\footnotesize\itshape \textcolor{profublue}{\textbf{Explicación:}} 75/100: casi el mismo significado; The Agent answer gives all the important points from the Expected response and adds helpful extra details like how many members must be present and what happens after decisions are made. Nothing is wrong or missing. The extra information is useful but not required, so the answer matches well but includes more than needed.}

\end{minipage}

\vspace{2.5em}
\noindent
\begin{minipage}{\textwidth}
\textcolor{profublue}{\textbf{Usuario (29):}}\\[4pt]
\fcolorbox{profublue!30}{profublue!3}{%
  \begin{minipage}{0.95\textwidth}
  \smallskip
  ?A qui?n corresponde la interpretaci?n administrativa de los preceptos de esta Ley?
  \smallskip
  \end{minipage}}
\end{minipage}

\vspace{0.4em}
\noindent
\begin{minipage}{\textwidth}
\textcolor{profugold!80!black}{\textbf{Agente (29):}}\\[4pt]
\fcolorbox{profugold!40}{profugold!5}{%
  \begin{minipage}{0.95\textwidth}
  \smallskip
  Respuesta breve: La interpretación administrativa de los preceptos de la Ley de los Sistemas de Ahorro para el Retiro corresponde a la Secretaría de Hacienda y Crédito Público.\par Fundamento: Esto se establece en el Artículo 4o de la Ley del SAR, el cual señala expresamente que la interpretación administrativa de la Ley es facultad de la Secretaría de Hacienda y Crédito Público.\par Referencia: Ley del SAR [Página 3]
  \smallskip
  \end{minipage}}
\end{minipage}

\vspace{0.35em}
\noindent
\begin{minipage}{0.95\textwidth}
\setlength{\parskip}{0.25em}
\noindent{\footnotesize\itshape \textcolor{profublue}{\textbf{Respuesta esperada:}} A la Secretar?a de Hacienda y Cr?dito P?blico.}

\noindent{\footnotesize\itshape \textcolor{profublue}{\textbf{Método de evaluación:}} CompareMeaning}

\noindent{\footnotesize\itshape \textcolor{profublue}{\textbf{Calificación aprobatoria:}} 50}

\noindent{\footnotesize\itshape \textcolor{profublue}{\textbf{Calificación de la respuesta:}} Pass}

\noindent{\footnotesize\itshape \textcolor{profublue}{\textbf{Explicación:}} 75/100: casi el mismo significado; The Agent answer gives the correct main point that the Secretaría de Hacienda y Crédito Público is responsible for interpreting the law, just like the Expected response. It also adds useful extra details about the law article that supports this, which is good and not wrong. So, the Agent answer matches well but includes more information than needed.}

\end{minipage}

\vspace{2.5em}
\noindent
\begin{minipage}{\textwidth}
\textcolor{profublue}{\textbf{Usuario (30):}}\\[4pt]
\fcolorbox{profublue!30}{profublue!3}{%
  \begin{minipage}{0.95\textwidth}
  \smallskip
  ?Qu? obligaci?n de informaci?n peri?dica hacia los trabajadores establece el art?culo 18?
  \smallskip
  \end{minipage}}
\end{minipage}

\vspace{0.4em}
\noindent
\begin{minipage}{\textwidth}
\textcolor{profugold!80!black}{\textbf{Agente (30):}}\\[4pt]
\fcolorbox{profugold!40}{profugold!5}{%
  \begin{minipage}{0.95\textwidth}
  \smallskip
  **Obligación de información periódica hacia los trabajadores según el artículo 18 de la Ley del SAR**\par De acuerdo con el texto disponible, el artículo 18 establece que las administradoras y las prestadoras de servicio deben reportar la individualización de las aportaciones y rendimientos correspondientes a los trabajadores. Esta información debe ser calculada por el Instituto de Seguridad y Servicios Sociales de los Trabajadores del Estado, con base en la tasa de rendimiento que le proporcione el fondo de pensiones.\par **Fundamento y localización** - Artículo 18 - Ley del SAR [Página 14]
  \smallskip
  \end{minipage}}
\end{minipage}

\vspace{0.35em}
\noindent
\begin{minipage}{0.95\textwidth}
\setlength{\parskip}{0.25em}
\noindent{\footnotesize\itshape \textcolor{profublue}{\textbf{Respuesta esperada:}} Enviar estados de cuenta al menos tres veces al a?o en forma cuatrimestral y establecer servicios de informaci?n por Internet y atenci?n personalizada.}

\noindent{\footnotesize\itshape \textcolor{profublue}{\textbf{Método de evaluación:}} CompareMeaning}

\noindent{\footnotesize\itshape \textcolor{profublue}{\textbf{Calificación aprobatoria:}} 50}

\noindent{\footnotesize\itshape \textcolor{profublue}{\textbf{Calificación de la respuesta:}} Pass}

\noindent{\footnotesize\itshape \textcolor{profublue}{\textbf{Explicación:}} 50/100: algún significado compartido; The agent answer talks about how information is calculated and reported but misses the key points about sending account statements three times a year and offering internet and personal services. It also adds extra details that aren't needed. So, it only partly answers the question and leaves out important parts.}

\end{minipage}

\vspace{2.5em}
\noindent
\begin{minipage}{\textwidth}
\textcolor{profublue}{\textbf{Usuario (31):}}\\[4pt]
\fcolorbox{profublue!30}{profublue!3}{%
  \begin{minipage}{0.95\textwidth}
  \smallskip
  ?Qu? efecto tienen los derechos patrimoniales y corporativos de acciones adquiridas en contravenci?n al art?culo 21?
  \smallskip
  \end{minipage}}
\end{minipage}

\vspace{0.4em}
\noindent
\begin{minipage}{\textwidth}
\textcolor{profugold!80!black}{\textbf{Agente (31):}}\\[4pt]
\fcolorbox{profugold!40}{profugold!5}{%
  \begin{minipage}{0.95\textwidth}
  \smallskip
  **Respuesta breve y clara:**\par Cuando se adquieren acciones en contravención al artículo 21 de la Ley del SAR, los derechos patrimoniales y corporativos que corresponden a esas acciones quedan en suspenso. Esto significa que no pueden ejercerse hasta que se acredite que se ha obtenido la autorización o resolución correspondiente.\par **Fundamento y localización en la Ley:**\par - El texto relevante se encuentra en el artículo 21 Bis, que establece que las administradoras deben abstenerse de inscribir en el registro las transmisiones de acciones realizadas en contravención al artículo 21 y deben informar a la Comisión. - Específicamente, señala que los derechos patrimoniales y corporativos inherentes a las acciones adquiridas en contravención quedarán en suspenso y no podrán ser ejercidos hasta que se obtenga la autorización o resolución correspondiente.\par Referencia: Ley del SAR [Página 16]
  \smallskip
  \end{minipage}}
\end{minipage}

\vspace{0.35em}
\noindent
\begin{minipage}{0.95\textwidth}
\setlength{\parskip}{0.25em}
\noindent{\footnotesize\itshape \textcolor{profublue}{\textbf{Respuesta esperada:}} Quedan en suspenso y no pueden ejercerse hasta acreditar la autorizaci?n o resoluci?n correspondiente.}

\noindent{\footnotesize\itshape \textcolor{profublue}{\textbf{Método de evaluación:}} CompareMeaning}

\noindent{\footnotesize\itshape \textcolor{profublue}{\textbf{Calificación aprobatoria:}} 50}

\noindent{\footnotesize\itshape \textcolor{profublue}{\textbf{Calificación de la respuesta:}} Pass}

\noindent{\footnotesize\itshape \textcolor{profublue}{\textbf{Explicación:}} 75/100: casi el mismo significado; The agent answer gives the main point expected: the rights are suspended until proper authorization is shown. It also adds helpful legal details that explain why and how this happens. There are no mistakes or missing parts, just extra useful information.}

\end{minipage}

\vspace{2.5em}
\noindent
\begin{minipage}{\textwidth}
\textcolor{profublue}{\textbf{Usuario (32):}}\\[4pt]
\fcolorbox{profublue!30}{profublue!3}{%
  \begin{minipage}{0.95\textwidth}
  \smallskip
  ?Los cargos de los miembros de la Junta de Gobierno y del Comit? Consultivo y de Vigilancia son remunerados?
  \smallskip
  \end{minipage}}
\end{minipage}

\vspace{0.4em}
\noindent
\begin{minipage}{\textwidth}
\textcolor{profugold!80!black}{\textbf{Agente (32):}}\\[4pt]
\fcolorbox{profugold!40}{profugold!5}{%
  \begin{minipage}{0.95\textwidth}
  \smallskip
  Respuesta breve: Los cargos de los miembros de la Junta de Gobierno y del Comité Consultivo y de Vigilancia son honorarios, es decir, no reciben salario ni remuneración alguna por su desempeño.\par Fundamento: Esto se establece en el Artículo 17 de la Ley de los Sistemas de Ahorro para el Retiro, que señala explícitamente que dichos cargos no devengan salario o remuneración alguna.\par Referencia: Ley del SAR [Página 12] [1]\par [1]: cite:1 "Citation-1"
  \smallskip
  \end{minipage}}
\end{minipage}

\vspace{0.35em}
\noindent
\begin{minipage}{0.95\textwidth}
\setlength{\parskip}{0.25em}
\noindent{\footnotesize\itshape \textcolor{profublue}{\textbf{Respuesta esperada:}} No, son honorarios y no devengan salario o remuneraci?n alguna.}

\noindent{\footnotesize\itshape \textcolor{profublue}{\textbf{Método de evaluación:}} CompareMeaning}

\noindent{\footnotesize\itshape \textcolor{profublue}{\textbf{Calificación aprobatoria:}} 50}

\noindent{\footnotesize\itshape \textcolor{profublue}{\textbf{Calificación de la respuesta:}} Pass}

\noindent{\footnotesize\itshape \textcolor{profublue}{\textbf{Explicación:}} 75/100: casi el mismo significado; The Agent answer says the same thing as the Expected response: these members do not get paid a salary. It also adds helpful details about the law that supports this. There are no mistakes or missing parts, just extra useful information.}

\end{minipage}

\vspace{2.5em}
\noindent
\begin{minipage}{\textwidth}
\textcolor{profublue}{\textbf{Usuario (33):}}\\[4pt]
\fcolorbox{profublue!30}{profublue!3}{%
  \begin{minipage}{0.95\textwidth}
  \smallskip
  ?Qu? servicios pueden prestar las administradoras a las sociedades de inversi?n conforme al art?culo 32?
  \smallskip
  \end{minipage}}
\end{minipage}

\vspace{0.4em}
\noindent
\begin{minipage}{\textwidth}
\textcolor{profugold!80!black}{\textbf{Agente (33):}}\\[4pt]
\fcolorbox{profugold!40}{profugold!5}{%
  \begin{minipage}{0.95\textwidth}
  \smallskip
  **Respuesta breve y clara:** Conforme al artículo 32 de la Ley de los Sistemas de Ahorro para el Retiro (Ley del SAR), las administradoras pueden prestar a las sociedades de inversión los servicios de distribución y recompra de sus acciones. Además, para la guarda y administración de las acciones de las sociedades de inversión que operen, deben depositar dichos valores en una institución para el depósito de valores.\par **Fundamento y localización:** - Artículo 32: Servicios de distribución y recompra de acciones; obligación de depositar valores en institución autorizada. - Ley del SAR [Página 19].
  \smallskip
  \end{minipage}}
\end{minipage}

\vspace{0.35em}
\noindent
\begin{minipage}{0.95\textwidth}
\setlength{\parskip}{0.25em}
\noindent{\footnotesize\itshape \textcolor{profublue}{\textbf{Respuesta esperada:}} Servicios de distribuci?n y recompra de sus acciones.}

\noindent{\footnotesize\itshape \textcolor{profublue}{\textbf{Método de evaluación:}} CompareMeaning}

\noindent{\footnotesize\itshape \textcolor{profublue}{\textbf{Calificación aprobatoria:}} 50}

\noindent{\footnotesize\itshape \textcolor{profublue}{\textbf{Calificación de la respuesta:}} Pass}

\noindent{\footnotesize\itshape \textcolor{profublue}{\textbf{Explicación:}} 75/100: casi el mismo significado; The agent answer gives all the important information the expected response has about the services. It also adds more details about how shares must be kept safely, which is true and related but not asked for. So, the agent answer is mostly right and helpful but includes extra information that wasn't needed.}

\end{minipage}

\vspace{2.5em}
\noindent
\begin{minipage}{\textwidth}
\textcolor{profublue}{\textbf{Usuario (34):}}\\[4pt]
\fcolorbox{profublue!30}{profublue!3}{%
  \begin{minipage}{0.95\textwidth}
  \smallskip
  ?Cu?l es el n?mero m?nimo de administradores en el consejo de administraci?n de una AFORE?
  \smallskip
  \end{minipage}}
\end{minipage}

\vspace{0.4em}
\noindent
\begin{minipage}{\textwidth}
\textcolor{profugold!80!black}{\textbf{Agente (34):}}\\[4pt]
\fcolorbox{profugold!40}{profugold!5}{%
  \begin{minipage}{0.95\textwidth}
  \smallskip
  **Respuesta breve y clara:** El número mínimo de administradores en el consejo de administración de una AFORE es cinco. Es decir, el consejo de administración debe estar integrado por al menos cinco personas.\par **Fundamento legal y localización:** Esto se establece en el Artículo 20, fracción III, de la Ley de los Sistemas de Ahorro para el Retiro, que señala:   ``El número de sus administradores no será inferior a cinco y actuarán constituidos en consejo de administración.''\par La información se encuentra en las páginas 14 y 15 del texto.\par **Referencia:**   Ley del SAR [Páginas 14--15].
  \smallskip
  \end{minipage}}
\end{minipage}

\vspace{0.35em}
\noindent
\begin{minipage}{0.95\textwidth}
\setlength{\parskip}{0.25em}
\noindent{\footnotesize\itshape \textcolor{profublue}{\textbf{Respuesta esperada:}} No inferior a cinco.}

\noindent{\footnotesize\itshape \textcolor{profublue}{\textbf{Método de evaluación:}} CompareMeaning}

\noindent{\footnotesize\itshape \textcolor{profublue}{\textbf{Calificación aprobatoria:}} 50}

\noindent{\footnotesize\itshape \textcolor{profublue}{\textbf{Calificación de la respuesta:}} Pass}

\noindent{\footnotesize\itshape \textcolor{profublue}{\textbf{Explicación:}} 75/100: casi el mismo significado; The Agent answer gives the exact number expected (five) and adds helpful details about the law that supports this. It doesn't leave out anything important or add wrong information. So, it matches the expected answer very well.}

\end{minipage}

\vspace{2.5em}
\noindent
\begin{minipage}{\textwidth}
\textcolor{profublue}{\textbf{Usuario (35):}}\\[4pt]
\fcolorbox{profublue!30}{profublue!3}{%
  \begin{minipage}{0.95\textwidth}
  \smallskip
  ?Qu? entiende la Ley del SAR por Fondos de Previsi?n Social?
  \smallskip
  \end{minipage}}
\end{minipage}

\vspace{0.4em}
\noindent
\begin{minipage}{\textwidth}
\textcolor{profugold!80!black}{\textbf{Agente (35):}}\\[4pt]
\fcolorbox{profugold!40}{profugold!5}{%
  \begin{minipage}{0.95\textwidth}
  \smallskip
  Definición breve y clara: La Ley del SAR entiende por Fondos de Previsión Social a los fondos de pensiones o jubilaciones de personal, de primas de antigüedad, así como fondos de ahorro establecidos por empresas privadas, dependencias o entidades públicas federales, estatales o municipales, o por cualquier otra persona, como una prestación laboral a favor de los trabajadores.\par Fundamento: Esta definición se encuentra en el Artículo 3o, fracción V, de la Ley del SAR, y abarca desde la Página 1 hasta la Página 2.\par Referencia: Ley del SAR [Páginas 1--2].
  \smallskip
  \end{minipage}}
\end{minipage}

\vspace{0.35em}
\noindent
\begin{minipage}{0.95\textwidth}
\setlength{\parskip}{0.25em}
\noindent{\footnotesize\itshape \textcolor{profublue}{\textbf{Respuesta esperada:}} Fondos de pensiones, jubilaciones, primas de antig?edad o fondos de ahorro establecidos como prestaci?n laboral por empresas privadas o entidades p?blicas.}

\noindent{\footnotesize\itshape \textcolor{profublue}{\textbf{Método de evaluación:}} CompareMeaning}

\noindent{\footnotesize\itshape \textcolor{profublue}{\textbf{Calificación aprobatoria:}} 50}

\noindent{\footnotesize\itshape \textcolor{profublue}{\textbf{Calificación de la respuesta:}} Pass}

\noindent{\footnotesize\itshape \textcolor{profublue}{\textbf{Explicación:}} 75/100: casi el mismo significado; The Agent answer includes all the important points from the Expected response and adds some extra helpful details. It does not leave out anything important or give wrong information. Because it adds more useful information, it is a very good match but with some extra details.}

\end{minipage}

\vspace{2.5em}
\noindent
\begin{minipage}{\textwidth}
\textcolor{profublue}{\textbf{Usuario (36):}}\\[4pt]
\fcolorbox{profublue!30}{profublue!3}{%
  \begin{minipage}{0.95\textwidth}
  \smallskip
  ?Qu? funci?n tiene la Comisi?n respecto a la Base de Datos Nacional SAR?
  \smallskip
  \end{minipage}}
\end{minipage}

\vspace{0.4em}
\noindent
\begin{minipage}{\textwidth}
\textcolor{profugold!80!black}{\textbf{Agente (36):}}\\[4pt]
\fcolorbox{profugold!40}{profugold!5}{%
  \begin{minipage}{0.95\textwidth}
  \smallskip
  **Respuesta breve y clara:**\par La función de la Comisión Nacional del Sistema de Ahorro para el Retiro (Comisión) respecto a la Base de Datos Nacional SAR es la de coordinar, regular, supervisar y vigilar el funcionamiento de los sistemas de ahorro para el retiro, lo que incluye la gestión y supervisión de la Base de Datos Nacional SAR. Esta base contiene la información individual de cada trabajador y el registro de la administradora o institución de crédito en la que cada uno está afiliado.\par **Fundamento y localización en la Ley:**\par - El artículo 2o establece que la Comisión es responsable de la coordinación, regulación, supervisión y vigilancia de los sistemas de ahorro para el retiro. - El artículo 3o define la Base de Datos Nacional SAR y la Comisión. - El contenido relevante se encuentra en las páginas 1 y 2 del texto.\par **Referencia:**   Ley del SAR [Páginas 1--2]
  \smallskip
  \end{minipage}}
\end{minipage}

\vspace{0.35em}
\noindent
\begin{minipage}{0.95\textwidth}
\setlength{\parskip}{0.25em}
\noindent{\footnotesize\itshape \textcolor{profublue}{\textbf{Respuesta esperada:}} Administrarla y operarla, en su caso.}

\noindent{\footnotesize\itshape \textcolor{profublue}{\textbf{Método de evaluación:}} CompareMeaning}

\noindent{\footnotesize\itshape \textcolor{profublue}{\textbf{Calificación aprobatoria:}} 50}

\noindent{\footnotesize\itshape \textcolor{profublue}{\textbf{Calificación de la respuesta:}} Pass}

\noindent{\footnotesize\itshape \textcolor{profublue}{\textbf{Explicación:}} 50/100: algún significado compartido; The Expected response says the Commission's job is to run and manage the database. The Agent answer talks about many related roles and gives extra details but doesn't clearly say the Commission runs or manages the database. It leaves out this important part and adds more information than needed. So, the Agent answer partly matches but misses the main point and has extra details.}

\end{minipage}

\vspace{2.5em}
\noindent
\begin{minipage}{\textwidth}
\textcolor{profublue}{\textbf{Usuario (37):}}\\[4pt]
\fcolorbox{profublue!30}{profublue!3}{%
  \begin{minipage}{0.95\textwidth}
  \smallskip
  ?Qu? nivel de autoridad tiene el Presidente de la Comisi?n dentro de la instituci?n?
  \smallskip
  \end{minipage}}
\end{minipage}

\vspace{0.4em}
\noindent
\begin{minipage}{\textwidth}
\textcolor{profugold!80!black}{\textbf{Agente (37):}}\\[4pt]
\fcolorbox{profugold!40}{profugold!5}{%
  \begin{minipage}{0.95\textwidth}
  \smallskip
  Nivel de autoridad del Presidente de la Comisión\par Respuesta breve y clara: El Presidente de la Comisión es la máxima autoridad administrativa dentro de la institución. Tiene a su cargo la representación legal de la Comisión, dirige administrativamente la institución y ejerce las facultades que le otorga la Ley del SAR, así como aquellas que le delegue la Junta de Gobierno. Además, puede actuar directamente o a través de otros servidores públicos de la Comisión, según el reglamento interior o acuerdos delegatorios publicados oficialmente.\par Fundamento normativo: - Artículo 11: El Presidente es la máxima autoridad administrativa de la Comisión. - Artículo 12: Facultades y obligaciones del Presidente, incluyendo la representación legal, dirección administrativa y presentación de informes a la Junta de Gobierno.\par Referencia: Ley del SAR [Páginas 8--9]
  \smallskip
  \end{minipage}}
\end{minipage}

\vspace{0.35em}
\noindent
\begin{minipage}{0.95\textwidth}
\setlength{\parskip}{0.25em}
\noindent{\footnotesize\itshape \textcolor{profublue}{\textbf{Respuesta esperada:}} Es la m?xima autoridad administrativa de la Comisi?n.}

\noindent{\footnotesize\itshape \textcolor{profublue}{\textbf{Método de evaluación:}} CompareMeaning}

\noindent{\footnotesize\itshape \textcolor{profublue}{\textbf{Calificación aprobatoria:}} 50}

\noindent{\footnotesize\itshape \textcolor{profublue}{\textbf{Calificación de la respuesta:}} Pass}

\noindent{\footnotesize\itshape \textcolor{profublue}{\textbf{Explicación:}} 75/100: casi el mismo significado; The Agent answer gives all the important information the Expected response has and adds more useful details about the President's role. The extra information is correct and helps explain better, so the answer is very good. It includes everything needed and more, but the extra parts are not necessary for a simple answer.}

\end{minipage}

\vspace{2.5em}
\noindent
\begin{minipage}{\textwidth}
\textcolor{profublue}{\textbf{Usuario (38):}}\\[4pt]
\fcolorbox{profublue!30}{profublue!3}{%
  \begin{minipage}{0.95\textwidth}
  \smallskip
  ?Qu? son los Institutos de Seguridad Social para efectos de esta Ley?
  \smallskip
  \end{minipage}}
\end{minipage}

\vspace{0.4em}
\noindent
\begin{minipage}{\textwidth}
\textcolor{profugold!80!black}{\textbf{Agente (38):}}\\[4pt]
\fcolorbox{profugold!40}{profugold!5}{%
  \begin{minipage}{0.95\textwidth}
  \smallskip
  **Respuesta breve y clara:**\par Para efectos de la Ley de los Sistemas de Ahorro para el Retiro, los **Institutos de Seguridad Social** son el Instituto Mexicano del Seguro Social (IMSS), el Instituto del Fondo Nacional de la Vivienda para los Trabajadores (INFONAVIT), el Instituto de Seguridad y Servicios Sociales de los Trabajadores del Estado (ISSSTE) y otras instituciones de naturaleza análoga.\par **Fundamento y localización en la Ley:**\par - Capítulo I, Artículo 3o, fracción VI. - Ley del SAR [Páginas 1--2].
  \smallskip
  \end{minipage}}
\end{minipage}

\vspace{0.35em}
\noindent
\begin{minipage}{0.95\textwidth}
\setlength{\parskip}{0.25em}
\noindent{\footnotesize\itshape \textcolor{profublue}{\textbf{Respuesta esperada:}} El IMSS, INFONAVIT, ISSSTE y otras instituciones de naturaleza an?loga.}

\noindent{\footnotesize\itshape \textcolor{profublue}{\textbf{Método de evaluación:}} CompareMeaning}

\noindent{\footnotesize\itshape \textcolor{profublue}{\textbf{Calificación aprobatoria:}} 50}

\noindent{\footnotesize\itshape \textcolor{profublue}{\textbf{Calificación de la respuesta:}} Pass}

\noindent{\footnotesize\itshape \textcolor{profublue}{\textbf{Explicación:}} 75/100: casi el mismo significado; The Agent answer gives all the important parts the Expected response has and adds useful details like the law's name and where to find it in the law. It doesn't leave out anything or say anything wrong. Because of this extra helpful info, the Agent answer is a very good match but includes more than what was asked.}

\end{minipage}

\vspace{2.5em}
\noindent
\begin{minipage}{\textwidth}
\textcolor{profublue}{\textbf{Usuario (39):}}\\[4pt]
\fcolorbox{profublue!30}{profublue!3}{%
  \begin{minipage}{0.95\textwidth}
  \smallskip
  ?Qu? documentos deben presentar los solicitantes para obtener autorizaci?n como administradora?
  \smallskip
  \end{minipage}}
\end{minipage}

\vspace{0.4em}
\noindent
\begin{minipage}{\textwidth}
\textcolor{profugold!80!black}{\textbf{Agente (39):}}\\[4pt]
\fcolorbox{profugold!40}{profugold!5}{%
  \begin{minipage}{0.95\textwidth}
  \smallskip
  **Requisitos documentales para obtener autorización como administradora**\par Para que los solicitantes obtengan autorización para organizarse y operar como administradora de fondos para el retiro, deben presentar ante la Comisión los siguientes documentos:\par - Solicitud respectiva y proyecto de estatutos sociales. - Programa general de operación y funcionamiento, de divulgación de la información y de reinversión de utilidades, cumpliendo los requisitos mínimos que determine la Comisión. - Estado de situación patrimonial de los accionistas que detenten el control, abarcando los cinco años anteriores, conforme a lo que señale la Comisión. - Escrituras constitutivas de la sociedad y sus reformas, que deben ser aprobadas por la Comisión y posteriormente inscritas en el Registro Público de Comercio. - Copia certificada de las actas de asamblea y, cuando proceda, testimonio notarial de la protocolización de las mismas.\par **Fundamento legal:**   Estos requisitos se encuentran en el Artículo 19 de la Ley del SAR, que abarca las páginas 14 y 15 del texto consultado.   Ley del SAR [Páginas 14--15].
  \smallskip
  \end{minipage}}
\end{minipage}

\vspace{0.35em}
\noindent
\begin{minipage}{0.95\textwidth}
\setlength{\parskip}{0.25em}
\noindent{\footnotesize\itshape \textcolor{profublue}{\textbf{Respuesta esperada:}} La solicitud respectiva, el proyecto de estatutos sociales y un programa general de operaci?n y funcionamiento, divulgaci?n de informaci?n y reinversi?n de utilidades.}

\noindent{\footnotesize\itshape \textcolor{profublue}{\textbf{Método de evaluación:}} CompareMeaning}

\noindent{\footnotesize\itshape \textcolor{profublue}{\textbf{Calificación aprobatoria:}} 50}

\noindent{\footnotesize\itshape \textcolor{profublue}{\textbf{Calificación de la respuesta:}} Pass}

\noindent{\footnotesize\itshape \textcolor{profublue}{\textbf{Explicación:}} 75/100: casi el mismo significado; The Agent answer includes all the important documents the Expected response mentions and adds more details about other documents needed. It doesn't have wrong information but gives more than what was asked. So, it mostly matches but has extra details.}

\end{minipage}

\vspace{2.5em}
\noindent
\begin{minipage}{\textwidth}
\textcolor{profublue}{\textbf{Usuario (40):}}\\[4pt]
\fcolorbox{profublue!30}{profublue!3}{%
  \begin{minipage}{0.95\textwidth}
  \smallskip
  ?Qui?n es considerado Trabajador para efectos de esta Ley?
  \smallskip
  \end{minipage}}
\end{minipage}

\vspace{0.4em}
\noindent
\begin{minipage}{\textwidth}
\textcolor{profugold!80!black}{\textbf{Agente (40):}}\\[4pt]
\fcolorbox{profugold!40}{profugold!5}{%
  \begin{minipage}{0.95\textwidth}
  \smallskip
  **Respuesta breve y clara:**\par Para efectos de la Ley de los Sistemas de Ahorro para el Retiro (Ley del SAR), se considera **Trabajador** tanto a los trabajadores afiliados como a cualquier otra persona que tenga derecho a la apertura de una cuenta individual conforme a lo dispuesto por la Ley. Es decir, no sólo quienes están inscritos en el Instituto Mexicano del Seguro Social, sino también quienes, por cualquier motivo, tienen derecho a una cuenta individual bajo esta Ley.\par **Fundamento y localización en la Ley:**\par - La definición de "Trabajador" se encuentra en la fracción XII del artículo correspondiente, donde se establece:   - **Trabajador:** ``a los trabajadores afiliados, así como a cualquier otra persona que tenga derecho a la apertura de una cuenta individual en los términos de esta ley.'' - Esta información está ubicada en las páginas 2 y 3 del texto.\par **Referencia:** Ley del SAR [Páginas 2--3]
  \smallskip
  \end{minipage}}
\end{minipage}

\vspace{0.35em}
\noindent
\begin{minipage}{0.95\textwidth}
\setlength{\parskip}{0.25em}
\noindent{\footnotesize\itshape \textcolor{profublue}{\textbf{Respuesta esperada:}} Trabajadores afiliados y cualquier otra persona con derecho a abrir una cuenta individual en t?rminos de esta Ley.}

\noindent{\footnotesize\itshape \textcolor{profublue}{\textbf{Método de evaluación:}} CompareMeaning}

\noindent{\footnotesize\itshape \textcolor{profublue}{\textbf{Calificación aprobatoria:}} 50}

\noindent{\footnotesize\itshape \textcolor{profublue}{\textbf{Calificación de la respuesta:}} Pass}

\noindent{\footnotesize\itshape \textcolor{profublue}{\textbf{Explicación:}} 75/100: casi el mismo significado; The Agent answer correctly explains who is considered a worker according to the law, matching the key points in the Expected response. It also adds some extra helpful details about the law's article and who is included, which is good but more than what was asked. So, the answer is very good and mostly matches, with some extra information.}

\end{minipage}

\vspace{2.5em}
\noindent
\begin{minipage}{\textwidth}
\textcolor{profublue}{\textbf{Usuario (41):}}\\[4pt]
\fcolorbox{profublue!30}{profublue!3}{%
  \begin{minipage}{0.95\textwidth}
  \smallskip
  ?Qu? nombramientos deben informar las administradoras a la Comisi?n?
  \smallskip
  \end{minipage}}
\end{minipage}

\vspace{0.4em}
\noindent
\begin{minipage}{\textwidth}
\textcolor{profugold!80!black}{\textbf{Agente (41):}}\\[4pt]
\fcolorbox{profugold!40}{profugold!5}{%
  \begin{minipage}{0.95\textwidth}
  \smallskip
  **Respuesta breve y clara:** Las administradoras deben informar a la Comisión los nombramientos de los miembros de su consejo de administración, del director general, de los funcionarios de los dos niveles inmediatos siguientes y de sus comisarios.\par **Fundamento y localización en la Ley del SAR:** Esto se encuentra en el Artículo 20, fracción IV, que establece la obligación de informar estos nombramientos a la Comisión.\par **Referencia normativa:** Ley del SAR [Páginas 14--15].
  \smallskip
  \end{minipage}}
\end{minipage}

\vspace{0.35em}
\noindent
\begin{minipage}{0.95\textwidth}
\setlength{\parskip}{0.25em}
\noindent{\footnotesize\itshape \textcolor{profublue}{\textbf{Respuesta esperada:}} Miembros del consejo de administraci?n, director general, funcionarios de los dos niveles inmediatos siguientes y comisarios; y someter a aprobaci?n del Comit? Consultivo los nombramientos de consejeros independientes y contralor normativo.}

\noindent{\footnotesize\itshape \textcolor{profublue}{\textbf{Método de evaluación:}} CompareMeaning}

\noindent{\footnotesize\itshape \textcolor{profublue}{\textbf{Calificación aprobatoria:}} 50}

\noindent{\footnotesize\itshape \textcolor{profublue}{\textbf{Calificación de la respuesta:}} Pass}

\noindent{\footnotesize\itshape \textcolor{profublue}{\textbf{Explicación:}} 75/100: casi el mismo significado; The agent answer lists most of the appointments that must be reported, just like the expected response. But it leaves out that some appointments need approval from a special committee. It also adds legal references, which is helpful. Because it misses an important step, it is mostly correct but not complete.}

\end{minipage}

\vspace{2.5em}
\noindent
\begin{minipage}{\textwidth}
\textcolor{profublue}{\textbf{Usuario (42):}}\\[4pt]
\fcolorbox{profublue!30}{profublue!3}{%
  \begin{minipage}{0.95\textwidth}
  \smallskip
  ?C?mo define la Ley del SAR a los Participantes en los sistemas de ahorro para el retiro?
  \smallskip
  \end{minipage}}
\end{minipage}

\vspace{0.4em}
\noindent
\begin{minipage}{\textwidth}
\textcolor{profugold!80!black}{\textbf{Agente (42):}}\\[4pt]
\fcolorbox{profugold!40}{profugold!5}{%
  \begin{minipage}{0.95\textwidth}
  \smallskip
  Definición breve y clara: La Ley del SAR define a los participantes como las personas y entidades que intervienen en el funcionamiento de los sistemas de ahorro para el retiro, conforme a lo previsto en la propia Ley y en las leyes del Seguro Social, del Instituto del Fondo Nacional de la Vivienda para los Trabajadores y del Instituto de Seguridad y Servicios Sociales de los Trabajadores del Estado.\par Fundamento y localización: Esta definición se encuentra en el Artículo 1o., Capítulo I, Disposiciones Preliminares, y se complementa con el Artículo 3o., que enumera a los principales participantes, como las Administradoras de Fondos para el Retiro, la Comisión Nacional del Sistema de Ahorro para el Retiro, los Institutos de Seguridad Social, entre otros.\par Referencia: Ley del SAR [Páginas 1--2]
  \smallskip
  \end{minipage}}
\end{minipage}

\vspace{0.35em}
\noindent
\begin{minipage}{0.95\textwidth}
\setlength{\parskip}{0.25em}
\noindent{\footnotesize\itshape \textcolor{profublue}{\textbf{Respuesta esperada:}} Instituciones de cr?dito, AFORE, sociedades de inversi?n, empresas operadoras, prestadores de servicios complementarios y entidades receptoras previstas en el reglamento.}

\noindent{\footnotesize\itshape \textcolor{profublue}{\textbf{Método de evaluación:}} CompareMeaning}

\noindent{\footnotesize\itshape \textcolor{profublue}{\textbf{Calificación aprobatoria:}} 50}

\noindent{\footnotesize\itshape \textcolor{profublue}{\textbf{Calificación de la respuesta:}} Pass}

\noindent{\footnotesize\itshape \textcolor{profublue}{\textbf{Explicación:}} 50/100: algún significado compartido; The Agent answer explains who the participants are in general terms and mentions some related laws and institutions. But it does not clearly list the specific groups of participants that the Expected response includes. It misses important details about the exact types of participants. So, the Agent answer is partly correct but leaves out key information.}

\end{minipage}

\vspace{2.5em}
\noindent
\begin{minipage}{\textwidth}
\textcolor{profublue}{\textbf{Usuario (43):}}\\[4pt]
\fcolorbox{profublue!30}{profublue!3}{%
  \begin{minipage}{0.95\textwidth}
  \smallskip
  ?Qu? potestad tiene la Comisi?n respecto a la constituci?n y funcionamiento de los participantes del sistema?
  \smallskip
  \end{minipage}}
\end{minipage}

\vspace{0.4em}
\noindent
\begin{minipage}{\textwidth}
\textcolor{profugold!80!black}{\textbf{Agente (43):}}\\[4pt]
\fcolorbox{profugold!40}{profugold!5}{%
  \begin{minipage}{0.95\textwidth}
  \smallskip
  **Potestad de la Comisión sobre la constitución y funcionamiento de los participantes del sistema**\par La Comisión Nacional del Sistema de Ahorro para el Retiro (CONSAR) tiene la facultad de expedir disposiciones de carácter general que regulan la constitución, organización, funcionamiento, operaciones y participación de los participantes en los sistemas de ahorro para el retiro. Esto incluye a instituciones de crédito, administradoras de fondos para el retiro, sociedades de inversión especializadas, empresas operadoras y otras entidades relacionadas.\par Además, la Comisión puede regular la operación de los sistemas, emitir regulación prudencial, establecer reglas para retiros programados, y otorgar, modificar o revocar autorizaciones a administradoras y sociedades de inversión.\par **Fundamento legal:**   - Artículo 5o., fracciones I, II, III, IV, V y VI   - Ley del SAR [Páginas 2--3]
  \smallskip
  \end{minipage}}
\end{minipage}

\vspace{0.35em}
\noindent
\begin{minipage}{0.95\textwidth}
\setlength{\parskip}{0.25em}
\noindent{\footnotesize\itshape \textcolor{profublue}{\textbf{Respuesta esperada:}} Expedir disposiciones generales a las que deben sujetarse los participantes en cuanto a su constituci?n, organizaci?n, funcionamiento, operaciones y participaci?n.}

\noindent{\footnotesize\itshape \textcolor{profublue}{\textbf{Método de evaluación:}} CompareMeaning}

\noindent{\footnotesize\itshape \textcolor{profublue}{\textbf{Calificación aprobatoria:}} 50}

\noindent{\footnotesize\itshape \textcolor{profublue}{\textbf{Calificación de la respuesta:}} Pass}

\noindent{\footnotesize\itshape \textcolor{profublue}{\textbf{Explicación:}} 75/100: casi el mismo significado; The Agent answer includes all the important points from the Expected response and adds useful extra details about the Comisión's powers. It does not leave out anything important or say anything wrong. The extra information helps explain the topic better, so the answers match well.}

\end{minipage}

\vspace{2.5em}
\noindent
\begin{minipage}{\textwidth}
\textcolor{profublue}{\textbf{Usuario (44):}}\\[4pt]
\fcolorbox{profublue!30}{profublue!3}{%
  \begin{minipage}{0.95\textwidth}
  \smallskip
  ?Qu? mantiene la Comisi?n respecto a los agentes promotores de las administradoras?
  \smallskip
  \end{minipage}}
\end{minipage}

\vspace{0.4em}
\noindent
\begin{minipage}{\textwidth}
\textcolor{profugold!80!black}{\textbf{Agente (44):}}\\[4pt]
\fcolorbox{profugold!40}{profugold!5}{%
  \begin{minipage}{0.95\textwidth}
  \smallskip
  **Respuesta breve y clara:**\par La Comisión puede imponer multas a los agentes promotores de las administradoras que registren a un trabajador o soliciten el traspaso de su cuenta individual sin su consentimiento, o si dicho consentimiento se obtuvo mediante engaño, mala fe, documentos falsos, firmas falsificadas o la entrega de beneficios indebidos. Además, en caso de reincidencia y dependiendo de la gravedad, la Comisión puede suspender el registro del agente promotor por seis meses hasta un año, o incluso inhabilitarlo como agente promotor. Estas sanciones son independientes de otras responsabilidades administrativas, civiles o penales.\par **Fundamento y localización en la Ley del SAR:**\par - Artículo 100 B - Ley del SAR [Página 59]
  \smallskip
  \end{minipage}}
\end{minipage}

\vspace{0.35em}
\noindent
\begin{minipage}{0.95\textwidth}
\setlength{\parskip}{0.25em}
\noindent{\footnotesize\itshape \textcolor{profublue}{\textbf{Respuesta esperada:}} Un registro y puede suspenderlo o cancelarlo conforme a la Ley.}

\noindent{\footnotesize\itshape \textcolor{profublue}{\textbf{Método de evaluación:}} CompareMeaning}

\noindent{\footnotesize\itshape \textcolor{profublue}{\textbf{Calificación aprobatoria:}} 50}

\noindent{\footnotesize\itshape \textcolor{profublue}{\textbf{Calificación de la respuesta:}} Pass}

\noindent{\footnotesize\itshape \textcolor{profublue}{\textbf{Explicación:}} 75/100: casi el mismo significado; The expected response says the Commission can register, suspend, or cancel agents. The agent answer says this too but adds more details about fines and reasons for punishment. The extra details are correct and helpful but go beyond what was asked. So, the agent answer matches well but has extra information.}

\end{minipage}

\vspace{2.5em}
\noindent
\begin{minipage}{\textwidth}
\textcolor{profublue}{\textbf{Usuario (45):}}\\[4pt]
\fcolorbox{profublue!30}{profublue!3}{%
  \begin{minipage}{0.95\textwidth}
  \smallskip
  ?El contralor normativo debe asistir a sesiones de ?rganos internos?
  \smallskip
  \end{minipage}}
\end{minipage}

\vspace{0.4em}
\noindent
\begin{minipage}{\textwidth}
\textcolor{profugold!80!black}{\textbf{Agente (45):}}\\[4pt]
\fcolorbox{profugold!40}{profugold!5}{%
  \begin{minipage}{0.95\textwidth}
  \smallskip
  Respuesta breve: Sí, el contralor normativo debe asistir a las sesiones de los órganos internos, específicamente a las sesiones de consejo de administración de las administradoras, de las sociedades de inversión y del comité de inversión. Su participación es con voz pero sin voto.\par Fundamento: Esto se establece en el texto de la Ley del SAR, donde se indica que el contralor normativo debe asistir a dichas sesiones y participar en ellas, aunque no tiene derecho a voto. Esta obligación se encuentra en la página 19 del documento.\par Referencia: Ley del SAR [Página 19]
  \smallskip
  \end{minipage}}
\end{minipage}

\vspace{0.35em}
\noindent
\begin{minipage}{0.95\textwidth}
\setlength{\parskip}{0.25em}
\noindent{\footnotesize\itshape \textcolor{profublue}{\textbf{Respuesta esperada:}} S?, debe asistir a sesiones del consejo de administraci?n, de sociedades de inversi?n y del comit? de inversi?n, con voz pero sin voto.}

\noindent{\footnotesize\itshape \textcolor{profublue}{\textbf{Método de evaluación:}} CompareMeaning}

\noindent{\footnotesize\itshape \textcolor{profublue}{\textbf{Calificación aprobatoria:}} 50}

\noindent{\footnotesize\itshape \textcolor{profublue}{\textbf{Calificación de la respuesta:}} Pass}

\noindent{\footnotesize\itshape \textcolor{profublue}{\textbf{Explicación:}} 75/100: casi el mismo significado; The Agent answer gives all the important information the Expected response has. It says the contralor normativo must attend certain meetings and can speak but not vote. It also adds where this rule comes from, which is helpful and correct. There are no mistakes or missing parts, so the answers match very well.}

\end{minipage}

\vspace{2.5em}
\noindent
\begin{minipage}{\textwidth}
\textcolor{profublue}{\textbf{Usuario (46):}}\\[4pt]
\fcolorbox{profublue!30}{profublue!3}{%
  \begin{minipage}{0.95\textwidth}
  \smallskip
  ?D?nde deben depositar las administradoras los t?tulos de las sociedades de inversi?n que operen?
  \smallskip
  \end{minipage}}
\end{minipage}

\vspace{0.4em}
\noindent
\begin{minipage}{\textwidth}
\textcolor{profugold!80!black}{\textbf{Agente (46):}}\\[4pt]
\fcolorbox{profugold!40}{profugold!5}{%
  \begin{minipage}{0.95\textwidth}
  \smallskip
  Depósito de títulos de sociedades de inversión\par Las administradoras, para la guarda y administración de las acciones de las sociedades de inversión que operen, deben depositar dichos valores en una institución para el depósito de valores.\par Fundamento: - Artículo 32 - Ley del SAR [Páginas 19--20]
  \smallskip
  \end{minipage}}
\end{minipage}

\vspace{0.35em}
\noindent
\begin{minipage}{0.95\textwidth}
\setlength{\parskip}{0.25em}
\noindent{\footnotesize\itshape \textcolor{profublue}{\textbf{Respuesta esperada:}} En una instituci?n para el dep?sito de valores para la guarda y administraci?n de las acciones.}

\noindent{\footnotesize\itshape \textcolor{profublue}{\textbf{Método de evaluación:}} CompareMeaning}

\noindent{\footnotesize\itshape \textcolor{profublue}{\textbf{Calificación aprobatoria:}} 50}

\noindent{\footnotesize\itshape \textcolor{profublue}{\textbf{Calificación de la respuesta:}} Pass}

\noindent{\footnotesize\itshape \textcolor{profublue}{\textbf{Explicación:}} 75/100: casi el mismo significado; The Agent answer gives the same main information as the Expected response: the titles must be deposited in a securities deposit institution. It also adds helpful details like why and legal references, which do not change or leave out anything important. So, the Agent answer matches well and is complete.}

\end{minipage}

\vspace{2.5em}
\noindent
\begin{minipage}{\textwidth}
\textcolor{profublue}{\textbf{Usuario (47):}}\\[4pt]
\fcolorbox{profublue!30}{profublue!3}{%
  \begin{minipage}{0.95\textwidth}
  \smallskip
  ?Qu? sucede con transmisiones de acciones realizadas en contravenci?n al art?culo 21?
  \smallskip
  \end{minipage}}
\end{minipage}

\vspace{0.4em}
\noindent
\begin{minipage}{\textwidth}
\textcolor{profugold!80!black}{\textbf{Agente (47):}}\\[4pt]
\fcolorbox{profugold!40}{profugold!5}{%
  \begin{minipage}{0.95\textwidth}
  \smallskip
  **Respuesta breve y clara:**\par Cuando se realizan transmisiones de acciones en contravención al artículo 21 de la Ley de los Sistemas de Ahorro para el Retiro, los derechos patrimoniales y corporativos inherentes a esas acciones quedan en suspenso. Esto significa que no pueden ejercerse hasta que se acredite que se ha obtenido la autorización o resolución correspondiente. Además, las administradoras deben abstenerse de inscribir dichas transmisiones en el registro y deben informar a la Comisión dentro de los cinco días hábiles siguientes a que tengan conocimiento de ello.\par **Fundamento legal y localización en el texto:**\par - Artículo 21 Bis: Suspensión de derechos y obligación de informar a la Comisión. - Página 16 del documento.\par Referencia: Ley del SAR [Página 16].
  \smallskip
  \end{minipage}}
\end{minipage}

\vspace{0.35em}
\noindent
\begin{minipage}{0.95\textwidth}
\setlength{\parskip}{0.25em}
\noindent{\footnotesize\itshape \textcolor{profublue}{\textbf{Respuesta esperada:}} La administradora debe abstenerse de inscribirlas en el registro correspondiente e informar a la Comisi?n en cinco d?as h?biles.}

\noindent{\footnotesize\itshape \textcolor{profublue}{\textbf{Método de evaluación:}} CompareMeaning}

\noindent{\footnotesize\itshape \textcolor{profublue}{\textbf{Calificación aprobatoria:}} 50}

\noindent{\footnotesize\itshape \textcolor{profublue}{\textbf{Calificación de la respuesta:}} Pass}

\noindent{\footnotesize\itshape \textcolor{profublue}{\textbf{Explicación:}} 75/100: casi el mismo significado; The Agent answer gives all the important points from the Expected response and adds helpful extra details about what happens when shares are transferred against the rules. It does not leave out anything important or say anything wrong. The extra information helps explain the situation better.}

\end{minipage}

\vspace{2.5em}
\noindent
\begin{minipage}{\textwidth}
\textcolor{profublue}{\textbf{Usuario (48):}}\\[4pt]
\fcolorbox{profublue!30}{profublue!3}{%
  \begin{minipage}{0.95\textwidth}
  \smallskip
  ?Qu? requisito permanente de capital fijo deben cumplir las administradoras?
  \smallskip
  \end{minipage}}
\end{minipage}

\vspace{0.4em}
\noindent
\begin{minipage}{\textwidth}
\textcolor{profugold!80!black}{\textbf{Agente (48):}}\\[4pt]
\fcolorbox{profugold!40}{profugold!5}{%
  \begin{minipage}{0.95\textwidth}
  \smallskip
  **Respuesta breve y clara:**\par Las administradoras deben contar de manera permanente con un capital fijo sin derecho a retiro, totalmente pagado, que sea al menos igual al capital mínimo exigido por la Comisión, según las disposiciones de carácter general que ésta emita. Si el capital se reduce por debajo del mínimo exigido, la administradora está obligada a reconstituirlo en el plazo que determine la Comisión, el cual no puede exceder de 45 días naturales.\par **Fundamento normativo:**\par - Artículo 24: Establece el requisito permanente de capital fijo sin derecho a retiro, totalmente pagado, igual al mínimo exigido por la Comisión. - Artículo 27: Regula el destino de las inversiones con cargo al capital mínimo pagado exigido.\par **Referencia:** Ley del SAR [Páginas 17--18]
  \smallskip
  \end{minipage}}
\end{minipage}

\vspace{0.35em}
\noindent
\begin{minipage}{0.95\textwidth}
\setlength{\parskip}{0.25em}
\noindent{\footnotesize\itshape \textcolor{profublue}{\textbf{Respuesta esperada:}} Contar con capital fijo sin derecho a retiro totalmente pagado, al menos igual al capital m?nimo exigido por la Comisi?n.}

\noindent{\footnotesize\itshape \textcolor{profublue}{\textbf{Método de evaluación:}} CompareMeaning}

\noindent{\footnotesize\itshape \textcolor{profublue}{\textbf{Calificación aprobatoria:}} 50}

\noindent{\footnotesize\itshape \textcolor{profublue}{\textbf{Calificación de la respuesta:}} Pass}

\noindent{\footnotesize\itshape \textcolor{profublue}{\textbf{Explicación:}} 75/100: casi el mismo significado; The Agent answer gives all the important parts the Expected response has about the fixed capital requirement. It also adds useful extra details about what happens if the capital drops too low and mentions the laws that support this. The extra information is helpful and correct, so the Agent answer matches very well with the Expected response.}

\end{minipage}

\vspace{2.5em}
\noindent
\begin{minipage}{\textwidth}
\textcolor{profublue}{\textbf{Usuario (49):}}\\[4pt]
\fcolorbox{profublue!30}{profublue!3}{%
  \begin{minipage}{0.95\textwidth}
  \smallskip
  ?C?mo se determina el m?ximo de comisiones de las AFORE seg?n el art?culo 37 reformado en 2020?
  \smallskip
  \end{minipage}}
\end{minipage}

\vspace{0.4em}
\noindent
\begin{minipage}{\textwidth}
\textcolor{profugold!80!black}{\textbf{Agente (49):}}\\[4pt]
\fcolorbox{profugold!40}{profugold!5}{%
  \begin{minipage}{0.95\textwidth}
  \smallskip
  **Determinación del máximo de comisiones de las AFORE según el artículo 37 reformado en 2020**\par **Respuesta breve y clara:** El artículo 37 establece que las AFORE solo pueden cobrar comisiones a los trabajadores con cuenta individual conforme a las reglas generales que expida la Comisión. Para promover un mayor rendimiento neto a favor de los trabajadores, las comisiones por administración de cuentas individuales únicamente pueden cobrarse como un porcentaje sobre el valor de los activos administrados. Además, las AFORE solo pueden cobrar cuotas fijas por servicios específicos señalados en el reglamento de la Ley, y nunca por la administración de las cuentas.\par **Fundamento y localización:** - Artículo 37, párrafo reformado. - Ley del SAR [Páginas 20].
  \smallskip
  \end{minipage}}
\end{minipage}

\vspace{0.35em}
\noindent
\begin{minipage}{0.95\textwidth}
\setlength{\parskip}{0.25em}
\noindent{\footnotesize\itshape \textcolor{profublue}{\textbf{Respuesta esperada:}} Por el promedio aritm?tico de los cobros en sistemas de contribuci?n definida de Estados Unidos de Am?rica, Colombia y Chile, conforme a pol?ticas de la Junta de Gobierno.}

\noindent{\footnotesize\itshape \textcolor{profublue}{\textbf{Método de evaluación:}} CompareMeaning}

\noindent{\footnotesize\itshape \textcolor{profublue}{\textbf{Calificación aprobatoria:}} 50}

\noindent{\footnotesize\itshape \textcolor{profublue}{\textbf{Calificación de la respuesta:}} Fail}

\noindent{\footnotesize\itshape \textcolor{profublue}{\textbf{Explicación:}} 25/100: significado principalmente diferente; The Agent answer talks about the rules for AFORE commissions but leaves out the main way the maximum commission is set, which is by averaging charges from other countries. It doesn't say how the maximum is actually determined, so it misses important information.}

\end{minipage}

\vspace{2.5em}
\noindent
\begin{minipage}{\textwidth}
\textcolor{profublue}{\textbf{Usuario (50):}}\\[4pt]
\fcolorbox{profublue!30}{profublue!3}{%
  \begin{minipage}{0.95\textwidth}
  \smallskip
  ?Requieren autorizaci?n para invertir en empresas que presten servicios complementarios o auxiliares?
  \smallskip
  \end{minipage}}
\end{minipage}

\vspace{0.4em}
\noindent
\begin{minipage}{\textwidth}
\textcolor{profugold!80!black}{\textbf{Agente (50):}}\\[4pt]
\fcolorbox{profugold!40}{profugold!5}{%
  \begin{minipage}{0.95\textwidth}
  \smallskip
  Respuesta breve: Sí, las administradoras requieren autorización de la Comisión para invertir en empresas que les presten servicios complementarios o auxiliares en la realización de su objeto.\par Fundamento: Esto se establece en el Artículo 34, el cual señala que las administradoras necesitan autorización para invertir en dichas empresas. Además, estas empresas estarán sujetas a la regulación y supervisión de la Comisión, y la administradora será responsable solidaria de las sanciones que correspondan a estas empresas por motivo de su supervisión.\par Referencia: Ley del SAR [Páginas 19--20].
  \smallskip
  \end{minipage}}
\end{minipage}

\vspace{0.35em}
\noindent
\begin{minipage}{0.95\textwidth}
\setlength{\parskip}{0.25em}
\noindent{\footnotesize\itshape \textcolor{profublue}{\textbf{Respuesta esperada:}} S?, las administradoras requieren autorizaci?n de la Comisi?n.}

\noindent{\footnotesize\itshape \textcolor{profublue}{\textbf{Método de evaluación:}} CompareMeaning}

\noindent{\footnotesize\itshape \textcolor{profublue}{\textbf{Calificación aprobatoria:}} 50}

\noindent{\footnotesize\itshape \textcolor{profublue}{\textbf{Calificación de la respuesta:}} Pass}

\noindent{\footnotesize\itshape \textcolor{profublue}{\textbf{Explicación:}} 75/100: casi el mismo significado; The agent answer correctly says that authorization is needed, just like the expected response. It also adds more details about the law and rules, which are true but not required. So, the agent answer is mostly right but gives extra information that wasn't asked for.}

\end{minipage}

\vspace{2.5em}
\noindent
\begin{minipage}{\textwidth}
\textcolor{profublue}{\textbf{Usuario (51):}}\\[4pt]
\fcolorbox{profublue!30}{profublue!3}{%
  \begin{minipage}{0.95\textwidth}
  \smallskip
  ?Cu?ntos miembros integran el Comit? Consultivo y de Vigilancia y c?mo se distribuyen?
  \smallskip
  \end{minipage}}
\end{minipage}

\vspace{0.4em}
\noindent
\begin{minipage}{\textwidth}
\textcolor{profugold!80!black}{\textbf{Agente (51):}}\\[4pt]
\fcolorbox{profugold!40}{profugold!5}{%
  \begin{minipage}{0.95\textwidth}
  \smallskip
  **Respuesta breve y clara:**\par El Comité Consultivo y de Vigilancia está integrado por diecinueve miembros. La distribución es la siguiente: - Seis representantes de los trabajadores. - Seis representantes de los patrones. - El Presidente de la Comisión. - Un representante de cada una de las siguientes dependencias y entidades:   - Secretaría de Hacienda y Crédito Público.   - Secretaría del Trabajo y Previsión Social.   - Instituto Mexicano del Seguro Social.   - Instituto de Seguridad y Servicios Sociales de los Trabajadores del Estado.   - Instituto del Fondo Nacional de la Vivienda para los Trabajadores.   - Banco de México.\par **Fundamento y localización en la Ley del SAR:**\par Esta integración y distribución se establece en el Artículo 15 de la Ley de los Sistemas de Ahorro para el Retiro, que se encuentra alrededor de las páginas 10 y 11 del texto.\par **Referencia:**   Ley del SAR [Páginas 10--11]
  \smallskip
  \end{minipage}}
\end{minipage}

\vspace{0.35em}
\noindent
\begin{minipage}{0.95\textwidth}
\setlength{\parskip}{0.25em}
\noindent{\footnotesize\itshape \textcolor{profublue}{\textbf{Respuesta esperada:}} Diecinueve: seis representantes de trabajadores, seis de patrones, el Presidente de la Comisi?n y representantes de diversas dependencias e instituciones como SHCP, STPS, IMSS, ISSSTE, INFONAVIT y Banxico.}

\noindent{\footnotesize\itshape \textcolor{profublue}{\textbf{Método de evaluación:}} CompareMeaning}

\noindent{\footnotesize\itshape \textcolor{profublue}{\textbf{Calificación aprobatoria:}} 50}

\noindent{\footnotesize\itshape \textcolor{profublue}{\textbf{Calificación de la respuesta:}} Pass}

\noindent{\footnotesize\itshape \textcolor{profublue}{\textbf{Explicación:}} 75/100: casi el mismo significado; The Agent answer gives all the important details the Expected response has about the committee members and how they are divided. It also adds where this information is found in the law, which is useful extra information. There are no mistakes or missing parts, so the answer is very good and complete.}

\end{minipage}

\vspace{2.5em}
\noindent
\begin{minipage}{\textwidth}
\textcolor{profublue}{\textbf{Usuario (52):}}\\[4pt]
\fcolorbox{profublue!30}{profublue!3}{%
  \begin{minipage}{0.95\textwidth}
  \smallskip
  ?C?mo define la Ley el concepto de 'control' para efectos de participaci?n en una administradora?
  \smallskip
  \end{minipage}}
\end{minipage}

\vspace{0.4em}
\noindent
\begin{minipage}{\textwidth}
\textcolor{profugold!80!black}{\textbf{Agente (52):}}\\[4pt]
\fcolorbox{profugold!40}{profugold!5}{%
  \begin{minipage}{0.95\textwidth}
  \smallskip
  **Definición de 'control' en la Ley del SAR para participación en una administradora**\par La Ley del SAR define el concepto de 'control' para efectos de participación en una administradora como la capacidad de imponer, directa o indirectamente, decisiones en las asambleas generales de accionistas de la administradora; el mantener la titularidad de derechos que permitan, directa o indirectamente, ejercer el voto respecto de más del cincuenta por ciento del capital social de la administradora; o dirigir, directa o indirectamente, la administración, la estrategia o las principales políticas de la administradora, ya sea a través de la propiedad de valores o por cualquier otro acto jurídico.\par **Fundamento legal y localización** Esta definición se encuentra en el texto de la Ley del SAR, específicamente en el artículo relacionado con la participación y control en las administradoras, alrededor de la Página 16.\par Referencia: Ley del SAR [Página 16].
  \smallskip
  \end{minipage}}
\end{minipage}

\vspace{0.35em}
\noindent
\begin{minipage}{0.95\textwidth}
\setlength{\parskip}{0.25em}
\noindent{\footnotesize\itshape \textcolor{profublue}{\textbf{Respuesta esperada:}} Capacidad de imponer decisiones en asambleas o mantener derechos de voto sobre m?s del 50\% del capital social, o dirigir administraci?n, estrategia o pol?ticas principales.}

\noindent{\footnotesize\itshape \textcolor{profublue}{\textbf{Método de evaluación:}} CompareMeaning}

\noindent{\footnotesize\itshape \textcolor{profublue}{\textbf{Calificación aprobatoria:}} 50}

\noindent{\footnotesize\itshape \textcolor{profublue}{\textbf{Calificación de la respuesta:}} Pass}

\noindent{\footnotesize\itshape \textcolor{profublue}{\textbf{Explicación:}} 75/100: casi el mismo significado; The Agent answer includes all the important parts of the Expected response about what 'control' means. It explains the same ideas clearly and adds a bit more detail without leaving anything out or being wrong. It also mentions where the law says this, which is extra but not wrong. So, the Agent answer is very good and mostly matches the Expected response.}

\end{minipage}

\vspace{2.5em}
\noindent
\begin{minipage}{\textwidth}
\textcolor{profublue}{\textbf{Usuario (53):}}\\[4pt]
\fcolorbox{profublue!30}{profublue!3}{%
  \begin{minipage}{0.95\textwidth}
  \smallskip
  ?En qu? casos se requiere autorizaci?n de la Comisi?n para adquisici?n de acciones de una AFORE?
  \smallskip
  \end{minipage}}
\end{minipage}

\vspace{0.4em}
\noindent
\begin{minipage}{\textwidth}
\textcolor{profugold!80!black}{\textbf{Agente (53):}}\\[4pt]
\fcolorbox{profugold!40}{profugold!5}{%
  \begin{minipage}{0.95\textwidth}
  \smallskip
  **Respuesta breve y clara:**\par Se requiere autorización de la Comisión para adquirir acciones de una AFORE (Administradora de Fondos para el Retiro) en los siguientes casos:\par - Cuando la adquisición implique que el comprador tenga el 5\% o más del capital social de la administradora. - También se requiere autorización si la Comisión considera que varias personas físicas o morales actúan como una sola para efectos de la adquisición y alcanzan ese porcentaje. - La fusión de administradoras también debe ser autorizada por la Comisión. - Si la adquisición es menor al 5\%, solo se debe dar aviso previo a la Comisión y posteriormente informar cuando se realice la operación.\par **Fundamento legal y localización:**\par Esto se encuentra en el Artículo 23 de la Ley de los Sistemas de Ahorro para el Retiro, en el rango de páginas correspondiente.\par Ley del SAR [Páginas 17]
  \smallskip
  \end{minipage}}
\end{minipage}

\vspace{0.35em}
\noindent
\begin{minipage}{0.95\textwidth}
\setlength{\parskip}{0.25em}
\noindent{\footnotesize\itshape \textcolor{profublue}{\textbf{Respuesta esperada:}} Cuando la adquisici?n o incorporaci?n de nuevos accionistas implique participaci?n del adquirente en 5\% o m?s del capital social, o en fusiones de administradoras.}

\noindent{\footnotesize\itshape \textcolor{profublue}{\textbf{Método de evaluación:}} CompareMeaning}

\noindent{\footnotesize\itshape \textcolor{profublue}{\textbf{Calificación aprobatoria:}} 50}

\noindent{\footnotesize\itshape \textcolor{profublue}{\textbf{Calificación de la respuesta:}} Pass}

\noindent{\footnotesize\itshape \textcolor{profublue}{\textbf{Explicación:}} 75/100: casi el mismo significado; The Agent answer gives all the important information the Expected response has and adds more helpful details that are correct. It explains extra rules about groups acting together and what happens if the purchase is less than 5\%. It also mentions the law source. Because it adds useful info without leaving anything out or being wrong, it matches well but has some extra details.}

\end{minipage}

\vspace{2.5em}
\noindent
\begin{minipage}{\textwidth}
\textcolor{profublue}{\textbf{Usuario (54):}}\\[4pt]
\fcolorbox{profublue!30}{profublue!3}{%
  \begin{minipage}{0.95\textwidth}
  \smallskip
  ?Qu? requisitos generales deben cumplir los miembros del Comit? Consultivo y de Vigilancia?
  \smallskip
  \end{minipage}}
\end{minipage}

\vspace{0.4em}
\noindent
\begin{minipage}{\textwidth}
\textcolor{profugold!80!black}{\textbf{Agente (54):}}\\[4pt]
\fcolorbox{profugold!40}{profugold!5}{%
  \begin{minipage}{0.95\textwidth}
  \smallskip
  **Requisitos generales para ser miembro del Comité Consultivo y de Vigilancia**\par Los miembros del Comité Consultivo y de Vigilancia deben cumplir con los siguientes requisitos:\par - Ser ciudadanos mexicanos por nacimiento que no adquieran otra nacionalidad y estar en pleno goce de sus derechos civiles y políticos. - Tener conocimientos en materia financiera, jurídica o de seguridad social. - Acreditar el nombramiento respectivo de la dependencia, entidad u organización que los proponga.\par **Fundamento legal:**   Estos requisitos se encuentran en el Artículo 14 de la Ley del SAR, ubicado en las páginas 10 y siguientes del texto oficial.   Ley del SAR [Páginas 10--11]
  \smallskip
  \end{minipage}}
\end{minipage}

\vspace{0.35em}
\noindent
\begin{minipage}{0.95\textwidth}
\setlength{\parskip}{0.25em}
\noindent{\footnotesize\itshape \textcolor{profublue}{\textbf{Respuesta esperada:}} Ser mexicanos por nacimiento, tener conocimientos financieros/jur?dicos/seguridad social, acreditar nombramiento y no ser funcionario o consejero de un participante del sistema.}

\noindent{\footnotesize\itshape \textcolor{profublue}{\textbf{Método de evaluación:}} CompareMeaning}

\noindent{\footnotesize\itshape \textcolor{profublue}{\textbf{Calificación aprobatoria:}} 50}

\noindent{\footnotesize\itshape \textcolor{profublue}{\textbf{Calificación de la respuesta:}} Pass}

\noindent{\footnotesize\itshape \textcolor{profublue}{\textbf{Explicación:}} 50/100: algún significado compartido; The agent answer gives most of the important requirements but leaves out that members can't be officials or counselors of a system participant. It also adds some extra details and legal references that weren't asked for. Because it misses a key rule, it's only a partial match.}

\end{minipage}

\vspace{2.5em}
\noindent
\begin{minipage}{\textwidth}
\textcolor{profublue}{\textbf{Usuario (55):}}\\[4pt]
\fcolorbox{profublue!30}{profublue!3}{%
  \begin{minipage}{0.95\textwidth}
  \smallskip
  ?Qui?n es responsable por la debida prestaci?n de servicios complementarios o auxiliares donde participe una administradora?
  \smallskip
  \end{minipage}}
\end{minipage}

\vspace{0.4em}
\noindent
\begin{minipage}{\textwidth}
\textcolor{profugold!80!black}{\textbf{Agente (55):}}\\[4pt]
\fcolorbox{profugold!40}{profugold!5}{%
  \begin{minipage}{0.95\textwidth}
  \smallskip
  **Respuesta breve y clara:**\par La administradora es responsable de la debida prestación de los servicios complementarios o auxiliares en los que participe. Además, si tiene participación accionaria en empresas que presten estos servicios, la administradora será solidariamente responsable de las sanciones que correspondan a dichas empresas por motivo de su supervisión.\par **Fundamento y localización en la Ley del SAR:**\par - El artículo 34 establece que la administradora debe obtener autorización para invertir en empresas que presten servicios complementarios o auxiliares y que, aunque estas empresas estén sujetas a regulación y supervisión de la Comisión, la administradora sigue siendo responsable de la debida prestación de los servicios. - También se menciona que la administradora será solidariamente responsable de las sanciones aplicables a dichas empresas.\par Referencia: Ley del SAR [Páginas 19--20].
  \smallskip
  \end{minipage}}
\end{minipage}

\vspace{0.35em}
\noindent
\begin{minipage}{0.95\textwidth}
\setlength{\parskip}{0.25em}
\noindent{\footnotesize\itshape \textcolor{profublue}{\textbf{Respuesta esperada:}} La administradora es responsable y puede ser solidariamente responsable de sanciones vinculadas a esas empresas.}

\noindent{\footnotesize\itshape \textcolor{profublue}{\textbf{Método de evaluación:}} CompareMeaning}

\noindent{\footnotesize\itshape \textcolor{profublue}{\textbf{Calificación aprobatoria:}} 50}

\noindent{\footnotesize\itshape \textcolor{profublue}{\textbf{Calificación de la respuesta:}} Pass}

\noindent{\footnotesize\itshape \textcolor{profublue}{\textbf{Explicación:}} 75/100: casi el mismo significado; The Agent answer gives all the important information the Expected response has and adds helpful details about the law. It correctly explains who is responsible and when, without leaving anything out or adding wrong facts. The extra info helps understand the answer better.}

\end{minipage}

\vspace{2.5em}
\noindent
\begin{minipage}{\textwidth}
\textcolor{profublue}{\textbf{Usuario (56):}}\\[4pt]
\fcolorbox{profublue!30}{profublue!3}{%
  \begin{minipage}{0.95\textwidth}
  \smallskip
  ?Qu? informaci?n adicional deben incluir las administradoras en los estados de cuenta de trabajadores afiliados seg?n el art?culo 18 bis?
  \smallskip
  \end{minipage}}
\end{minipage}

\vspace{0.4em}
\noindent
\begin{minipage}{\textwidth}
\textcolor{profugold!80!black}{\textbf{Agente (56):}}\\[4pt]
\fcolorbox{profugold!40}{profugold!5}{%
  \begin{minipage}{0.95\textwidth}
  \smallskip
  **Respuesta breve y clara:**\par Las administradoras deben incluir en los estados de cuenta de los trabajadores afiliados, sin costo adicional, la siguiente información adicional: - El salario base de cotización. - El número de días laborados declarados ante el Instituto Mexicano del Seguro Social para efecto del pago de cuotas.\par Además, en caso de discrepancia entre el salario recibido, su forma de integración o los días laborados por el trabajador con los declarados por el patrón, el trabajador puede denunciarlo ante las autoridades competentes.\par **Fundamento y localización en la Ley del SAR:**\par Esta obligación está establecida en el Artículo 18 bis, que se encuentra en la Página 14 del texto de la Ley del SAR.\par Referencia: Ley del SAR [Página 14] [1]\par [1]: cite:1 "Citation-1"
  \smallskip
  \end{minipage}}
\end{minipage}

\vspace{0.35em}
\noindent
\begin{minipage}{0.95\textwidth}
\setlength{\parskip}{0.25em}
\noindent{\footnotesize\itshape \textcolor{profublue}{\textbf{Respuesta esperada:}} El salario base de cotizaci?n y el n?mero de d?as laborados declarados ante el IMSS, sin costo adicional.}

\noindent{\footnotesize\itshape \textcolor{profublue}{\textbf{Método de evaluación:}} CompareMeaning}

\noindent{\footnotesize\itshape \textcolor{profublue}{\textbf{Calificación aprobatoria:}} 50}

\noindent{\footnotesize\itshape \textcolor{profublue}{\textbf{Calificación de la respuesta:}} Pass}

\noindent{\footnotesize\itshape \textcolor{profublue}{\textbf{Explicación:}} 75/100: casi el mismo significado; The Agent answer gives all the important information the Expected response has, like the salary base and days worked reported to IMSS, and adds helpful extra details about workers' rights and legal references. It doesn't leave out anything important or say anything wrong. So, it matches very well and adds useful information.}

\end{minipage}

\vspace{2.5em}
\noindent
\begin{minipage}{\textwidth}
\textcolor{profublue}{\textbf{Usuario (57):}}\\[4pt]
\fcolorbox{profublue!30}{profublue!3}{%
  \begin{minipage}{0.95\textwidth}
  \smallskip
  ?Qu? car?cter tiene la Ley de los Sistemas de Ahorro para el Retiro?
  \smallskip
  \end{minipage}}
\end{minipage}

\vspace{0.4em}
\noindent
\begin{minipage}{\textwidth}
\textcolor{profugold!80!black}{\textbf{Agente (57):}}\\[4pt]
\fcolorbox{profugold!40}{profugold!5}{%
  \begin{minipage}{0.95\textwidth}
  \smallskip
  **Respuesta breve y clara:**\par La Ley de los Sistemas de Ahorro para el Retiro tiene carácter de orden público e interés social. Esto significa que sus disposiciones son obligatorias y buscan proteger el bienestar colectivo, regulando el funcionamiento de los sistemas de ahorro para el retiro y sus participantes.\par **Fundamento y localización en la Ley:**\par - El carácter de la Ley se establece en el Artículo 1o, Capítulo I, Disposiciones Preliminares. - Se encuentra en la Página 1 del texto.\par **Referencia:**   Ley del SAR [Página 1]
  \smallskip
  \end{minipage}}
\end{minipage}

\vspace{0.35em}
\noindent
\begin{minipage}{0.95\textwidth}
\setlength{\parskip}{0.25em}
\noindent{\footnotesize\itshape \textcolor{profublue}{\textbf{Respuesta esperada:}} Es de orden p?blico e inter?s social.}

\noindent{\footnotesize\itshape \textcolor{profublue}{\textbf{Método de evaluación:}} CompareMeaning}

\noindent{\footnotesize\itshape \textcolor{profublue}{\textbf{Calificación aprobatoria:}} 50}

\noindent{\footnotesize\itshape \textcolor{profublue}{\textbf{Calificación de la respuesta:}} Pass}

\noindent{\footnotesize\itshape \textcolor{profublue}{\textbf{Explicación:}} 75/100: casi el mismo significado; The Agent answer gives the main point the Expected response has: the law is of public order and social interest. It also adds helpful details like where to find this in the law and what it means. These extra details are correct and useful, so the Agent answer matches well but includes more than needed.}

\end{minipage}

\vspace{2.5em}
\noindent
\begin{minipage}{\textwidth}
\textcolor{profublue}{\textbf{Usuario (58):}}\\[4pt]
\fcolorbox{profublue!30}{profublue!3}{%
  \begin{minipage}{0.95\textwidth}
  \smallskip
  ?Qu? perfil deben tener los consejeros independientes de las administradoras?
  \smallskip
  \end{minipage}}
\end{minipage}

\vspace{0.4em}
\noindent
\begin{minipage}{\textwidth}
\textcolor{profugold!80!black}{\textbf{Agente (58):}}\\[4pt]
\fcolorbox{profugold!40}{profugold!5}{%
  \begin{minipage}{0.95\textwidth}
  \smallskip
  **Perfil de los consejeros independientes de las administradoras**\par Los consejeros independientes de las administradoras deben cumplir con los siguientes requisitos:\par - Ser personas de reconocido prestigio en materia financiera, económica, jurídica o de seguridad social, y contar con experiencia profesional previa en la materia de al menos cinco años. - Acreditar ante la Comisión solvencia moral, así como capacidad técnica y administrativa. - No ser cónyuge ni tener parentesco por afinidad, civil o consanguíneo dentro del segundo grado, ni vínculo laboral con los accionistas de control o principales funcionarios de las administradoras. Tampoco deben ser accionistas o empleados de ninguna empresa del grupo financiero o corporativo al que pertenezca el accionista de control mayoritario de la administradora. - No prestar servicios personales a los institutos de seguridad social ni haberlos prestado durante los doce meses anteriores a su contratación. - Residir en territorio nacional. - Contar con la aprobación del Comité Consultivo y de Vigilancia de la Comisión.\par Además, los consejeros independientes no pueden ejercer simultáneamente su función en más de una administradora.\par **Fundamento legal:**   Estos requisitos se encuentran en el Artículo 50 de la Ley de los Sistemas de Ahorro para el Retiro, abarcando las páginas 30 y 31.   Ley del SAR [Páginas 30--31].
  \smallskip
  \end{minipage}}
\end{minipage}

\vspace{0.35em}
\noindent
\begin{minipage}{0.95\textwidth}
\setlength{\parskip}{0.25em}
\noindent{\footnotesize\itshape \textcolor{profublue}{\textbf{Respuesta esperada:}} Ser expertos en materia financiera, econ?mica, jur?dica o de seguridad social y no tener nexo patrimonial ni v?nculo laboral con accionistas de control o funcionarios.}

\noindent{\footnotesize\itshape \textcolor{profublue}{\textbf{Método de evaluación:}} CompareMeaning}

\noindent{\footnotesize\itshape \textcolor{profublue}{\textbf{Calificación aprobatoria:}} 50}

\noindent{\footnotesize\itshape \textcolor{profublue}{\textbf{Calificación de la respuesta:}} Pass}

\noindent{\footnotesize\itshape \textcolor{profublue}{\textbf{Explicación:}} 75/100: casi el mismo significado; The agent answer includes all the important points the expected response has and adds more useful details. It does not leave out anything important or say anything wrong. Because it gives extra information, it is a very good match but with some extra details.}

\end{minipage}

\vspace{2.5em}
\noindent
\begin{minipage}{\textwidth}
\textcolor{profublue}{\textbf{Usuario (59):}}\\[4pt]
\fcolorbox{profublue!30}{profublue!3}{%
  \begin{minipage}{0.95\textwidth}
  \smallskip
  ?Qu? reglas puede dictar la Comisi?n relacionadas con retiros programados?
  \smallskip
  \end{minipage}}
\end{minipage}

\vspace{0.4em}
\noindent
\begin{minipage}{\textwidth}
\textcolor{profugold!80!black}{\textbf{Agente (59):}}\\[4pt]
\fcolorbox{profugold!40}{profugold!5}{%
  \begin{minipage}{0.95\textwidth}
  \smallskip
  **Respuesta breve y clara:**\par La Comisión Nacional del Sistema de Ahorro para el Retiro (CONSAR) tiene la facultad de emitir reglas de carácter general para la operación y pago de los retiros programados. Esto significa que puede establecer disposiciones que regulen cómo se realizan y pagan los retiros programados dentro de los sistemas de ahorro para el retiro.\par **Fundamento y localización en la Ley del SAR:**\par - Capítulo II ``De la Comisión Nacional del Sistema de Ahorro para el Retiro'' - Artículo 5o, fracción IV: ``Emitir reglas de carácter general para la operación y pago de los retiros programados.'' - Ley del SAR [Páginas 2--3]
  \smallskip
  \end{minipage}}
\end{minipage}

\vspace{0.35em}
\noindent
\begin{minipage}{0.95\textwidth}
\setlength{\parskip}{0.25em}
\noindent{\footnotesize\itshape \textcolor{profublue}{\textbf{Respuesta esperada:}} Reglas de car?cter general para la operaci?n y pago de los retiros programados.}

\noindent{\footnotesize\itshape \textcolor{profublue}{\textbf{Método de evaluación:}} CompareMeaning}

\noindent{\footnotesize\itshape \textcolor{profublue}{\textbf{Calificación aprobatoria:}} 50}

\noindent{\footnotesize\itshape \textcolor{profublue}{\textbf{Calificación de la respuesta:}} Pass}

\noindent{\footnotesize\itshape \textcolor{profublue}{\textbf{Explicación:}} 75/100: casi el mismo significado; The agent answer gives all the important information the expected response has and adds helpful details about the law and where the rules come from. It doesn't leave out anything or say anything wrong. The extra details make the answer clearer and more complete.}

\end{minipage}

\vspace{2.5em}
\noindent
\begin{minipage}{\textwidth}
\textcolor{profublue}{\textbf{Usuario (60):}}\\[4pt]
\fcolorbox{profublue!30}{profublue!3}{%
  \begin{minipage}{0.95\textwidth}
  \smallskip
  ?Cu?les son los ?rganos de gobierno de la Comisi?n?
  \smallskip
  \end{minipage}}
\end{minipage}

\vspace{0.4em}
\noindent
\begin{minipage}{\textwidth}
\textcolor{profugold!80!black}{\textbf{Agente (60):}}\\[4pt]
\fcolorbox{profugold!40}{profugold!5}{%
  \begin{minipage}{0.95\textwidth}
  \smallskip
  Respuesta breve: Los órganos de gobierno de la Comisión, según la Ley de los Sistemas de Ahorro para el Retiro, son la Junta de Gobierno, la Presidencia y el Comité Consultivo y de Vigilancia.\par Fundamento: Esto se encuentra en el Artículo 6o., Sección II, donde se establece explícitamente la estructura de los órganos de gobierno de la Comisión.\par Referencia: Ley del SAR [Página 6]
  \smallskip
  \end{minipage}}
\end{minipage}

\vspace{0.35em}
\noindent
\begin{minipage}{0.95\textwidth}
\setlength{\parskip}{0.25em}
\noindent{\footnotesize\itshape \textcolor{profublue}{\textbf{Respuesta esperada:}} La Junta de Gobierno, la Presidencia y el Comit? Consultivo y de Vigilancia.}

\noindent{\footnotesize\itshape \textcolor{profublue}{\textbf{Método de evaluación:}} CompareMeaning}

\noindent{\footnotesize\itshape \textcolor{profublue}{\textbf{Calificación aprobatoria:}} 50}

\noindent{\footnotesize\itshape \textcolor{profublue}{\textbf{Calificación de la respuesta:}} Pass}

\noindent{\footnotesize\itshape \textcolor{profublue}{\textbf{Explicación:}} 75/100: casi el mismo significado; The agent answer gives all the main parts of the Commission's governing bodies just like the expected response. It also adds helpful details about the law and where the information comes from. This extra information is correct and useful, so the answer matches well and is complete.}

\end{minipage}

\vspace{2.5em}
\noindent
\begin{minipage}{\textwidth}
\textcolor{profublue}{\textbf{Usuario (61):}}\\[4pt]
\fcolorbox{profublue!30}{profublue!3}{%
  \begin{minipage}{0.95\textwidth}
  \smallskip
  ?Qu? tipo de regulaci?n puede emitir la Comisi?n dentro de su competencia?
  \smallskip
  \end{minipage}}
\end{minipage}

\vspace{0.4em}
\noindent
\begin{minipage}{\textwidth}
\textcolor{profugold!80!black}{\textbf{Agente (61):}}\\[4pt]
\fcolorbox{profugold!40}{profugold!5}{%
  \begin{minipage}{0.95\textwidth}
  \smallskip
  **Respuesta breve y clara:**\par La Comisión Nacional del Sistema de Ahorro para el Retiro (Comisión) tiene la facultad de coordinar, regular, supervisar y vigilar los sistemas de ahorro para el retiro y sus participantes, conforme a lo establecido en la Ley del SAR. Esto significa que puede emitir regulaciones, disposiciones y lineamientos dentro de su competencia para asegurar el correcto funcionamiento y supervisión de dichos sistemas.\par **Fundamento y localización en la Ley:**\par - El fundamento se encuentra en el Artículo 2o, que señala expresamente las facultades de la Comisión para regular, coordinar, supervisar y vigilar los sistemas de ahorro para el retiro. - Este contenido está localizado en las primeras páginas del texto.\par **Referencia:** Ley del SAR [Páginas 1--2]
  \smallskip
  \end{minipage}}
\end{minipage}

\vspace{0.35em}
\noindent
\begin{minipage}{0.95\textwidth}
\setlength{\parskip}{0.25em}
\noindent{\footnotesize\itshape \textcolor{profublue}{\textbf{Respuesta esperada:}} Regulaci?n prudencial aplicable a los participantes en los sistemas de ahorro para el retiro.}

\noindent{\footnotesize\itshape \textcolor{profublue}{\textbf{Método de evaluación:}} CompareMeaning}

\noindent{\footnotesize\itshape \textcolor{profublue}{\textbf{Calificación aprobatoria:}} 50}

\noindent{\footnotesize\itshape \textcolor{profublue}{\textbf{Calificación de la respuesta:}} Pass}

\noindent{\footnotesize\itshape \textcolor{profublue}{\textbf{Explicación:}} 75/100: casi el mismo significado; The Expected response says the Commission can issue 'prudential regulation' for participants. The Agent answer explains the Commission's powers and says it can make rules but doesn't use the exact term 'prudential regulation.' It also adds extra details about the law. So, the Agent answer is mostly right but misses the exact key phrase and has some extra info.}

\end{minipage}

\vspace{2.5em}
\noindent
\begin{minipage}{\textwidth}
\textcolor{profublue}{\textbf{Usuario (62):}}\\[4pt]
\fcolorbox{profublue!30}{profublue!3}{%
  \begin{minipage}{0.95\textwidth}
  \smallskip
  ?Qu? puede hacer el trabajador si detecta discrepancias entre su salario/d?as laborados y lo declarado por el patr?n?
  \smallskip
  \end{minipage}}
\end{minipage}

\vspace{0.4em}
\noindent
\begin{minipage}{\textwidth}
\textcolor{profugold!80!black}{\textbf{Agente (62):}}\\[4pt]
\fcolorbox{profugold!40}{profugold!5}{%
  \begin{minipage}{0.95\textwidth}
  \smallskip
  Respuesta breve y clara: Si un trabajador detecta discrepancias entre su salario o los días laborados y lo que el patrón ha declarado, puede denunciar esta situación ante las autoridades competentes. La Ley del SAR establece que las administradoras deben incluir en los estados de cuenta el salario base de cotización y el número de días laborados declarados ante el IMSS, y en caso de diferencias, el trabajador tiene derecho a presentar una denuncia.\par Fundamento legal: Esto se encuentra en el Artículo 18 bis, que señala expresamente el derecho del trabajador a denunciar discrepancias entre el salario recibido, su integración o los días laborados, y lo declarado por el patrón.\par Referencia: Ley del SAR [Páginas 14--15]
  \smallskip
  \end{minipage}}
\end{minipage}

\vspace{0.35em}
\noindent
\begin{minipage}{0.95\textwidth}
\setlength{\parskip}{0.25em}
\noindent{\footnotesize\itshape \textcolor{profublue}{\textbf{Respuesta esperada:}} Denunciarlo ante las autoridades competentes.}

\noindent{\footnotesize\itshape \textcolor{profublue}{\textbf{Método de evaluación:}} CompareMeaning}

\noindent{\footnotesize\itshape \textcolor{profublue}{\textbf{Calificación aprobatoria:}} 50}

\noindent{\footnotesize\itshape \textcolor{profublue}{\textbf{Calificación de la respuesta:}} Pass}

\noindent{\footnotesize\itshape \textcolor{profublue}{\textbf{Explicación:}} 75/100: casi el mismo significado; The Agent answer correctly says the worker can report the problem to the right authorities, which is what the Expected response says. It also gives more helpful details about the law and the worker's rights. This extra information is useful and does not change or leave out anything important.}

\end{minipage}

\vspace{2.5em}
\noindent
\begin{minipage}{\textwidth}
\textcolor{profublue}{\textbf{Usuario (63):}}\\[4pt]
\fcolorbox{profublue!30}{profublue!3}{%
  \begin{minipage}{0.95\textwidth}
  \smallskip
  ?Qu? limitaci?n de nombres y s?mbolos aplica a la denominaci?n de una AFORE?
  \smallskip
  \end{minipage}}
\end{minipage}

\vspace{0.4em}
\noindent
\begin{minipage}{\textwidth}
\textcolor{profugold!80!black}{\textbf{Agente (63):}}\\[4pt]
\fcolorbox{profugold!40}{profugold!5}{%
  \begin{minipage}{0.95\textwidth}
  \smallskip
  Limitaciones para la denominación de una AFORE\par Respuesta breve y clara: La Ley del SAR establece que las Administradoras de Fondos para el Retiro (AFORE) deben incluir en su denominación la expresión ``Administradora de Fondos para el Retiro'' o su abreviatura ``AFORE''. Además, está prohibido utilizar en su nombre expresiones en idioma extranjero, el nombre de alguna asociación religiosa o política, así como símbolos religiosos o patrios que sean objeto de devoción o culto público.\par Fundamento legal y localización: Esta limitación se encuentra en el Artículo 20, fracción I, de la Ley de los Sistemas de Ahorro para el Retiro, abarcando las páginas 14 y 15 del texto.\par Referencia: Ley del SAR [Páginas 14--15].
  \smallskip
  \end{minipage}}
\end{minipage}

\vspace{0.35em}
\noindent
\begin{minipage}{0.95\textwidth}
\setlength{\parskip}{0.25em}
\noindent{\footnotesize\itshape \textcolor{profublue}{\textbf{Respuesta esperada:}} No deben usar expresiones en idioma extranjero ni el nombre de asociaciones religiosas o pol?ticas, ni s?mbolos religiosos o patrios objeto de culto p?blico.}

\noindent{\footnotesize\itshape \textcolor{profublue}{\textbf{Método de evaluación:}} CompareMeaning}

\noindent{\footnotesize\itshape \textcolor{profublue}{\textbf{Calificación aprobatoria:}} 50}

\noindent{\footnotesize\itshape \textcolor{profublue}{\textbf{Calificación de la respuesta:}} Pass}

\noindent{\footnotesize\itshape \textcolor{profublue}{\textbf{Explicación:}} 75/100: casi el mismo significado; The Agent answer includes all the important limits from the Expected response and adds more useful details about the required name and legal source. It doesn't leave out anything important and doesn't have wrong info. It just gives a bit more than needed, which is helpful but not necessary.}

\end{minipage}

\vspace{2.5em}
\noindent
\begin{minipage}{\textwidth}
\textcolor{profublue}{\textbf{Usuario (64):}}\\[4pt]
\fcolorbox{profublue!30}{profublue!3}{%
  \begin{minipage}{0.95\textwidth}
  \smallskip
  ?Qui?n nombra al contralor normativo y a qui?n reporta?
  \smallskip
  \end{minipage}}
\end{minipage}

\vspace{0.4em}
\noindent
\begin{minipage}{\textwidth}
\textcolor{profugold!80!black}{\textbf{Agente (64):}}\\[4pt]
\fcolorbox{profugold!40}{profugold!5}{%
  \begin{minipage}{0.95\textwidth}
  \smallskip
  **Respuesta breve y clara:**\par El contralor normativo es nombrado por la asamblea de accionistas de la administradora. Este funcionario reporta únicamente al consejo de administración y a la asamblea de accionistas de la administradora correspondiente, sin estar subordinado a ningún otro órgano social ni funcionario de la administradora.\par **Fundamento y localización en la Ley del SAR:**\par - Artículo 30: Establece que el contralor normativo debe ser nombrado por la asamblea de accionistas, y que reporta exclusivamente al consejo de administración y a la asamblea de accionistas. - Página 18.\par **Referencia:**   Ley del SAR [Página 18]
  \smallskip
  \end{minipage}}
\end{minipage}

\vspace{0.35em}
\noindent
\begin{minipage}{0.95\textwidth}
\setlength{\parskip}{0.25em}
\noindent{\footnotesize\itshape \textcolor{profublue}{\textbf{Respuesta esperada:}} Lo nombra la asamblea de accionistas y reporta ?nicamente al consejo de administraci?n y a la asamblea.}

\noindent{\footnotesize\itshape \textcolor{profublue}{\textbf{Método de evaluación:}} CompareMeaning}

\noindent{\footnotesize\itshape \textcolor{profublue}{\textbf{Calificación aprobatoria:}} 50}

\noindent{\footnotesize\itshape \textcolor{profublue}{\textbf{Calificación de la respuesta:}} Pass}

\noindent{\footnotesize\itshape \textcolor{profublue}{\textbf{Explicación:}} 75/100: casi el mismo significado; The agent answer gives all the important information the expected response has about who appoints the contralor and who they report to. It also adds useful extra details and legal references that are correct and helpful. There are no mistakes or missing parts, so the agent answer is very good and complete.}

\end{minipage}

\vspace{2.5em}
\noindent
\begin{minipage}{\textwidth}
\textcolor{profublue}{\textbf{Usuario (65):}}\\[4pt]
\fcolorbox{profublue!30}{profublue!3}{%
  \begin{minipage}{0.95\textwidth}
  \smallskip
  ?Qu? atribuci?n tiene la Comisi?n en materia sancionatoria?
  \smallskip
  \end{minipage}}
\end{minipage}

\vspace{0.4em}
\noindent
\begin{minipage}{\textwidth}
\textcolor{profugold!80!black}{\textbf{Agente (65):}}\\[4pt]
\fcolorbox{profugold!40}{profugold!5}{%
  \begin{minipage}{0.95\textwidth}
  \smallskip
  **Atribuciones de la Comisión en materia sancionatoria**\par La Comisión tiene la facultad de imponer sanciones a las personas que participan en los sistemas de ahorro para el retiro. Estas atribuciones incluyen:\par - Determinar la amonestación, suspensión o remoción de consejeros, contralores normativos, directivos, comisarios, apoderados, funcionarios y demás personas que presten servicios a las administradoras o sociedades de inversión. - Inhabilitar a dichas personas para desempeñar empleos, cargos o comisiones dentro del sistema financiero mexicano o en entidades que participen en los sistemas de ahorro para el retiro, por un periodo de seis meses a diez años, cuando incurran de manera grave o reiterada en infracciones a las leyes y disposiciones normativas. - Para imponer la inhabilitación, la Comisión debe considerar la gravedad de la infracción, el nivel jerárquico, antecedentes, antigüedad, reincidencia y el monto de los daños o beneficios económicos derivados de la infracción.\par **Fundamento legal:**   Estas atribuciones se encuentran en el Artículo 52 de la Ley del SAR, y el procedimiento está detallado en las páginas correspondientes.\par **Referencia:**   Ley del SAR [Páginas 31--32]
  \smallskip
  \end{minipage}}
\end{minipage}

\vspace{0.35em}
\noindent
\begin{minipage}{0.95\textwidth}
\setlength{\parskip}{0.25em}
\noindent{\footnotesize\itshape \textcolor{profublue}{\textbf{Respuesta esperada:}} Imponer multas y sanciones y emitir opini?n a la autoridad competente sobre delitos previstos en la Ley.}

\noindent{\footnotesize\itshape \textcolor{profublue}{\textbf{Método de evaluación:}} CompareMeaning}

\noindent{\footnotesize\itshape \textcolor{profublue}{\textbf{Calificación aprobatoria:}} 50}

\noindent{\footnotesize\itshape \textcolor{profublue}{\textbf{Calificación de la respuesta:}} Pass}

\noindent{\footnotesize\itshape \textcolor{profublue}{\textbf{Explicación:}} 50/100: algún significado compartido; The agent answer explains many details about the Commission's sanction powers but leaves out an important part: that the Commission also gives opinions to authorities about crimes. It also adds extra details that the question didn't ask for. So, the agent answer is partly right but misses some key points.}

\end{minipage}

\vspace{2.5em}
\noindent
\begin{minipage}{\textwidth}
\textcolor{profublue}{\textbf{Usuario (66):}}\\[4pt]
\fcolorbox{profublue!30}{profublue!3}{%
  \begin{minipage}{0.95\textwidth}
  \smallskip
  ?Qu? periodicidad tienen las sesiones ordinarias de la Junta de Gobierno?
  \smallskip
  \end{minipage}}
\end{minipage}

\vspace{0.4em}
\noindent
\begin{minipage}{\textwidth}
\textcolor{profugold!80!black}{\textbf{Agente (66):}}\\[4pt]
\fcolorbox{profugold!40}{profugold!5}{%
  \begin{minipage}{0.95\textwidth}
  \smallskip
  Periodicidad de las sesiones ordinarias de la Junta de Gobierno\par Respuesta breve: Las sesiones ordinarias de la Junta de Gobierno se celebran de manera bimestral, es decir, cada dos meses. Además, pueden realizarse en cualquier momento si son convocadas por su Presidente o por el Presidente de la Comisión.\par Fundamento: Esta información se encuentra en el Artículo 9o, donde se establece la periodicidad bimestral y la posibilidad de convocatorias extraordinarias.\par Referencia: Ley del SAR [Página 7]
  \smallskip
  \end{minipage}}
\end{minipage}

\vspace{0.35em}
\noindent
\begin{minipage}{0.95\textwidth}
\setlength{\parskip}{0.25em}
\noindent{\footnotesize\itshape \textcolor{profublue}{\textbf{Respuesta esperada:}} Son bimestrales, y pueden celebrarse en cualquier tiempo cuando sean convocadas.}

\noindent{\footnotesize\itshape \textcolor{profublue}{\textbf{Método de evaluación:}} CompareMeaning}

\noindent{\footnotesize\itshape \textcolor{profublue}{\textbf{Calificación aprobatoria:}} 50}

\noindent{\footnotesize\itshape \textcolor{profublue}{\textbf{Calificación de la respuesta:}} Pass}

\noindent{\footnotesize\itshape \textcolor{profublue}{\textbf{Explicación:}} 100/100: mismo significado; The Agent answer gives all the important information the Expected response has. It says the meetings happen every two months and can be called anytime. It also adds helpful details about the law and explains the timing clearly. There are no mistakes or missing parts.}

\end{minipage}

\vspace{2.5em}
\noindent
\begin{minipage}{\textwidth}
\textcolor{profublue}{\textbf{Usuario (67):}}\\[4pt]
\fcolorbox{profublue!30}{profublue!3}{%
  \begin{minipage}{0.95\textwidth}
  \smallskip
  ?C?mo define la Ley del SAR a la Cuenta Individual?
  \smallskip
  \end{minipage}}
\end{minipage}

\vspace{0.4em}
\noindent
\begin{minipage}{\textwidth}
\textcolor{profugold!80!black}{\textbf{Agente (67):}}\\[4pt]
\fcolorbox{profugold!40}{profugold!5}{%
  \begin{minipage}{0.95\textwidth}
  \smallskip
  Definición breve y clara: La Ley del SAR define la Cuenta Individual como aquella de la que es titular un trabajador, en la cual se depositan las cuotas obrero patronales y estatales, así como sus rendimientos. Además, en ella se registran las aportaciones a los fondos de ahorro para el retiro.\par Fundamento y localización: Esta definición se encuentra en el Artículo 3o, fracción III bis, y abarca las páginas 1 y 2 del texto oficial.\par Referencia: Ley del SAR [Páginas 1--2]
  \smallskip
  \end{minipage}}
\end{minipage}

\vspace{0.35em}
\noindent
\begin{minipage}{0.95\textwidth}
\setlength{\parskip}{0.25em}
\noindent{\footnotesize\itshape \textcolor{profublue}{\textbf{Respuesta esperada:}} La cuenta de titularidad del trabajador donde se depositan cuotas obrero-patronales y estatales, rendimientos, aportaciones de vivienda y dem?s recursos permitidos.}

\noindent{\footnotesize\itshape \textcolor{profublue}{\textbf{Método de evaluación:}} CompareMeaning}

\noindent{\footnotesize\itshape \textcolor{profublue}{\textbf{Calificación aprobatoria:}} 50}

\noindent{\footnotesize\itshape \textcolor{profublue}{\textbf{Calificación de la respuesta:}} Pass}

\noindent{\footnotesize\itshape \textcolor{profublue}{\textbf{Explicación:}} 50/100: algún significado compartido; The agent answer explains the account well but leaves out that it also includes housing contributions and other allowed resources. It adds some extra details about retirement savings funds that the expected answer doesn't mention. So, it mostly gets the idea but misses some important parts and adds some extra info.}

\end{minipage}

\vspace{2.5em}
\noindent
\begin{minipage}{\textwidth}
\textcolor{profublue}{\textbf{Usuario (68):}}\\[4pt]
\fcolorbox{profublue!30}{profublue!3}{%
  \begin{minipage}{0.95\textwidth}
  \smallskip
  ?Qu? obligaci?n tienen los auditores externos de las administradoras frente a la Comisi?n?
  \smallskip
  \end{minipage}}
\end{minipage}

\vspace{0.4em}
\noindent
\begin{minipage}{\textwidth}
\textcolor{profugold!80!black}{\textbf{Agente (68):}}\\[4pt]
\fcolorbox{profugold!40}{profugold!5}{%
  \begin{minipage}{0.95\textwidth}
  \smallskip
  **Obligación de los auditores externos frente a la Comisión**\par Los auditores externos de las administradoras tienen la obligación de entregar a la Comisión la información que ésta les solicite sobre la situación de dichas entidades financieras. Además, deben informar a la Comisión sobre las irregularidades graves que encuentren en el desempeño de su labor.\par **Fundamento legal y localización** - Artículo 31 de la Ley del SAR. - Ley del SAR [Páginas 19--20].
  \smallskip
  \end{minipage}}
\end{minipage}

\vspace{0.35em}
\noindent
\begin{minipage}{0.95\textwidth}
\setlength{\parskip}{0.25em}
\noindent{\footnotesize\itshape \textcolor{profublue}{\textbf{Respuesta esperada:}} Entregar la informaci?n que solicite la Comisi?n e informar sobre irregularidades graves detectadas.}

\noindent{\footnotesize\itshape \textcolor{profublue}{\textbf{Método de evaluación:}} CompareMeaning}

\noindent{\footnotesize\itshape \textcolor{profublue}{\textbf{Calificación aprobatoria:}} 50}

\noindent{\footnotesize\itshape \textcolor{profublue}{\textbf{Calificación de la respuesta:}} Pass}

\noindent{\footnotesize\itshape \textcolor{profublue}{\textbf{Explicación:}} 75/100: casi el mismo significado; The Agent answer includes all the important points from the Expected response about what auditors must do. It also adds useful extra details like the law that supports these duties. Nothing important is left out, and the extra information helps explain the answer better.}

\end{minipage}

\vspace{2.5em}
\noindent
\begin{minipage}{\textwidth}
\textcolor{profublue}{\textbf{Usuario (69):}}\\[4pt]
\fcolorbox{profublue!30}{profublue!3}{%
  \begin{minipage}{0.95\textwidth}
  \smallskip
  ?La Comisi?n puede otorgar o revocar autorizaciones? ?A qui?nes aplica principalmente?
  \smallskip
  \end{minipage}}
\end{minipage}

\vspace{0.4em}
\noindent
\begin{minipage}{\textwidth}
\textcolor{profugold!80!black}{\textbf{Agente (69):}}\\[4pt]
\fcolorbox{profugold!40}{profugold!5}{%
  \begin{minipage}{0.95\textwidth}
  \smallskip
  Respuesta breve: Sí, la Comisión Nacional del Sistema de Ahorro para el Retiro (CONSAR) tiene la facultad de otorgar, modificar o revocar autorizaciones. Estas facultades aplican principalmente a las administradoras de fondos para el retiro y a las sociedades de inversión especializadas de fondos para el retiro.\par Fundamento: Esto se encuentra en el Artículo 5o., fracción VI, que establece que la Comisión puede ``otorgar, modificar o revocar las autorizaciones a que se refiere esta ley, a las administradoras y sociedades de inversión''.\par Referencia: Ley del SAR [Página 3].
  \smallskip
  \end{minipage}}
\end{minipage}

\vspace{0.35em}
\noindent
\begin{minipage}{0.95\textwidth}
\setlength{\parskip}{0.25em}
\noindent{\footnotesize\itshape \textcolor{profublue}{\textbf{Respuesta esperada:}} S?, puede otorgar, modificar o revocar autorizaciones a administradoras y sociedades de inversi?n.}

\noindent{\footnotesize\itshape \textcolor{profublue}{\textbf{Método de evaluación:}} CompareMeaning}

\noindent{\footnotesize\itshape \textcolor{profublue}{\textbf{Calificación aprobatoria:}} 50}

\noindent{\footnotesize\itshape \textcolor{profublue}{\textbf{Calificación de la respuesta:}} Pass}

\noindent{\footnotesize\itshape \textcolor{profublue}{\textbf{Explicación:}} 75/100: casi el mismo significado; The Agent answer correctly says the Commission can give, change, or take away permissions, just like the Expected response. It also says who this applies to, matching the Expected response. The Agent answer adds a legal reference, which is extra but useful information. There are no mistakes or missing parts, so the Agent answer is very good and matches well.}

\end{minipage}

\vspace{2.5em}
\noindent
\begin{minipage}{\textwidth}
\textcolor{profublue}{\textbf{Usuario (70):}}\\[4pt]
\fcolorbox{profublue!30}{profublue!3}{%
  \begin{minipage}{0.95\textwidth}
  \smallskip
  ?Qu? es el Rendimiento Neto seg?n la Ley del SAR?
  \smallskip
  \end{minipage}}
\end{minipage}

\vspace{0.4em}
\noindent
\begin{minipage}{\textwidth}
\textcolor{profugold!80!black}{\textbf{Agente (70):}}\\[4pt]
\fcolorbox{profugold!40}{profugold!5}{%
  \begin{minipage}{0.95\textwidth}
  \smallskip
  **Definición sencilla de Rendimiento Neto según la Ley del SAR**\par El Rendimiento Neto, de acuerdo con la Ley de los Sistemas de Ahorro para el Retiro, es el indicador que refleja los rendimientos obtenidos por los trabajadores por la inversión de sus recursos en las Sociedades de Inversión, menos las comisiones que se les cobran. Es decir, muestra cuánto ganan realmente los trabajadores después de descontar las comisiones.\par La metodología para calcular estos indicadores debe ser autorizada por la Junta de Gobierno de la Comisión Nacional del Sistema de Ahorro para el Retiro, y en dicha metodología se establece el periodo para su cálculo.\par **Fundamento legal y localización** - Artículo 3o, fracción V bis. - Ley del SAR [Páginas 1--2].
  \smallskip
  \end{minipage}}
\end{minipage}

\vspace{0.35em}
\noindent
\begin{minipage}{0.95\textwidth}
\setlength{\parskip}{0.25em}
\noindent{\footnotesize\itshape \textcolor{profublue}{\textbf{Respuesta esperada:}} Indicadores que reflejan rendimientos menos comisiones obtenidos por la inversi?n de recursos en Sociedades de Inversi?n, con metodolog?a autorizada por la Junta de Gobierno.}

\noindent{\footnotesize\itshape \textcolor{profublue}{\textbf{Método de evaluación:}} CompareMeaning}

\noindent{\footnotesize\itshape \textcolor{profublue}{\textbf{Calificación aprobatoria:}} 50}

\noindent{\footnotesize\itshape \textcolor{profublue}{\textbf{Calificación de la respuesta:}} Pass}

\noindent{\footnotesize\itshape \textcolor{profublue}{\textbf{Explicación:}} 75/100: casi el mismo significado; The agent answer includes all the important points from the expected response and adds clear explanations and legal references. It doesn't leave out anything important and doesn't say anything wrong. It just gives more helpful details than needed.}

\end{minipage}

\vspace{2.5em}
\noindent
\begin{minipage}{\textwidth}
\textcolor{profublue}{\textbf{Usuario (71):}}\\[4pt]
\fcolorbox{profublue!30}{profublue!3}{%
  \begin{minipage}{0.95\textwidth}
  \smallskip
  ?Qu? instituciones de seguridad social menciona expresamente el art?culo 1 como marco relacionado?
  \smallskip
  \end{minipage}}
\end{minipage}

\vspace{0.4em}
\noindent
\begin{minipage}{\textwidth}
\textcolor{profugold!80!black}{\textbf{Agente (71):}}\\[4pt]
\fcolorbox{profugold!40}{profugold!5}{%
  \begin{minipage}{0.95\textwidth}
  \smallskip
  Respuesta breve y clara: El artículo 1 de la Ley de los Sistemas de Ahorro para el Retiro menciona expresamente como marco relacionado las siguientes instituciones de seguridad social: - El Instituto Mexicano del Seguro Social (IMSS) - El Instituto del Fondo Nacional de la Vivienda para los Trabajadores (INFONAVIT) - El Instituto de Seguridad y Servicios Sociales de los Trabajadores del Estado (ISSSTE)\par Estas instituciones se consideran dentro del marco regulatorio de la Ley del SAR.\par Fundamento y localización: Esto se encuentra en el Artículo 1 y en las definiciones del Artículo 3, fracción VI, en las Páginas 1 y 2 del texto.\par Referencia: Ley del SAR [Páginas 1--2]
  \smallskip
  \end{minipage}}
\end{minipage}

\vspace{0.35em}
\noindent
\begin{minipage}{0.95\textwidth}
\setlength{\parskip}{0.25em}
\noindent{\footnotesize\itshape \textcolor{profublue}{\textbf{Respuesta esperada:}} El IMSS, el INFONAVIT y el ISSSTE, adem?s de las leyes del Seguro Social, INFONAVIT e ISSSTE.}

\noindent{\footnotesize\itshape \textcolor{profublue}{\textbf{Método de evaluación:}} CompareMeaning}

\noindent{\footnotesize\itshape \textcolor{profublue}{\textbf{Calificación aprobatoria:}} 50}

\noindent{\footnotesize\itshape \textcolor{profublue}{\textbf{Calificación de la respuesta:}} Pass}

\noindent{\footnotesize\itshape \textcolor{profublue}{\textbf{Explicación:}} 75/100: casi el mismo significado; The Agent answer lists the main institutions mentioned in the question, just like the Expected response. But it leaves out the part about the related laws, which is important. It also adds some extra helpful details that don't cause confusion. So, the Agent answer is mostly right but misses some key info about the laws.}

\end{minipage}

\vspace{2.5em}
\noindent
\begin{minipage}{\textwidth}
\textcolor{profublue}{\textbf{Usuario (72):}}\\[4pt]
\fcolorbox{profublue!30}{profublue!3}{%
  \begin{minipage}{0.95\textwidth}
  \smallskip
  ?Cu?ndo deben presentar las administradoras sus comisiones anuales para autorizaci?n?
  \smallskip
  \end{minipage}}
\end{minipage}

\vspace{0.4em}
\noindent
\begin{minipage}{\textwidth}
\textcolor{profugold!80!black}{\textbf{Agente (72):}}\\[4pt]
\fcolorbox{profugold!40}{profugold!5}{%
  \begin{minipage}{0.95\textwidth}
  \smallskip
  **Respuesta breve y clara:** Las administradoras deben presentar sus comisiones anuales para autorización ante la Junta de Gobierno de la Comisión dentro de los primeros diez días hábiles del mes de noviembre, para que sean aplicadas en el año calendario siguiente. Si no presentan en ese plazo, estarán obligadas a cobrar la comisión más baja autorizada hasta que presenten su solicitud y sea autorizada. \par **Fundamento y localización en la Ley del SAR:** Esta obligación se encuentra en el texto que regula la presentación y autorización de comisiones, especificando el plazo y las consecuencias en caso de omisión. El contenido relevante se localiza en la página 21 del documento.\par **Referencia:**   Ley del SAR [Página 21]
  \smallskip
  \end{minipage}}
\end{minipage}

\vspace{0.35em}
\noindent
\begin{minipage}{0.95\textwidth}
\setlength{\parskip}{0.25em}
\noindent{\footnotesize\itshape \textcolor{profublue}{\textbf{Respuesta esperada:}} Cada a?o dentro de los primeros diez d?as h?biles de noviembre, para aplicarlas en el a?o calendario siguiente.}

\noindent{\footnotesize\itshape \textcolor{profublue}{\textbf{Método de evaluación:}} CompareMeaning}

\noindent{\footnotesize\itshape \textcolor{profublue}{\textbf{Calificación aprobatoria:}} 50}

\noindent{\footnotesize\itshape \textcolor{profublue}{\textbf{Calificación de la respuesta:}} Pass}

\noindent{\footnotesize\itshape \textcolor{profublue}{\textbf{Explicación:}} 75/100: casi el mismo significado; The Agent answer gives all the important details from the Expected response about when the commissions must be presented and for what year. It also adds useful extra information about who they present to and what happens if they are late. This extra information is correct and helpful, so the Agent answer matches well but includes more than needed.}

\end{minipage}

\vspace{2.5em}
\noindent
\begin{minipage}{\textwidth}
\textcolor{profublue}{\textbf{Usuario (73):}}\\[4pt]
\fcolorbox{profublue!30}{profublue!3}{%
  \begin{minipage}{0.95\textwidth}
  \smallskip
  Menciona una funci?n espec?fica del contralor normativo prevista en la Ley.
  \smallskip
  \end{minipage}}
\end{minipage}

\vspace{0.4em}
\noindent
\begin{minipage}{\textwidth}
\textcolor{profugold!80!black}{\textbf{Agente (73):}}\\[4pt]
\fcolorbox{profugold!40}{profugold!5}{%
  \begin{minipage}{0.95\textwidth}
  \smallskip
  Respuesta breve: Una función específica del contralor normativo prevista en la Ley de los Sistemas de Ahorro para el Retiro es verificar que se cumpla el programa de autorregulación de la administradora, el cual incluye actividades, normas y acciones correctivas para garantizar el cumplimiento normativo, la operación eficiente y la protección de los intereses de los trabajadores, así como evitar conflictos de interés y el uso indebido de información privilegiada.\par Fundamento: Esta función se encuentra en el Artículo 30, fracción I, de la Ley del SAR, ubicada en la Página 18 del texto.\par Referencia: Ley del SAR [Página 18]
  \smallskip
  \end{minipage}}
\end{minipage}

\vspace{0.35em}
\noindent
\begin{minipage}{0.95\textwidth}
\setlength{\parskip}{0.25em}
\noindent{\footnotesize\itshape \textcolor{profublue}{\textbf{Respuesta esperada:}} Verificar el cumplimiento del programa de autorregulaci?n y proponer modificaciones para prevenir conflictos de inter?s y uso indebido de informaci?n.}

\noindent{\footnotesize\itshape \textcolor{profublue}{\textbf{Método de evaluación:}} CompareMeaning}

\noindent{\footnotesize\itshape \textcolor{profublue}{\textbf{Calificación aprobatoria:}} 50}

\noindent{\footnotesize\itshape \textcolor{profublue}{\textbf{Calificación de la respuesta:}} Pass}

\noindent{\footnotesize\itshape \textcolor{profublue}{\textbf{Explicación:}} 50/100: algún significado compartido; The Agent answer correctly says the controller checks the self-regulation program, like the Expected response. But it leaves out that the controller also suggests changes to avoid conflicts and misuse of information. The Agent answer adds extra details but misses this important part. So, it partly matches but misses some key info.}

\end{minipage}

\vspace{2.5em}
\noindent
\begin{minipage}{\textwidth}
\textcolor{profublue}{\textbf{Usuario (74):}}\\[4pt]
\fcolorbox{profublue!30}{profublue!3}{%
  \begin{minipage}{0.95\textwidth}
  \smallskip
  ?A qu? inter?s deben atender exclusivamente las administradoras al invertir recursos de trabajadores?
  \smallskip
  \end{minipage}}
\end{minipage}

\vspace{0.4em}
\noindent
\begin{minipage}{\textwidth}
\textcolor{profugold!80!black}{\textbf{Agente (74):}}\\[4pt]
\fcolorbox{profugold!40}{profugold!5}{%
  \begin{minipage}{0.95\textwidth}
  \smallskip
  Respuesta breve y clara: Las administradoras deben atender exclusivamente al interés de preservar el valor adquisitivo del ahorro de los trabajadores al invertir sus recursos. Es decir, su principal objetivo debe ser proteger el poder de compra de los fondos de los trabajadores, eligiendo valores y activos que cumplan con esta finalidad.\par Fundamento y localización: Esto se establece en el Artículo 47 de la Ley de los Sistemas de Ahorro para el Retiro, que señala que la cartera de inversión debe estar integrada fundamentalmente por valores cuyas características específicas preserven el valor adquisitivo del ahorro de los trabajadores. Ley del SAR [Página 28].
  \smallskip
  \end{minipage}}
\end{minipage}

\vspace{0.35em}
\noindent
\begin{minipage}{0.95\textwidth}
\setlength{\parskip}{0.25em}
\noindent{\footnotesize\itshape \textcolor{profublue}{\textbf{Respuesta esperada:}} Al inter?s de los trabajadores, buscando adecuada rentabilidad y seguridad en las inversiones.}

\noindent{\footnotesize\itshape \textcolor{profublue}{\textbf{Método de evaluación:}} CompareMeaning}

\noindent{\footnotesize\itshape \textcolor{profublue}{\textbf{Calificación aprobatoria:}} 50}

\noindent{\footnotesize\itshape \textcolor{profublue}{\textbf{Calificación de la respuesta:}} Pass}

\noindent{\footnotesize\itshape \textcolor{profublue}{\textbf{Explicación:}} 50/100: algún significado compartido; The agent answer talks about protecting the value of workers' savings but leaves out that investments should also aim for good returns and safety. It gives some correct details but misses important parts of the expected answer. So, it partly matches but is missing key information.}

\end{minipage}

\vspace{2.5em}
\noindent
\begin{minipage}{\textwidth}
\textcolor{profublue}{\textbf{Usuario (75):}}\\[4pt]
\fcolorbox{profublue!30}{profublue!3}{%
  \begin{minipage}{0.95\textwidth}
  \smallskip
  ?Qui?n tiene a su cargo la coordinaci?n, regulaci?n, supervisi?n y vigilancia de los sistemas de ahorro para el retiro?
  \smallskip
  \end{minipage}}
\end{minipage}

\vspace{0.4em}
\noindent
\begin{minipage}{\textwidth}
\textcolor{profugold!80!black}{\textbf{Agente (75):}}\\[4pt]
\fcolorbox{profugold!40}{profugold!5}{%
  \begin{minipage}{0.95\textwidth}
  \smallskip
  Respuesta breve: La coordinación, regulación, supervisión y vigilancia de los sistemas de ahorro para el retiro está a cargo de la Comisión Nacional del Sistema de Ahorro para el Retiro (CONSAR), que es un órgano administrativo desconcentrado de la Secretaría de Hacienda y Crédito Público, con autonomía técnica y facultades ejecutivas.\par Fundamento: Esto se establece en el Artículo 2o de la Ley de los Sistemas de Ahorro para el Retiro, ubicado en el Capítulo I, Disposiciones Preliminares, y puede consultarse en las páginas 1 y 2 del texto.\par Referencia: Ley del SAR [Páginas 1--2]
  \smallskip
  \end{minipage}}
\end{minipage}

\vspace{0.35em}
\noindent
\begin{minipage}{0.95\textwidth}
\setlength{\parskip}{0.25em}
\noindent{\footnotesize\itshape \textcolor{profublue}{\textbf{Respuesta esperada:}} La Comisi?n Nacional del Sistema de Ahorro para el Retiro (CONSAR), ?rgano administrativo desconcentrado de la SHCP con autonom?a t?cnica y facultades ejecutivas.}

\noindent{\footnotesize\itshape \textcolor{profublue}{\textbf{Método de evaluación:}} CompareMeaning}

\noindent{\footnotesize\itshape \textcolor{profublue}{\textbf{Calificación aprobatoria:}} 50}

\noindent{\footnotesize\itshape \textcolor{profublue}{\textbf{Calificación de la respuesta:}} Pass}

\noindent{\footnotesize\itshape \textcolor{profublue}{\textbf{Explicación:}} 75/100: casi el mismo significado; The agent answer gives all the important information expected: who is in charge and their role. It also adds helpful details about the law and where to find it. There are no mistakes or extra unrelated information. So, the agent answer matches very well with what was expected.}

\end{minipage}

\vspace{2.5em}
\noindent
\begin{minipage}{\textwidth}
\textcolor{profublue}{\textbf{Usuario (76):}}\\[4pt]
\fcolorbox{profublue!30}{profublue!3}{%
  \begin{minipage}{0.95\textwidth}
  \smallskip
  ?Cu?l es una facultad clave del Presidente en procesos judiciales donde participe la Comisi?n?
  \smallskip
  \end{minipage}}
\end{minipage}

\vspace{0.4em}
\noindent
\begin{minipage}{\textwidth}
\textcolor{profugold!80!black}{\textbf{Agente (76):}}\\[4pt]
\fcolorbox{profugold!40}{profugold!5}{%
  \begin{minipage}{0.95\textwidth}
  \smallskip
  **Respuesta breve y clara:**\par Una facultad clave del Presidente de la Comisión en procesos judiciales es tener a su cargo la representación legal de la Comisión. Esto incluye ejercer acciones, defensas, producir alegatos, ofrecer pruebas, interponer recursos, presentar desistimientos y realizar todos los actos procesales que correspondan, incluso en juicios de amparo. El Presidente puede delegar estas funciones a Vicepresidentes o Directores Generales mediante acuerdos delegatorios.\par **Fundamento y localización en la Ley del SAR:**\par - Artículo 12, fracción I: ``Tener a su cargo la representación legal de la Comisión y el ejercicio de sus facultades\ldots{} En los procedimientos judiciales, administrativos o laborales en los que la Comisión sea parte o pueda resultar afectada, el Presidente\ldots{} ejercerá las acciones, excepciones y defensas, producirá alegatos, ofrecerá pruebas, interpondrá los recursos que procedan, podrá presentar desistimientos, y en general realizará todos los actos procesales que correspondan a la Comisión o a sus órganos, incluyendo en los juicios de amparo la presentación de los informes de ley\ldots{}'' - El texto se encuentra en las Páginas 8 y siguientes.\par **Referencia:** Ley del SAR [Páginas 8--9]
  \smallskip
  \end{minipage}}
\end{minipage}

\vspace{0.35em}
\noindent
\begin{minipage}{0.95\textwidth}
\setlength{\parskip}{0.25em}
\noindent{\footnotesize\itshape \textcolor{profublue}{\textbf{Respuesta esperada:}} Ejercitar acciones, defensas, ofrecer pruebas e interponer recursos, incluyendo en juicios de amparo la presentaci?n de informes de ley.}

\noindent{\footnotesize\itshape \textcolor{profublue}{\textbf{Método de evaluación:}} CompareMeaning}

\noindent{\footnotesize\itshape \textcolor{profublue}{\textbf{Calificación aprobatoria:}} 50}

\noindent{\footnotesize\itshape \textcolor{profublue}{\textbf{Calificación de la respuesta:}} Pass}

\noindent{\footnotesize\itshape \textcolor{profublue}{\textbf{Explicación:}} 75/100: casi el mismo significado; The Agent answer includes all the important points from the Expected response and adds more helpful details about the President's role. It doesn't leave out anything important or say anything wrong. Because it gives extra correct information, it matches well but has some extra details.}

\end{minipage}

\vspace{2.5em}
\noindent
\begin{minipage}{\textwidth}
\textcolor{profublue}{\textbf{Usuario (77):}}\\[4pt]
\fcolorbox{profublue!30}{profublue!3}{%
  \begin{minipage}{0.95\textwidth}
  \smallskip
  ?Cu?l es la finalidad del Comit? Consultivo y de Vigilancia?
  \smallskip
  \end{minipage}}
\end{minipage}

\vspace{0.4em}
\noindent
\begin{minipage}{\textwidth}
\textcolor{profugold!80!black}{\textbf{Agente (77):}}\\[4pt]
\fcolorbox{profugold!40}{profugold!5}{%
  \begin{minipage}{0.95\textwidth}
  \smallskip
  Finalidad del Comité Consultivo y de Vigilancia\par Respuesta breve y clara: El Comité Consultivo y de Vigilancia tiene como finalidad principal conocer, vigilar y emitir opiniones sobre los asuntos relevantes de los sistemas de ahorro para el retiro. Entre sus funciones destacan: vigilar el desarrollo de los sistemas para prevenir conflictos de interés y prácticas monopólicas, conocer sobre la administración de cuentas individuales, autorizar la constitución y modificación de administradoras y sociedades de inversión, aprobar nombramientos de contralores normativos y consejeros independientes, y emitir opiniones sobre reglas generales de inversión y comisiones.\par Fundamento y localización: La finalidad y facultades del Comité Consultivo y de Vigilancia se encuentran en el Artículo 16 de la Ley de los Sistemas de Ahorro para el Retiro, que abarca desde la vigilancia, conocimiento y opinión sobre diversos aspectos operativos y normativos de los sistemas de ahorro para el retiro.\par Referencia: Ley del SAR [Páginas 10--11]
  \smallskip
  \end{minipage}}
\end{minipage}

\vspace{0.35em}
\noindent
\begin{minipage}{0.95\textwidth}
\setlength{\parskip}{0.25em}
\noindent{\footnotesize\itshape \textcolor{profublue}{\textbf{Respuesta esperada:}} Velar por los intereses de los sectores obrero, patronal y gobierno, asegurando armon?a y equilibrio para el mejor funcionamiento de los sistemas.}

\noindent{\footnotesize\itshape \textcolor{profublue}{\textbf{Método de evaluación:}} CompareMeaning}

\noindent{\footnotesize\itshape \textcolor{profublue}{\textbf{Calificación aprobatoria:}} 50}

\noindent{\footnotesize\itshape \textcolor{profublue}{\textbf{Calificación de la respuesta:}} Pass}

\noindent{\footnotesize\itshape \textcolor{profublue}{\textbf{Explicación:}} 50/100: algún significado compartido; The Expected response says the committee's main goal is to protect the interests of workers, employers, and the government and keep things balanced. The Agent answer talks about what the committee does and its legal rules but doesn't clearly say this main goal. So, the Agent answer is partly right but misses the most important point.}

\end{minipage}

\vspace{2.5em}
\noindent
\begin{minipage}{\textwidth}
\textcolor{profublue}{\textbf{Usuario (78):}}\\[4pt]
\fcolorbox{profublue!30}{profublue!3}{%
  \begin{minipage}{0.95\textwidth}
  \smallskip
  ?Qui?n nombra al Presidente de la Comisi?n?
  \smallskip
  \end{minipage}}
\end{minipage}

\vspace{0.4em}
\noindent
\begin{minipage}{\textwidth}
\textcolor{profugold!80!black}{\textbf{Agente (78):}}\\[4pt]
\fcolorbox{profugold!40}{profugold!5}{%
  \begin{minipage}{0.95\textwidth}
  \smallskip
  Respuesta breve: El Presidente de la Comisión es nombrado por el Secretario de Hacienda y Crédito Público.\par Fundamento: Esto se establece en el Artículo 10 de la Ley de los Sistemas de Ahorro para el Retiro, donde se indica expresamente que el nombramiento corresponde al Secretario de Hacienda y Crédito Público. El artículo también detalla los requisitos que debe cumplir el Presidente.\par Referencia: Ley del SAR [Páginas 7--8] [1]\par [1]: cite:1 "Citation-1"
  \smallskip
  \end{minipage}}
\end{minipage}

\vspace{0.35em}
\noindent
\begin{minipage}{0.95\textwidth}
\setlength{\parskip}{0.25em}
\noindent{\footnotesize\itshape \textcolor{profublue}{\textbf{Respuesta esperada:}} El Secretario de Hacienda y Cr?dito P?blico.}

\noindent{\footnotesize\itshape \textcolor{profublue}{\textbf{Método de evaluación:}} CompareMeaning}

\noindent{\footnotesize\itshape \textcolor{profublue}{\textbf{Calificación aprobatoria:}} 50}

\noindent{\footnotesize\itshape \textcolor{profublue}{\textbf{Calificación de la respuesta:}} Pass}

\noindent{\footnotesize\itshape \textcolor{profublue}{\textbf{Explicación:}} 75/100: casi el mismo significado; The Agent answer gives the main fact the user asked for: who appoints the President of the Commission. It also adds some extra helpful details about the law that the appointment follows. This extra information is correct and useful, so the answer is very good but includes more than just the simple answer.}

\end{minipage}

\vspace{2.5em}
\noindent
\begin{minipage}{\textwidth}
\textcolor{profublue}{\textbf{Usuario (79):}}\\[4pt]
\fcolorbox{profublue!30}{profublue!3}{%
  \begin{minipage}{0.95\textwidth}
  \smallskip
  ?Cu?l es el objeto de la Ley de los Sistemas de Ahorro para el Retiro?
  \smallskip
  \end{minipage}}
\end{minipage}

\vspace{0.4em}
\noindent
\begin{minipage}{\textwidth}
\textcolor{profugold!80!black}{\textbf{Agente (79):}}\\[4pt]
\fcolorbox{profugold!40}{profugold!5}{%
  \begin{minipage}{0.95\textwidth}
  \smallskip
  **Objeto de la Ley de los Sistemas de Ahorro para el Retiro**\par La Ley de los Sistemas de Ahorro para el Retiro tiene como objetivo regular el funcionamiento de los sistemas de ahorro para el retiro y sus participantes, conforme a lo previsto en la propia Ley y en las leyes del Seguro Social, del Instituto del Fondo Nacional de la Vivienda para los Trabajadores y del Instituto de Seguridad y Servicios Sociales de los Trabajadores del Estado.\par **Fundamento legal:**   El objeto se establece en el Artículo 1o, Capítulo I, Disposiciones Preliminares.   Ley del SAR [Página 1]
  \smallskip
  \end{minipage}}
\end{minipage}

\vspace{0.35em}
\noindent
\begin{minipage}{0.95\textwidth}
\setlength{\parskip}{0.25em}
\noindent{\footnotesize\itshape \textcolor{profublue}{\textbf{Respuesta esperada:}} Regular el funcionamiento de los sistemas de ahorro para el retiro y sus participantes previstos en esta Ley y en las leyes de seguridad social aplicables.}

\noindent{\footnotesize\itshape \textcolor{profublue}{\textbf{Método de evaluación:}} CompareMeaning}

\noindent{\footnotesize\itshape \textcolor{profublue}{\textbf{Calificación aprobatoria:}} 50}

\noindent{\footnotesize\itshape \textcolor{profublue}{\textbf{Calificación de la respuesta:}} Pass}

\noindent{\footnotesize\itshape \textcolor{profublue}{\textbf{Explicación:}} 75/100: casi el mismo significado; The Agent answer gives all the important information the Expected response has about the law's purpose. It also adds useful details about related laws and where the rule is found in the law. This extra information helps understand the answer better and does not leave out anything important or add wrong facts.}

\end{minipage}

\vspace{2.5em}
\noindent
\begin{minipage}{\textwidth}
\textcolor{profublue}{\textbf{Usuario (80):}}\\[4pt]
\fcolorbox{profublue!30}{profublue!3}{%
  \begin{minipage}{0.95\textwidth}
  \smallskip
  ?Qui?nes son Partes Independientes conforme a la Ley del SAR?
  \smallskip
  \end{minipage}}
\end{minipage}

\vspace{0.4em}
\noindent
\begin{minipage}{\textwidth}
\textcolor{profugold!80!black}{\textbf{Agente (80):}}\\[4pt]
\fcolorbox{profugold!40}{profugold!5}{%
  \begin{minipage}{0.95\textwidth}
  \smallskip
  **Respuesta breve y clara:** Las Partes Independientes, conforme a la Ley de los Sistemas de Ahorro para el Retiro (Ley del SAR), son aquellas personas morales que no tienen nexo patrimonial con una administradora. Es decir, no existe relación de propiedad o control entre la administradora y dichas personas morales.\par **Fundamento legal y localización:** La definición de Partes Independientes se encuentra en la fracción VIII bis del artículo correspondiente, donde se establece que:   - "Partes Independientes, a las personas morales que no tengan nexo patrimonial con una administradora."\par Esta información está ubicada en la Ley del SAR [Páginas 2--3].
  \smallskip
  \end{minipage}}
\end{minipage}

\vspace{0.35em}
\noindent
\begin{minipage}{0.95\textwidth}
\setlength{\parskip}{0.25em}
\noindent{\footnotesize\itshape \textcolor{profublue}{\textbf{Respuesta esperada:}} Personas morales que no tengan nexo patrimonial con una administradora.}

\noindent{\footnotesize\itshape \textcolor{profublue}{\textbf{Método de evaluación:}} CompareMeaning}

\noindent{\footnotesize\itshape \textcolor{profublue}{\textbf{Calificación aprobatoria:}} 50}

\noindent{\footnotesize\itshape \textcolor{profublue}{\textbf{Calificación de la respuesta:}} Pass}

\noindent{\footnotesize\itshape \textcolor{profublue}{\textbf{Explicación:}} 75/100: casi el mismo significado; The agent answer gives the main idea expected: that independent parties are companies with no ownership ties to an administrator. It also adds helpful extra details and legal references. There are no mistakes or missing parts, just more useful information than needed.}

\end{minipage}

\vspace{2.5em}
\noindent
\begin{minipage}{\textwidth}
\textcolor{profublue}{\textbf{Usuario (81):}}\\[4pt]
\fcolorbox{profublue!30}{profublue!3}{%
  \begin{minipage}{0.95\textwidth}
  \smallskip
  Menciona una actividad que forma parte del objeto de las administradoras.
  \smallskip
  \end{minipage}}
\end{minipage}

\vspace{0.4em}
\noindent
\begin{minipage}{\textwidth}
\textcolor{profugold!80!black}{\textbf{Agente (81):}}\\[4pt]
\fcolorbox{profugold!40}{profugold!5}{%
  \begin{minipage}{0.95\textwidth}
  \smallskip
  Respuesta breve: Una de las actividades que forma parte del objeto de las administradoras es administrar y operar cuentas individuales de los trabajadores.\par Fundamento: Esto se encuentra en el Artículo 18, fracción I, donde se establece que las administradoras son entidades financieras dedicadas de manera habitual y profesional a administrar las cuentas individuales y canalizar los recursos de las subcuentas que las integran, así como a administrar sociedades de inversión.\par Referencia: Ley del SAR [Página 12]
  \smallskip
  \end{minipage}}
\end{minipage}

\vspace{0.35em}
\noindent
\begin{minipage}{0.95\textwidth}
\setlength{\parskip}{0.25em}
\noindent{\footnotesize\itshape \textcolor{profublue}{\textbf{Respuesta esperada:}} Abrir, administrar y operar cuentas individuales de los trabajadores.}

\noindent{\footnotesize\itshape \textcolor{profublue}{\textbf{Método de evaluación:}} CompareMeaning}

\noindent{\footnotesize\itshape \textcolor{profublue}{\textbf{Calificación aprobatoria:}} 50}

\noindent{\footnotesize\itshape \textcolor{profublue}{\textbf{Calificación de la respuesta:}} Pass}

\noindent{\footnotesize\itshape \textcolor{profublue}{\textbf{Explicación:}} 50/100: algún significado compartido; The Agent answer mostly matches the Expected response but leaves out that administradoras also "open" accounts, which is important. It adds helpful legal details but misses this key part, so it's mostly correct but not complete.}

\end{minipage}

\vspace{2.5em}
\noindent
\begin{minipage}{\textwidth}
\textcolor{profublue}{\textbf{Usuario (82):}}\\[4pt]
\fcolorbox{profublue!30}{profublue!3}{%
  \begin{minipage}{0.95\textwidth}
  \smallskip
  ?Qu? informes debe presentar el Presidente a la Junta de Gobierno?
  \smallskip
  \end{minipage}}
\end{minipage}

\vspace{0.4em}
\noindent
\begin{minipage}{\textwidth}
\textcolor{profugold!80!black}{\textbf{Agente (82):}}\\[4pt]
\fcolorbox{profugold!40}{profugold!5}{%
  \begin{minipage}{0.95\textwidth}
  \smallskip
  **Respuesta breve y clara:**\par El Presidente debe presentar a la Junta de Gobierno dos tipos de informes: - Un informe semestral sobre la situación que guardan los sistemas de ahorro para el retiro. - Un informe anual sobre las labores desarrolladas por la Comisión.\par Además, debe informar a la Junta de Gobierno sobre todos los asuntos relativos al funcionamiento de los sistemas de ahorro para el retiro y proponer medidas pertinentes cuando se presenten hechos o situaciones que afecten el buen funcionamiento de los mismos.\par **Fundamento y localización en la Ley del SAR:**\par Esto se encuentra en el Artículo 12, fracción III, donde se detallan las facultades y obligaciones del Presidente respecto a la presentación de informes ante la Junta de Gobierno.\par Referencia: Ley del SAR [Páginas 8]
  \smallskip
  \end{minipage}}
\end{minipage}

\vspace{0.35em}
\noindent
\begin{minipage}{0.95\textwidth}
\setlength{\parskip}{0.25em}
\noindent{\footnotesize\itshape \textcolor{profublue}{\textbf{Respuesta esperada:}} Un informe semestral sobre la situaci?n de los sistemas y un informe anual sobre las labores desarrolladas por la Comisi?n.}

\noindent{\footnotesize\itshape \textcolor{profublue}{\textbf{Método de evaluación:}} CompareMeaning}

\noindent{\footnotesize\itshape \textcolor{profublue}{\textbf{Calificación aprobatoria:}} 50}

\noindent{\footnotesize\itshape \textcolor{profublue}{\textbf{Calificación de la respuesta:}} Pass}

\noindent{\footnotesize\itshape \textcolor{profublue}{\textbf{Explicación:}} 75/100: casi el mismo significado; The Agent answer gives all the important information the Expected response has about the reports the President must give. It also adds helpful extra details about other duties the President has, which are true and related. This means the Agent answer is complete and correct but includes more information than asked.}

\end{minipage}

\vspace{2.5em}
\noindent
\begin{minipage}{\textwidth}
\textcolor{profublue}{\textbf{Usuario (83):}}\\[4pt]
\fcolorbox{profublue!30}{profublue!3}{%
  \begin{minipage}{0.95\textwidth}
  \smallskip
  ?Qui?n cubre los gastos de establecimiento y organizaci?n de las sociedades de inversi?n administradas?
  \smallskip
  \end{minipage}}
\end{minipage}

\vspace{0.4em}
\noindent
\begin{minipage}{\textwidth}
\textcolor{profugold!80!black}{\textbf{Agente (83):}}\\[4pt]
\fcolorbox{profugold!40}{profugold!5}{%
  \begin{minipage}{0.95\textwidth}
  \smallskip
  **Respuesta breve y clara**\par Los gastos de establecimiento y organización de las sociedades de inversión administradas deben ser cubiertos por las administradoras, utilizando sus propios ingresos. Es decir, no se trasladan estos gastos a los trabajadores ni a las sociedades de inversión mismas.\par **Fundamento legal y localización**\par Esto se establece en el Artículo 33 de la Ley de los Sistemas de Ahorro para el Retiro, que señala expresamente que las administradoras cubrirán con cargo a sus ingresos todos los gastos necesarios para la operación de las sociedades de inversión que administren.\par Referencia: Ley del SAR [Páginas 19--20].
  \smallskip
  \end{minipage}}
\end{minipage}

\vspace{0.35em}
\noindent
\begin{minipage}{0.95\textwidth}
\setlength{\parskip}{0.25em}
\noindent{\footnotesize\itshape \textcolor{profublue}{\textbf{Respuesta esperada:}} Las administradoras, con cargo a sus ingresos.}

\noindent{\footnotesize\itshape \textcolor{profublue}{\textbf{Método de evaluación:}} CompareMeaning}

\noindent{\footnotesize\itshape \textcolor{profublue}{\textbf{Calificación aprobatoria:}} 50}

\noindent{\footnotesize\itshape \textcolor{profublue}{\textbf{Calificación de la respuesta:}} Pass}

\noindent{\footnotesize\itshape \textcolor{profublue}{\textbf{Explicación:}} 75/100: casi el mismo significado; The Agent answer gives the main point that administrators pay the costs using their own money, just like the Expected response says. It also adds helpful details about the law and who does not pay these costs. The extra information is correct and useful, so the Agent answer matches well and has no mistakes.}

\end{minipage}

\vspace{2.5em}
\noindent
\begin{minipage}{\textwidth}
\textcolor{profublue}{\textbf{Usuario (84):}}\\[4pt]
\fcolorbox{profublue!30}{profublue!3}{%
  \begin{minipage}{0.95\textwidth}
  \smallskip
  ?Qu? obligaci?n tienen las administradoras si cometen actos dolosos contrarios a la Ley que afecten patrimonialmente a trabajadores?
  \smallskip
  \end{minipage}}
\end{minipage}

\vspace{0.4em}
\noindent
\begin{minipage}{\textwidth}
\textcolor{profugold!80!black}{\textbf{Agente (84):}}\\[4pt]
\fcolorbox{profugold!40}{profugold!5}{%
  \begin{minipage}{0.95\textwidth}
  \smallskip
  Obligación de las administradoras ante actos dolosos que afecten patrimonialmente a los trabajadores\par Las administradoras que cometan actos dolosos contrarios a la Ley y que, como consecuencia directa, produzcan una afectación patrimonial a los trabajadores, están obligadas a reparar el daño causado. Es decir, deben resarcir el perjuicio económico que hayan provocado a los trabajadores por dichos actos.\par Fundamento legal: Esta obligación se encuentra establecida en el Artículo 36 de la Ley de los Sistemas de Ahorro para el Retiro, que señala expresamente la responsabilidad directa de las administradoras en estos casos.\par Referencia: Ley del SAR [Páginas 19--20]
  \smallskip
  \end{minipage}}
\end{minipage}

\vspace{0.35em}
\noindent
\begin{minipage}{0.95\textwidth}
\setlength{\parskip}{0.25em}
\noindent{\footnotesize\itshape \textcolor{profublue}{\textbf{Respuesta esperada:}} Reparar el da?o causado.}

\noindent{\footnotesize\itshape \textcolor{profublue}{\textbf{Método de evaluación:}} CompareMeaning}

\noindent{\footnotesize\itshape \textcolor{profublue}{\textbf{Calificación aprobatoria:}} 50}

\noindent{\footnotesize\itshape \textcolor{profublue}{\textbf{Calificación de la respuesta:}} Pass}

\noindent{\footnotesize\itshape \textcolor{profublue}{\textbf{Explicación:}} 75/100: casi el mismo significado; The expected response says the administrators must fix the damage they cause. The agent answer says the same thing and also explains which law says this. It gives more helpful details without leaving anything out or being wrong. So, the agent answer matches well and adds useful information.}

\end{minipage}

\vspace{2.5em}
\noindent
\begin{minipage}{\textwidth}
\textcolor{profublue}{\textbf{Usuario (85):}}\\[4pt]
\fcolorbox{profublue!30}{profublue!3}{%
  \begin{minipage}{0.95\textwidth}
  \smallskip
  Menciona un requisito para ser Presidente de la Comisi?n.
  \smallskip
  \end{minipage}}
\end{minipage}

\vspace{0.4em}
\noindent
\begin{minipage}{\textwidth}
\textcolor{profugold!80!black}{\textbf{Agente (85):}}\\[4pt]
\fcolorbox{profugold!40}{profugold!5}{%
  \begin{minipage}{0.95\textwidth}
  \smallskip
  **Respuesta breve y clara:** Uno de los requisitos para ser Presidente de la Comisión es ser ciudadano mexicano por nacimiento que no adquiera otra nacionalidad y estar en pleno goce de sus derechos civiles y políticos.\par **Fundamento y localización en la Ley del SAR:** Este requisito se encuentra en el Artículo 10, fracción I, de la Ley de los Sistemas de Ahorro para el Retiro, ubicado en las páginas 7 y 8.\par **Referencia:**   Ley del SAR [Páginas 7--8]
  \smallskip
  \end{minipage}}
\end{minipage}

\vspace{0.35em}
\noindent
\begin{minipage}{0.95\textwidth}
\setlength{\parskip}{0.25em}
\noindent{\footnotesize\itshape \textcolor{profublue}{\textbf{Respuesta esperada:}} Ser ciudadano mexicano por nacimiento y estar en pleno goce de derechos; adem?s de contar con experiencia reconocida en materia econ?mica, financiera, jur?dica o de seguridad social.}

\noindent{\footnotesize\itshape \textcolor{profublue}{\textbf{Método de evaluación:}} CompareMeaning}

\noindent{\footnotesize\itshape \textcolor{profublue}{\textbf{Calificación aprobatoria:}} 50}

\noindent{\footnotesize\itshape \textcolor{profublue}{\textbf{Calificación de la respuesta:}} Pass}

\noindent{\footnotesize\itshape \textcolor{profublue}{\textbf{Explicación:}} 50/100: algún significado compartido; The Agent answer gives some correct requirements to be President of the Commission, like being a Mexican citizen and having full rights. But it leaves out an important part about needing experience in certain areas. It also adds where the law says this, which is helpful but not asked for. So, the Agent answer is partly right but misses some key info.}

\end{minipage}

\vspace{2.5em}
\noindent
\begin{minipage}{\textwidth}
\textcolor{profublue}{\textbf{Usuario (86):}}\\[4pt]
\fcolorbox{profublue!30}{profublue!3}{%
  \begin{minipage}{0.95\textwidth}
  \smallskip
  Menciona una facultad clave de la Comisi?n en relaci?n con la operaci?n de los sistemas de ahorro para el retiro.
  \smallskip
  \end{minipage}}
\end{minipage}

\vspace{0.4em}
\noindent
\begin{minipage}{\textwidth}
\textcolor{profugold!80!black}{\textbf{Agente (86):}}\\[4pt]
\fcolorbox{profugold!40}{profugold!5}{%
  \begin{minipage}{0.95\textwidth}
  \smallskip
  **Facultad clave de la Comisión en la operación de los sistemas de ahorro para el retiro**\par La Comisión Nacional del Sistema de Ahorro para el Retiro tiene la facultad de regular, mediante la expedición de disposiciones de carácter general, todo lo relativo a la operación de los sistemas de ahorro para el retiro. Esto incluye la recepción, depósito, transmisión y administración de las cuotas y aportaciones, así como el manejo e intercambio de información entre las dependencias, entidades y participantes, determinando los procedimientos para su buen funcionamiento.\par **Fundamento legal y localización** - Artículo 5o., fracción I - Ley del SAR [Páginas 2--3]
  \smallskip
  \end{minipage}}
\end{minipage}

\vspace{0.35em}
\noindent
\begin{minipage}{0.95\textwidth}
\setlength{\parskip}{0.25em}
\noindent{\footnotesize\itshape \textcolor{profublue}{\textbf{Respuesta esperada:}} Regular mediante disposiciones de car?cter general la operaci?n de los sistemas, incluyendo recepci?n, dep?sito, transmisi?n y administraci?n de cuotas y aportaciones.}

\noindent{\footnotesize\itshape \textcolor{profublue}{\textbf{Método de evaluación:}} CompareMeaning}

\noindent{\footnotesize\itshape \textcolor{profublue}{\textbf{Calificación aprobatoria:}} 50}

\noindent{\footnotesize\itshape \textcolor{profublue}{\textbf{Calificación de la respuesta:}} Pass}

\noindent{\footnotesize\itshape \textcolor{profublue}{\textbf{Explicación:}} 75/100: casi el mismo significado; The agent answer gives all the important points the expected response has, like the commission's power to regulate the retirement savings systems and manage contributions. It also adds extra helpful details about how the system works and legal references. These extra details don't change or leave out anything important, so the answer is mostly complete but has some extra information.}

\end{minipage}

\vspace{2.5em}
\noindent
\begin{minipage}{\textwidth}
\textcolor{profublue}{\textbf{Usuario (87):}}\\[4pt]
\fcolorbox{profublue!30}{profublue!3}{%
  \begin{minipage}{0.95\textwidth}
  \smallskip
  ?Qu? forma societaria deben tener las administradoras?
  \smallskip
  \end{minipage}}
\end{minipage}

\vspace{0.4em}
\noindent
\begin{minipage}{\textwidth}
\textcolor{profugold!80!black}{\textbf{Agente (87):}}\\[4pt]
\fcolorbox{profugold!40}{profugold!5}{%
  \begin{minipage}{0.95\textwidth}
  \smallskip
  **Respuesta breve y clara:** Las administradoras deben constituirse como sociedades anónimas de capital variable. Además, en su denominación o a continuación de ésta, deben incluir la expresión ``Administradora de Fondos para el Retiro'' o su abreviatura ``AFORE''. No pueden usar expresiones en idioma extranjero, nombres de asociaciones religiosas o políticas, ni símbolos religiosos o patrios que sean objeto de devoción o culto público.\par **Fundamento y localización en la Ley del SAR:** Esto se establece en el Artículo 20, fracción I, de la Ley de los Sistemas de Ahorro para el Retiro, que se encuentra en las páginas 14 y 15 del texto.\par **Referencia:**   Ley del SAR [Páginas 14--15]
  \smallskip
  \end{minipage}}
\end{minipage}

\vspace{0.35em}
\noindent
\begin{minipage}{0.95\textwidth}
\setlength{\parskip}{0.25em}
\noindent{\footnotesize\itshape \textcolor{profublue}{\textbf{Respuesta esperada:}} Ser sociedades an?nimas de capital variable y usar en su denominaci?n 'Administradora de Fondos para el Retiro' o la abreviatura 'AFORE'.}

\noindent{\footnotesize\itshape \textcolor{profublue}{\textbf{Método de evaluación:}} CompareMeaning}

\noindent{\footnotesize\itshape \textcolor{profublue}{\textbf{Calificación aprobatoria:}} 50}

\noindent{\footnotesize\itshape \textcolor{profublue}{\textbf{Calificación de la respuesta:}} Pass}

\noindent{\footnotesize\itshape \textcolor{profublue}{\textbf{Explicación:}} 75/100: casi el mismo significado; The Agent answer gives all the important points the Expected response has and adds more useful details about what names and symbols are not allowed. It also says where the law says this. The extra information is correct and helpful but not required. So, the Agent answer matches well but includes more than needed.}

\end{minipage}

\vspace{2.5em}
\noindent
\begin{minipage}{\textwidth}
\textcolor{profublue}{\textbf{Usuario (88):}}\\[4pt]
\fcolorbox{profublue!30}{profublue!3}{%
  \begin{minipage}{0.95\textwidth}
  \smallskip
  ?En qu? debe invertirse el importe restante del capital m?nimo pagado exigido de las administradoras seg?n el art?culo 27?
  \smallskip
  \end{minipage}}
\end{minipage}

\vspace{0.4em}
\noindent
\begin{minipage}{\textwidth}
\textcolor{profugold!80!black}{\textbf{Agente (88):}}\\[4pt]
\fcolorbox{profugold!40}{profugold!5}{%
  \begin{minipage}{0.95\textwidth}
  \smallskip
  **Respuesta breve y clara:**\par El importe restante del capital mínimo pagado exigido de las administradoras debe invertirse en acciones de las sociedades de inversión que administren.\par **Fundamento legal y localización:**\par Esto se establece en el Artículo 27, fracción II, de la Ley de los Sistemas de Ahorro para el Retiro, que señala:\par - ``El importe restante del capital mínimo pagado exigido deberá invertirse en acciones de las sociedades de inversión que administren.''\par La referencia se encuentra en la Ley del SAR [Página 17].
  \smallskip
  \end{minipage}}
\end{minipage}

\vspace{0.35em}
\noindent
\begin{minipage}{0.95\textwidth}
\setlength{\parskip}{0.25em}
\noindent{\footnotesize\itshape \textcolor{profublue}{\textbf{Respuesta esperada:}} En acciones de las sociedades de inversi?n que administren.}

\noindent{\footnotesize\itshape \textcolor{profublue}{\textbf{Método de evaluación:}} CompareMeaning}

\noindent{\footnotesize\itshape \textcolor{profublue}{\textbf{Calificación aprobatoria:}} 50}

\noindent{\footnotesize\itshape \textcolor{profublue}{\textbf{Calificación de la respuesta:}} Pass}

\noindent{\footnotesize\itshape \textcolor{profublue}{\textbf{Explicación:}} 75/100: casi el mismo significado; The Agent answer gives the same key information as the Expected response about where the remaining capital should be invested. It also adds helpful legal details without leaving anything out or adding wrong facts. So, the Agent answer matches well and is complete.}

\end{minipage}

\vspace{2.5em}
\noindent
\begin{minipage}{\textwidth}
\textcolor{profublue}{\textbf{Usuario (89):}}\\[4pt]
\fcolorbox{profublue!30}{profublue!3}{%
  \begin{minipage}{0.95\textwidth}
  \smallskip
  ?C?mo define la Ley a las Sociedades de Inversi?n en este contexto?
  \smallskip
  \end{minipage}}
\end{minipage}

\vspace{0.4em}
\noindent
\begin{minipage}{\textwidth}
\textcolor{profugold!80!black}{\textbf{Agente (89):}}\\[4pt]
\fcolorbox{profugold!40}{profugold!5}{%
  \begin{minipage}{0.95\textwidth}
  \smallskip
  **Definición de Sociedades de Inversión según la Ley del SAR**\par Las sociedades de inversión, en el contexto de la Ley de los Sistemas de Ahorro para el Retiro, son entidades administradas y operadas por las administradoras. Su objetivo principal es invertir los recursos provenientes de las cuentas individuales que reciben conforme a las leyes de seguridad social y la propia Ley del SAR. Además, pueden invertir recursos de las administradoras, aportaciones destinadas a fondos de previsión social, aportaciones voluntarias y complementarias de retiro de trabajadores y patrones, así como otros recursos que puedan ser depositados en las cuentas individuales según la Ley.\par **Fundamento legal y localización** - Artículo 39, Sección II, Capítulo correspondiente. - Ley del SAR [Página 23] [1]\par [1]: cite:1 "Citation-1"
  \smallskip
  \end{minipage}}
\end{minipage}

\vspace{0.35em}
\noindent
\begin{minipage}{0.95\textwidth}
\setlength{\parskip}{0.25em}
\noindent{\footnotesize\itshape \textcolor{profublue}{\textbf{Respuesta esperada:}} Las sociedades de inversi?n especializadas de fondos para el retiro.}

\noindent{\footnotesize\itshape \textcolor{profublue}{\textbf{Método de evaluación:}} CompareMeaning}

\noindent{\footnotesize\itshape \textcolor{profublue}{\textbf{Calificación aprobatoria:}} 50}

\noindent{\footnotesize\itshape \textcolor{profublue}{\textbf{Calificación de la respuesta:}} Pass}

\noindent{\footnotesize\itshape \textcolor{profublue}{\textbf{Explicación:}} 75/100: casi el mismo significado; The Expected response says that the law defines investment companies as those specialized in retirement funds. The Agent answer explains this and adds more details about how these companies work and what they do. The extra information is correct but not needed to answer the question. So, the Agent answer matches well but includes more than asked.}

\end{minipage}

\vspace{2.5em}
\noindent
\begin{minipage}{\textwidth}
\textcolor{profublue}{\textbf{Usuario (90):}}\\[4pt]
\fcolorbox{profublue!30}{profublue!3}{%
  \begin{minipage}{0.95\textwidth}
  \smallskip
  ?Qu? tipo de comisiones pueden cobrar las administradoras a los trabajadores con cuenta individual?
  \smallskip
  \end{minipage}}
\end{minipage}

\vspace{0.4em}
\noindent
\begin{minipage}{\textwidth}
\textcolor{profugold!80!black}{\textbf{Agente (90):}}\\[4pt]
\fcolorbox{profugold!40}{profugold!5}{%
  \begin{minipage}{0.95\textwidth}
  \smallskip
  **Respuesta breve y clara:**\par Las administradoras solo pueden cobrar a los trabajadores con cuenta individual comisiones que estén establecidas conforme a las reglas generales que expida la Comisión. Específicamente, las comisiones por administración de cuentas individuales únicamente pueden cobrarse como un porcentaje sobre el valor de los activos administrados. Además, las administradoras pueden cobrar cuotas fijas únicamente por servicios señalados en el reglamento de la Ley, pero nunca por la administración de las cuentas.\par **Fundamento y localización en la Ley del SAR:**\par - Artículo 37: Establece que las comisiones por administración solo pueden ser un porcentaje sobre el valor de los activos administrados y que las cuotas fijas solo pueden cobrarse por servicios específicos del reglamento, no por la administración de cuentas. - El contenido se encuentra en las páginas 19 y 20.\par **Referencia:**   Ley del SAR [Páginas 19--20]
  \smallskip
  \end{minipage}}
\end{minipage}

\vspace{0.35em}
\noindent
\begin{minipage}{0.95\textwidth}
\setlength{\parskip}{0.25em}
\noindent{\footnotesize\itshape \textcolor{profublue}{\textbf{Respuesta esperada:}} Solo las comisiones con cargo a esas cuentas que se establezcan conforme a reglas de car?cter general de la Comisi?n.}

\noindent{\footnotesize\itshape \textcolor{profublue}{\textbf{Método de evaluación:}} CompareMeaning}

\noindent{\footnotesize\itshape \textcolor{profublue}{\textbf{Calificación aprobatoria:}} 50}

\noindent{\footnotesize\itshape \textcolor{profublue}{\textbf{Calificación de la respuesta:}} Pass}

\noindent{\footnotesize\itshape \textcolor{profublue}{\textbf{Explicación:}} 75/100: casi el mismo significado; The Expected response says commissions can only be charged if set by general rules from the Commission. The Agent answer says this too and adds more details about how commissions work and when fixed fees can be charged. These extra details are correct and helpful, not wrong or off-topic. So, the Agent answer matches well but has more information than needed.}

\end{minipage}

\vspace{2.5em}
\noindent
\begin{minipage}{\textwidth}
\textcolor{profublue}{\textbf{Usuario (91):}}\\[4pt]
\fcolorbox{profublue!30}{profublue!3}{%
  \begin{minipage}{0.95\textwidth}
  \smallskip
  Menciona una facultad del Comit? Consultivo y de Vigilancia.
  \smallskip
  \end{minipage}}
\end{minipage}

\vspace{0.4em}
\noindent
\begin{minipage}{\textwidth}
\textcolor{profugold!80!black}{\textbf{Agente (91):}}\\[4pt]
\fcolorbox{profugold!40}{profugold!5}{%
  \begin{minipage}{0.95\textwidth}
  \smallskip
  Respuesta breve: Una facultad del Comité Consultivo y de Vigilancia es vigilar el desarrollo de los sistemas de ahorro para el retiro para prevenir posibles situaciones que presenten conflicto de interés y prácticas monopólicas.\par Fundamento: Esta facultad se encuentra en el Artículo 16, fracción II, de la Ley de los Sistemas de Ahorro para el Retiro, ubicada en las páginas 10--11.\par Referencia: Ley del SAR [Páginas 10--11].
  \smallskip
  \end{minipage}}
\end{minipage}

\vspace{0.35em}
\noindent
\begin{minipage}{0.95\textwidth}
\setlength{\parskip}{0.25em}
\noindent{\footnotesize\itshape \textcolor{profublue}{\textbf{Respuesta esperada:}} Emitir opini?n respecto de reglas generales del r?gimen de inversi?n y del r?gimen de comisiones y su aplicaci?n.}

\noindent{\footnotesize\itshape \textcolor{profublue}{\textbf{Método de evaluación:}} CompareMeaning}

\noindent{\footnotesize\itshape \textcolor{profublue}{\textbf{Calificación aprobatoria:}} 50}

\noindent{\footnotesize\itshape \textcolor{profublue}{\textbf{Calificación de la respuesta:}} Fail}

\noindent{\footnotesize\itshape \textcolor{profublue}{\textbf{Explicación:}} 0/100: significado diferente; The agent answer talks about a different duty than the one asked for. It gives a correct duty but not the one expected. So, it doesn't answer the question properly and gives information that doesn't match what was asked.}

\end{minipage}

\vspace{2.5em}
\noindent
\begin{minipage}{\textwidth}
\textcolor{profublue}{\textbf{Usuario (92):}}\\[4pt]
\fcolorbox{profublue!30}{profublue!3}{%
  \begin{minipage}{0.95\textwidth}
  \smallskip
  ?Qu? son las administradoras en t?rminos de la Ley del SAR?
  \smallskip
  \end{minipage}}
\end{minipage}

\vspace{0.4em}
\noindent
\begin{minipage}{\textwidth}
\textcolor{profugold!80!black}{\textbf{Agente (92):}}\\[4pt]
\fcolorbox{profugold!40}{profugold!5}{%
  \begin{minipage}{0.95\textwidth}
  \smallskip
  **Respuesta breve y clara**\par Las administradoras, según la Ley de los Sistemas de Ahorro para el Retiro (Ley del SAR), son entidades financieras que se dedican de manera habitual y profesional a administrar las cuentas individuales de los trabajadores y canalizar los recursos de las subcuentas que las integran. Además, administran sociedades de inversión relacionadas con estos fondos. Su función principal es gestionar los recursos para obtener una adecuada rentabilidad y seguridad en las inversiones, siempre atendiendo exclusivamente al interés de los trabajadores.\par **Fundamento legal y localización**\par - Capítulo III ``De los Participantes en los Sistemas de Ahorro para el Retiro'' - Sección I ``De las Administradoras de Fondos para el Retiro'' - Artículo 18: Define a las administradoras y sus funciones principales.\par Referencia: Ley del SAR [Página 12]
  \smallskip
  \end{minipage}}
\end{minipage}

\vspace{0.35em}
\noindent
\begin{minipage}{0.95\textwidth}
\setlength{\parskip}{0.25em}
\noindent{\footnotesize\itshape \textcolor{profublue}{\textbf{Respuesta esperada:}} Entidades financieras que administran cuentas individuales y canalizan recursos de subcuentas, adem?s de administrar sociedades de inversi?n.}

\noindent{\footnotesize\itshape \textcolor{profublue}{\textbf{Método de evaluación:}} CompareMeaning}

\noindent{\footnotesize\itshape \textcolor{profublue}{\textbf{Calificación aprobatoria:}} 50}

\noindent{\footnotesize\itshape \textcolor{profublue}{\textbf{Calificación de la respuesta:}} Pass}

\noindent{\footnotesize\itshape \textcolor{profublue}{\textbf{Explicación:}} 75/100: casi el mismo significado; The agent answer gives all the important points from the expected response and adds more helpful details about what administradoras do and where the law says this. There are no mistakes, but it includes extra information that wasn't asked for. So, it matches well but isn't exactly the same.}

\end{minipage}

\vspace{2.5em}
\noindent
\begin{minipage}{\textwidth}
\textcolor{profublue}{\textbf{Usuario (93):}}\\[4pt]
\fcolorbox{profublue!30}{profublue!3}{%
  \begin{minipage}{0.95\textwidth}
  \smallskip
  ?Puede el Presidente delegar facultades? ?En qui?nes?
  \smallskip
  \end{minipage}}
\end{minipage}

\vspace{0.4em}
\noindent
\begin{minipage}{\textwidth}
\textcolor{profugold!80!black}{\textbf{Agente (93):}}\\[4pt]
\fcolorbox{profugold!40}{profugold!5}{%
  \begin{minipage}{0.95\textwidth}
  \smallskip
  **Respuesta breve y clara**\par Sí, el Presidente de la Comisión puede delegar facultades. Según la Ley del SAR, esta delegación puede hacerse directamente o a través de los servidores públicos de la Comisión, conforme al reglamento interior o mediante acuerdos delegatorios que deben publicarse en el Diario Oficial de la Federación. Específicamente, en procedimientos judiciales, administrativos o laborales, el Presidente puede ejercer acciones directamente o por medio de los Vicepresidentes o Directores Generales que designe en los acuerdos delegatorios.\par **Fundamento y localización en la Ley**\par - Artículo 11: Establece que el Presidente puede ejercer facultades directamente o a través de servidores públicos de la Comisión, según el reglamento interior o acuerdos delegatorios publicados. - Artículo 12, fracción I: Señala que en procedimientos judiciales, administrativos o laborales, el Presidente puede delegar facultades en Vicepresidentes o Directores Generales mediante acuerdos delegatorios.\par **Referencia** Ley del SAR [Páginas 8--9]
  \smallskip
  \end{minipage}}
\end{minipage}

\vspace{0.35em}
\noindent
\begin{minipage}{0.95\textwidth}
\setlength{\parskip}{0.25em}
\noindent{\footnotesize\itshape \textcolor{profublue}{\textbf{Respuesta esperada:}} S?, puede delegar ciertas facultades en Vicepresidentes y Directores Generales mediante acuerdo publicado en el DOF, conforme a la Ley.}

\noindent{\footnotesize\itshape \textcolor{profublue}{\textbf{Método de evaluación:}} CompareMeaning}

\noindent{\footnotesize\itshape \textcolor{profublue}{\textbf{Calificación aprobatoria:}} 50}

\noindent{\footnotesize\itshape \textcolor{profublue}{\textbf{Calificación de la respuesta:}} Pass}

\noindent{\footnotesize\itshape \textcolor{profublue}{\textbf{Explicación:}} 75/100: casi el mismo significado; The Agent answer correctly says the President can delegate powers to Vice Presidents and Directors General, just like the Expected response. It also gives more helpful details about how and when this can happen, without leaving out anything important or adding wrong information.}

\end{minipage}

\vspace{2.5em}
\noindent
\begin{minipage}{\textwidth}
\textcolor{profublue}{\textbf{Usuario (94):}}\\[4pt]
\fcolorbox{profublue!30}{profublue!3}{%
  \begin{minipage}{0.95\textwidth}
  \smallskip
  ?Qu? obligaci?n informativa tiene la Comisi?n hacia el Congreso?
  \smallskip
  \end{minipage}}
\end{minipage}

\vspace{0.4em}
\noindent
\begin{minipage}{\textwidth}
\textcolor{profugold!80!black}{\textbf{Agente (94):}}\\[4pt]
\fcolorbox{profugold!40}{profugold!5}{%
  \begin{minipage}{0.95\textwidth}
  \smallskip
  Obligación informativa de la Comisión hacia el Congreso\par Respuesta breve: La Ley del SAR establece que la Comisión tiene la obligación de informar y reportar sobre actos, operaciones y servicios relacionados con los sistemas de ahorro para el retiro, siguiendo las disposiciones de carácter general. Estos reportes deben elaborarse y presentarse conforme a las modalidades y lineamientos que determine la Secretaría de Hacienda y Crédito Público, incluyendo la periodicidad y los sistemas de transmisión de la información. Además, la Comisión puede ser requerida para recabar y proporcionar información relevante a las autoridades competentes.\par Fundamento y localización: Esta obligación se encuentra regulada en el texto de la Ley del SAR, en el artículo que detalla los reportes y la transmisión de información, así como la facultad de la Secretaría de Hacienda para requerir información por conducto de la Comisión. El contenido relevante se localiza en la Página 65.\par Referencia: Ley del SAR [Página 65] [1]\par [1]: cite:1 "Citation-1"
  \smallskip
  \end{minipage}}
\end{minipage}

\vspace{0.35em}
\noindent
\begin{minipage}{0.95\textwidth}
\setlength{\parskip}{0.25em}
\noindent{\footnotesize\itshape \textcolor{profublue}{\textbf{Respuesta esperada:}} Rendir un informe trimestral sobre la situaci?n de los Sistemas de Ahorro para el Retiro con apartados espec?ficos sobre inversi?n, valores extranjeros y protecci?n de recursos.}

\noindent{\footnotesize\itshape \textcolor{profublue}{\textbf{Método de evaluación:}} CompareMeaning}

\noindent{\footnotesize\itshape \textcolor{profublue}{\textbf{Calificación aprobatoria:}} 50}

\noindent{\footnotesize\itshape \textcolor{profublue}{\textbf{Calificación de la respuesta:}} Pass}

\noindent{\footnotesize\itshape \textcolor{profublue}{\textbf{Explicación:}} 50/100: algún significado compartido; The Expected response says the Commission must give a report every three months with specific parts about investments and protecting resources. The Agent answer says the Commission must report but doesn't say it's every three months or mention those specific parts. It also adds some extra details about the law and other reporting duties. So, the Agent answer partly answers the question but misses important details about the report's timing and content.}

\end{minipage}

\vspace{2.5em}
\noindent
\begin{minipage}{\textwidth}
\textcolor{profublue}{\textbf{Usuario (95):}}\\[4pt]
\fcolorbox{profublue!30}{profublue!3}{%
  \begin{minipage}{0.95\textwidth}
  \smallskip
  ?Qu? facultad tiene la Junta de Gobierno respecto a comisiones consideradas excesivas?
  \smallskip
  \end{minipage}}
\end{minipage}

\vspace{0.4em}
\noindent
\begin{minipage}{\textwidth}
\textcolor{profugold!80!black}{\textbf{Agente (95):}}\\[4pt]
\fcolorbox{profugold!40}{profugold!5}{%
  \begin{minipage}{0.95\textwidth}
  \smallskip
  Respuesta breve y clara: La Junta de Gobierno tiene la facultad de negar la autorización de comisiones si considera que son excesivas para los intereses de los trabajadores. Para ello, puede exigir información adicional, pedir aclaraciones o adecuaciones, y si lo estima necesario, denegar la autorización. Además, dicta políticas y criterios sobre comisiones, puede emitir recomendaciones a las administradoras y debe hacer públicas las razones de la negativa, salvo que la información sea reservada o confidencial.\par Fundamento y localización: Esta facultad se encuentra en el artículo que regula la autorización de comisiones, donde se establece que la Junta de Gobierno puede denegar la autorización si las comisiones son excesivas, considerando factores como el monto de los activos, la estructura de costos, el nivel de otras comisiones en el mercado y otros elementos pertinentes. También se menciona la obligación de fundamentar y motivar la resolución, así como la publicación de las razones de la negativa.\par Referencia: Ley del SAR [Páginas 20--21].
  \smallskip
  \end{minipage}}
\end{minipage}

\vspace{0.35em}
\noindent
\begin{minipage}{0.95\textwidth}
\setlength{\parskip}{0.25em}
\noindent{\footnotesize\itshape \textcolor{profublue}{\textbf{Respuesta esperada:}} Puede exigir informaci?n adicional, pedir adecuaciones o denegar la autorizaci?n si son excesivas para los intereses de los trabajadores.}

\noindent{\footnotesize\itshape \textcolor{profublue}{\textbf{Método de evaluación:}} CompareMeaning}

\noindent{\footnotesize\itshape \textcolor{profublue}{\textbf{Calificación aprobatoria:}} 50}

\noindent{\footnotesize\itshape \textcolor{profublue}{\textbf{Calificación de la respuesta:}} Pass}

\noindent{\footnotesize\itshape \textcolor{profublue}{\textbf{Explicación:}} 75/100: casi el mismo significado; The Agent answer includes all the important points from the Expected response and adds more helpful details about the Junta de Gobierno's powers. It doesn't leave out anything important and doesn't have wrong information. The extra details are useful but go beyond what was asked.}

\end{minipage}

\vspace{2.5em}
\noindent
\begin{minipage}{\textwidth}
\textcolor{profublue}{\textbf{Usuario (96):}}\\[4pt]
\fcolorbox{profublue!30}{profublue!3}{%
  \begin{minipage}{0.95\textwidth}
  \smallskip
  ?Qu? regla aplica a las inversiones con cargo al capital m?nimo pagado exigido de las administradoras?
  \smallskip
  \end{minipage}}
\end{minipage}

\vspace{0.4em}
\noindent
\begin{minipage}{\textwidth}
\textcolor{profugold!80!black}{\textbf{Agente (96):}}\\[4pt]
\fcolorbox{profugold!40}{profugold!5}{%
  \begin{minipage}{0.95\textwidth}
  \smallskip
  **Regla para inversiones con cargo al capital mínimo pagado exigido de las administradoras**\par Las inversiones que realicen las administradoras con cargo al capital mínimo pagado exigido se sujetan a las siguientes reglas:\par - No podrá exceder del 40\% del capital mínimo pagado exigido el importe de las inversiones en mobiliario y equipo, inmuebles, derechos reales que no sean de garantía o gastos de instalación, más el importe de las inversiones en el capital de empresas que les presten servicios complementarios o auxiliares. - El importe restante del capital mínimo pagado exigido debe invertirse en acciones de las sociedades de inversión que administren. - La Comisión puede autorizar que el porcentaje de la primera regla sea mayor, pero nunca podrá exceder del 60\%.\par **Fundamento legal:**   Artículo 27, Ley del SAR [Página 17].
  \smallskip
  \end{minipage}}
\end{minipage}

\vspace{0.35em}
\noindent
\begin{minipage}{0.95\textwidth}
\setlength{\parskip}{0.25em}
\noindent{\footnotesize\itshape \textcolor{profublue}{\textbf{Respuesta esperada:}} El importe en mobiliario, equipo, inmuebles, derechos reales no de garant?a y gastos de instalaci?n, m?s inversiones en empresas de servicios auxiliares, no debe exceder 40\% del capital m?nimo pagado.}

\noindent{\footnotesize\itshape \textcolor{profublue}{\textbf{Método de evaluación:}} CompareMeaning}

\noindent{\footnotesize\itshape \textcolor{profublue}{\textbf{Calificación aprobatoria:}} 50}

\noindent{\footnotesize\itshape \textcolor{profublue}{\textbf{Calificación de la respuesta:}} Pass}

\noindent{\footnotesize\itshape \textcolor{profublue}{\textbf{Explicación:}} 75/100: casi el mismo significado; The Agent answer includes all the important parts of the Expected response and adds more helpful details that are correct. It explains the 40\% limit and also mentions other rules about where the rest of the money should go and possible exceptions. This extra information is useful and does not contradict anything, so the answer is very good but a bit more detailed than needed.}

\end{minipage}

\vspace{2.5em}
\noindent
\begin{minipage}{\textwidth}
\textcolor{profublue}{\textbf{Usuario (97):}}\\[4pt]
\fcolorbox{profublue!30}{profublue!3}{%
  \begin{minipage}{0.95\textwidth}
  \smallskip
  ?Qu? miembros forman parte de la Junta de Gobierno adem?s del Secretario de Hacienda y el Presidente de la Comisi?n?
  \smallskip
  \end{minipage}}
\end{minipage}

\vspace{0.4em}
\noindent
\begin{minipage}{\textwidth}
\textcolor{profugold!80!black}{\textbf{Agente (97):}}\\[4pt]
\fcolorbox{profugold!40}{profugold!5}{%
  \begin{minipage}{0.95\textwidth}
  \smallskip
  **Miembros de la Junta de Gobierno además del Secretario de Hacienda y el Presidente de la Comisión**\par La Junta de Gobierno está integrada, además del Secretario de Hacienda y el Presidente de la Comisión, por los siguientes miembros:\par - Dos vicepresidentes de la Comisión. - Trece vocales, que son:   - El Secretario del Trabajo y Previsión Social.   - El Gobernador del Banco de México.   - El Subsecretario de Hacienda y Crédito Público.   - El Director General del Instituto Mexicano del Seguro Social.   - El Director General del Instituto del Fondo Nacional de la Vivienda para los Trabajadores.   - El Director General del Instituto de Seguridad y Servicios Sociales de los Trabajadores del Estado.   - El Presidente de la Comisión Nacional Bancaria y de Valores.   - El Presidente de la Comisión Nacional de Seguros y Fianzas.   - Cinco vocales designados por el Secretario de Hacienda y Crédito Público: cuatro representantes de organizaciones nacionales de trabajadores y uno de organizaciones de patrones, que formen parte del Comité Consultivo y de Vigilancia y ostenten la mayor representatividad.\par **Fundamento legal:**   Esta integración se encuentra en el Artículo 7o de la Ley del SAR, en la Sección II ``De los Órganos de Gobierno'', abarcando la información desde la Página 5 y continuando en la Página 6.\par **Referencia:**   Ley del SAR [Páginas 5--6]
  \smallskip
  \end{minipage}}
\end{minipage}

\vspace{0.35em}
\noindent
\begin{minipage}{0.95\textwidth}
\setlength{\parskip}{0.25em}
\noindent{\footnotesize\itshape \textcolor{profublue}{\textbf{Respuesta esperada:}} Incluye dos vicepresidentes de la Comisi?n y trece vocales, entre ellos secretarios y titulares de instituciones como STPS, Banxico, IMSS, INFONAVIT, ISSSTE, CNBV y CNSF.}

\noindent{\footnotesize\itshape \textcolor{profublue}{\textbf{Método de evaluación:}} CompareMeaning}

\noindent{\footnotesize\itshape \textcolor{profublue}{\textbf{Calificación aprobatoria:}} 50}

\noindent{\footnotesize\itshape \textcolor{profublue}{\textbf{Calificación de la respuesta:}} Pass}

\noindent{\footnotesize\itshape \textcolor{profublue}{\textbf{Explicación:}} 75/100: casi el mismo significado; The Agent answer gives all the important members listed in the Expected response and adds more details about who they are. It does not leave out anything important or say anything wrong. It also includes where the information comes from, which helps make the answer clear and complete.}

\end{minipage}

\vspace{2.5em}
\noindent
\begin{minipage}{\textwidth}
\textcolor{profublue}{\textbf{Usuario (98):}}\\[4pt]
\fcolorbox{profublue!30}{profublue!3}{%
  \begin{minipage}{0.95\textwidth}
  \smallskip
  ?Qu? elementos m?nimos deben contener los contratos de administraci?n de fondos para el retiro?
  \smallskip
  \end{minipage}}
\end{minipage}

\vspace{0.4em}
\noindent
\begin{minipage}{\textwidth}
\textcolor{profugold!80!black}{\textbf{Agente (98):}}\\[4pt]
\fcolorbox{profugold!40}{profugold!5}{%
  \begin{minipage}{0.95\textwidth}
  \smallskip
  **Elementos mínimos de los contratos de administración de fondos para el retiro**\par De acuerdo con la Ley de los Sistemas de Ahorro para el Retiro, los contratos de administración de fondos para el retiro deben contener, como mínimo, los siguientes elementos:\par - El objeto del contrato. - El tipo de trabajador con el que se celebra, conforme a las definiciones del artículo 3o. de la Ley. - Las obligaciones específicas de la administradora. - La elección de las sociedades de inversión por parte del trabajador. - La estructura y el cobro de comisiones por los servicios prestados por la administradora. - La responsabilidad de la administradora por sus actos y los de las sociedades de inversión que administre. - La vigencia del contrato y sus causas de terminación.\par **Fundamento legal:**   Estos requisitos se encuentran en el artículo correspondiente de la Ley del SAR, en la página 18.\par Ley del SAR [Página 18]
  \smallskip
  \end{minipage}}
\end{minipage}

\vspace{0.35em}
\noindent
\begin{minipage}{0.95\textwidth}
\setlength{\parskip}{0.25em}
\noindent{\footnotesize\itshape \textcolor{profublue}{\textbf{Respuesta esperada:}} Objeto del contrato, tipo de trabajador, obligaciones de la administradora, elecci?n de sociedades de inversi?n, estructura de comisiones, responsabilidad, vigencia y causas de terminaci?n.}

\noindent{\footnotesize\itshape \textcolor{profublue}{\textbf{Método de evaluación:}} CompareMeaning}

\noindent{\footnotesize\itshape \textcolor{profublue}{\textbf{Calificación aprobatoria:}} 50}

\noindent{\footnotesize\itshape \textcolor{profublue}{\textbf{Calificación de la respuesta:}} Pass}

\noindent{\footnotesize\itshape \textcolor{profublue}{\textbf{Explicación:}} 75/100: casi el mismo significado; The Agent answer includes all the important parts the Expected response has and adds useful extra details about the law source. It doesn't leave anything out or say anything wrong. So, it matches very well.}

\end{minipage}

\vspace{2.5em}
\noindent
\begin{minipage}{\textwidth}
\textcolor{profublue}{\textbf{Usuario (99):}}\\[4pt]
\fcolorbox{profublue!30}{profublue!3}{%
  \begin{minipage}{0.95\textwidth}
  \smallskip
  ?Qu? obligaci?n de reserva especial tienen las administradoras conforme al art?culo 28?
  \smallskip
  \end{minipage}}
\end{minipage}

\vspace{0.4em}
\noindent
\begin{minipage}{\textwidth}
\textcolor{profugold!80!black}{\textbf{Agente (99):}}\\[4pt]
\fcolorbox{profugold!40}{profugold!5}{%
  \begin{minipage}{0.95\textwidth}
  \smallskip
  Obligación de reserva especial de las administradoras según el artículo 28\par Respuesta breve y clara: Las administradoras están obligadas a constituir y mantener una reserva especial, la cual debe estar invertida en las acciones de cada una de las sociedades de inversión que administren. El monto y la composición de esta reserva especial serán determinados por la Junta de Gobierno mediante disposiciones de carácter general, considerando el capital suscrito y pagado por los trabajadores y la naturaleza de cada sociedad de inversión.\par Fundamento: - Artículo 28, Ley del SAR [Página 17]
  \smallskip
  \end{minipage}}
\end{minipage}

\vspace{0.35em}
\noindent
\begin{minipage}{0.95\textwidth}
\setlength{\parskip}{0.25em}
\noindent{\footnotesize\itshape \textcolor{profublue}{\textbf{Respuesta esperada:}} Constituir y mantener una reserva especial invertida en acciones de cada sociedad de inversi?n que administren, en monto y composici?n definidos por la Junta de Gobierno.}

\noindent{\footnotesize\itshape \textcolor{profublue}{\textbf{Método de evaluación:}} CompareMeaning}

\noindent{\footnotesize\itshape \textcolor{profublue}{\textbf{Calificación aprobatoria:}} 50}

\noindent{\footnotesize\itshape \textcolor{profublue}{\textbf{Calificación de la respuesta:}} Pass}

\noindent{\footnotesize\itshape \textcolor{profublue}{\textbf{Explicación:}} 75/100: casi el mismo significado; The Agent answer includes all the important points from the Expected response and adds helpful extra details that explain the rules better. It does not leave out anything or give wrong information. So, it matches very well with the Expected response.}

\end{minipage}

\vspace{2.5em}
\noindent
\begin{minipage}{\textwidth}
\textcolor{profublue}{\textbf{Usuario (100):}}\\[4pt]
\fcolorbox{profublue!30}{profublue!3}{%
  \begin{minipage}{0.95\textwidth}
  \smallskip
  ?Puede la Comisi?n autorizar un porcentaje mayor al 40\% para inversiones se?aladas en el art?culo 27?
  \smallskip
  \end{minipage}}
\end{minipage}

\vspace{0.4em}
\noindent
\begin{minipage}{\textwidth}
\textcolor{profugold!80!black}{\textbf{Agente (100):}}\\[4pt]
\fcolorbox{profugold!40}{profugold!5}{%
  \begin{minipage}{0.95\textwidth}
  \smallskip
  Respuesta breve: Sí, la Comisión puede autorizar un porcentaje mayor al 40\% para las inversiones señaladas en la fracción I del artículo 27, pero este porcentaje no puede exceder del 60\%.\par Fundamento: Esto se encuentra expresamente en el propio artículo 27, el cual establece que la Comisión podrá autorizar un porcentaje mayor al 40\% para las inversiones en mobiliario, equipo, inmuebles, derechos reales no de garantía o gastos de instalación, así como en el capital de empresas que presten servicios complementarios o auxiliares, pero con un límite máximo del 60\%.\par Referencia: Ley del SAR [Página 17].
  \smallskip
  \end{minipage}}
\end{minipage}

\vspace{0.35em}
\noindent
\begin{minipage}{0.95\textwidth}
\setlength{\parskip}{0.25em}
\noindent{\footnotesize\itshape \textcolor{profublue}{\textbf{Respuesta esperada:}} S?, puede autorizar un porcentaje mayor sin exceder del 60\%.}

\noindent{\footnotesize\itshape \textcolor{profublue}{\textbf{Método de evaluación:}} CompareMeaning}

\noindent{\footnotesize\itshape \textcolor{profublue}{\textbf{Calificación aprobatoria:}} 50}

\noindent{\footnotesize\itshape \textcolor{profublue}{\textbf{Calificación de la respuesta:}} Pass}

\noindent{\footnotesize\itshape \textcolor{profublue}{\textbf{Explicación:}} 75/100: casi el mismo significado; The Agent answer correctly says the Commission can allow more than 40\% but not over 60\%, just like the Expected response. It also gives more details about the types of investments and the law, which helps but isn't required. So, the Agent answer matches well but has extra helpful info.}

\end{minipage}

\vspace{2em}
{\color{profugold}\rule{\textwidth}{1.5pt}}

\begin{center}
\textit{Fin del reporte de evaluación.}
\end{center}

\end{document}
